\renewcommand{\textbf}[1]{#1}
\chapter{Onde está o laboratório e os seus cientistas?}
\label{capitulo1}

\epigrafe{%
A Pesquisa é a zona para a qual são arrastados humanos e não humanos, onde ao longo das eras foi feito o mais extraordinário dos experimentos coletivos para distinguir, em tempo real, o ``cosmos'' da ``desordem'' sem que ninguém, cientista ou ``estudioso da ciência'', pudesse saber de antemão qual seria a resposta provisória. Talvez, afinal das contas, os estudos de ciência sejam Anticiência.}%
{Bruno Latour, \textit{A ciência em ação: como seguir cientistas e engenheiros sociedade afora}, 2011.}

Em dezembro de 2022, estive no México para participar do congresso da \textit{Society for Social Studies of Science} (4S), onde apresentei pela primeira vez os resultados preliminares desta pesquisa. Além da experiência de conhecer o país e participar de um evento internacional na área de estudos sociais da ciência e tecnologia, viajei com um grupo diverso de pesquisadores e amigos que contribuíram para aprofundar a reflexão sobre o que significa fazer ciências sociais e investigar como a ciência é produzida em outros áreas do conhecimento. Clarissa, uma das companheiras da viagem e do doutorado na Unicamp, investigava um tema muito diferente sobre o trabalho de antropólogas que menstruam em campo. Ainda assim, compartilhávamos uma preocupação: a ideia de ``não ferir os sentimentos estabelecidos'', uma reflexão de Isabelle Stengers,\footnote{Stengers propõe em \textit{Another Science is Possible} uma ``ciência lenta'' (\textit{slow science}) como resistência às pressões de produtividade que degradam a qualidade do pensamento científico \parencite{stengers2018another}. A defesa da lentidão pressupõe condições de autonomia que nem sempre são possíveis. Esta tese documenta como a temporalidade acadêmica pode se tornar condição de vulnerabilidade quando articulada às lógicas temporais corporativas. Enquanto universidades públicas operam em horizontes de décadas, corporações decidem em ciclos de curto prazo segundo lógicas de retorno sobre investimentos. A capa da edição inglesa traz um caracol, imagem que ressoa com o lema zapatista ``lento, pero avanzo'' e que esta tese adota na dedicatória.} ``como estratégia para criar alianças com interlocutores cientistas, engenheiros e áreas afins''. Longe de se configurar como atitude conformista, trata-se justamente do oposto: ``uma postura calcada apenas no confronto e na denúncia das mazelas da tecnociência é infrutífera, pois produz um resultado alienante e perverso: oculta para nós, pessoas que pesquisam desde as ciências humanas, nossa responsabilidade enquanto coabitantes desse mesmo espaço onde estão os cientistas que denunciamos'' \parencite[p. 314]{costa2024manchando}. Isso não significa que as denúncias sejam irrelevantes, mas ``reconhecer nossa responsabilidade como coabitantes do sistema de produção de conhecimento científico é também aprender a construir e sustentar diálogos, por mais difíceis que sejam'' \parencite[p. 314]{costa2024manchando}.\footnote{Costa desenvolve o conceito de \textit{armadilha} como estratégia metodológica e política para que corpos marcados produzam conhecimento crítico dentro de estruturas acadêmicas sem serem devorados por elas. A armadilha opera em duas etapas, primeiro, encanta através da beleza estética; depois, revela que as imagens fractais são feitas de sangue menstrual, vazando questões feministas de forma que não sejam imediatamente rejeitadas. A etnografia também pode operar como armadilha. Seguir os atores sem julgamento prévio, ganhar confiança, descrever o que é dito e feito para depois tornar visíveis tensões e contradições que os próprios atores não nomeiam \parencite{costa2024manchando}.}

Quase três anos depois do congresso no México, em novembro de 2025, a IBM fechou seu laboratório de pesquisa no Brasil, o único da América Latina, operando há 15 anos com mais de 100 pesquisadores. Em dezembro de 2025, encerrou-se o convênio C4AI-IBM. Desde o início desta pesquisa eu me perguntava até quando aquilo duraria. Na época, havia expectativa por parte do C4AI de que se estendesse por dez anos. Acabou em cinco. A confiança com que meus interlocutores falavam sobre a participação da IBM e a solidez dos investimentos exemplifica de um modo brutal a ``economia especulativa de promessas'' que Stengers critica \parencite{stengers2018another}. Eu, no entanto, repetia para mim mesma: é uma empresa, quero dizer, funciona através de lógicas muito distintas das da universidade. 

O encerramento chegou quando eu entrava na fase final de escrita desta tese e obrigou-me a rever algumas questões que atravessam esta pesquisa. No entanto, acabou por ser também um presente etnográfico sobre o registro de uma parceria público-privada controversa, tensa, contraditória e, a seu modo, bem-sucedida na pesquisa em IA em uma universidade pública brasileira. Essa situação me aproximou do estudo de \textcite{Latour1993Aramis} sobre o Aramis, um sistema de transporte que nunca chegou a existir. Latour conduziu sua pesquisa entre dezembro de 1987 e janeiro de 1989 com o intuito de investigar por que o projeto havia sido cancelado em fevereiro de 1987.\footnote{Minha posição temporal distingue-se: iniciei a pesquisa em março de 2022, quando a parceria parecia consolidada; o encerramento em dezembro de 2025 não estava previsto. Registrei processos enquanto ocorriam, capturando tensões que permaneciam latentes porque a rede ainda resistia.} Tanto o Aramis quanto o C4AI são casos de parcerias público-privadas onde atores com interesses diversos mobilizam promessas de longo prazo que se desfazem quando cálculos estratégicos corporativos mudam. Em fevereiro de 1987, a Matra distribuiu uma nota ordenando ``pare tudo''; em agosto de 2025, a IBM comunicou internamente que o convênio não seria renovado. No Aramis, a retirada resultou no cancelamento completo após quase vinte anos de investimentos públicos.\footnote{O que restou do Aramis foi um protótipo no saguão da empresa, uma casca vazia sem trilhos, sem eletricidade, sem passageiros: ``um objeto de museu'' antes mesmo de ter sido um objeto no mundo. Latour pergunta quando uma tecnologia se torna um objeto. ``Nunca'', ele responde, ``exceto quando porções extraídas da instituição são colocadas em exposição dentro de museus técnicos'' \parencite[p. 9]{Latour1993Aramis}. É só no ferro-velho que um objeto técnico finalmente se torna objeto. Antes disso, é sempre uma instituição, uma rede de humanos e não humanos, uma transação em curso. Law radicaliza esse argumento ao propor vocabulário para pensar objetos que nunca alcançam fechamento, investigando ``as vidas incertas e complexas dos objetos em um mundo onde não há fechamento'' \parencite[p. 59]{Law2004After}. Objetos são múltiplos mesmo quando parecem singulares, e essa multiplicidade persiste sem se resolver com estabilização.} No C4AI, a saída da IBM não acabou com o Centro, mas forçou sua reconfiguração como a perda de financiamento corporativo, incertezas sobre projetos em andamento e o esvaziamento das expectativas construídas ao longo de cinco anos. A promessa de dez anos durou cinco. O Centro continua a existir, mas a parceria IBM-C4AI tornou-se objeto do passado; uma relação que deixou rastros em artigos científicos, infraestrutura computacional e promessas não cumpridas. Esta tese começou como uma etnografia de um arranjo sociotécnico aparentemente estável e terminou descrevendo justamente a sua instabilidade e mudanças em um campo volátil e de difícil acesso, e o que isso pode vir a dizer sobre a pesquisa pública em IA no Brasil.

A fotografia da Figura~\ref{fig:dgx1-openai} circulou pela internet depois que Musk a publicou no Twitter em 2016, e tornou-se uma das imagens de referência da história recente da inteligência artificial: Jensen Huang (CEO da Nvidia) entrega a Elon Musk (co-fundador da OpenAI em 2015) um servidor de US\$~129.000 com oito GPUs, e Musk assina o chassi da supermáquina. Ao fundo, uma frase na parede, parcialmente legível, anuncia que o destino do mundo depende do uso que se faz do conhecimento. Menos de uma década depois, a arquitetura de treinamento inaugurada por aquela máquina produziria sistemas capazes de gerar texto, imagem, código e voz com fluência suficiente para redesenhar mercados de trabalho, práticas científicas, sistemas educacionais e disputas geopolíticas. Treinar esses sistemas exige acesso a dezenas, centenas ou milhares de GPUs; a infraestrutura que a cena de 2016 inaugura tornou-se condição de entrada para qualquer laboratório que pretenda produzir modelos competitivos, e a busca por mais poder computacional e mais dados funcionou, por anos, como lei não escrita do desenvolvimento de modelos de linguagem. Em janeiro de 2026, o lançamento do \textit{DeepSeek}, modelo chinês que atingiu desempenho comparável aos sistemas americanos de ponta com apenas uma fração dos recursos computacionais e dos dados, abalou esse paradigma e reabriu a questão sobre quais são, afinal, as condições necessárias para fazer pesquisa em IA. O Spira, projeto desenvolvido no laboratório de processamento de linguagem natural do IME-USP antes mesmo de o C4AI existir como instituição, chegou a essa condição por outra via através de uma doação pontual de GPU da própria Nvidia, registrada nos agradecimentos do primeiro artigo do projeto \parencite{Casanova2021DeepLearning}, sem contrato nem vínculo institucional. O C4AI, por sua vez, operou ao longo de cinco anos com um cluster de GPUs financiado pela USP-Fapes-IBM; em dezembro de 2025, dias após o fim da parceria com a IBM, a USP inaugurou o JAIRU, o maior \textit{cluster} (conjunto de GPUs que trabalham em paralelo) de IA da América Latina, com 96 GPUs Nvidia B200 e um investimento de aproximadamente R\$~40 milhões. As possibilidades que se abriram com aquela máquina são inseparáveis dos problemas que trouxeram, a concentração de poder computacional em poucas corporações, a dependência infraestrutural de países periféricos, a extração de trabalho não remunerado de milhões de usuários cujos dados alimentam os modelos, a opacidade dos sistemas e dificuldade de responsabilização quando erram, entre outros problemas. 

%%%%%%%%%%%%%%%%%%%%%%%%%%
% TRANSIÇÃO PARA TIMELINE
%%%%%%%%%%%%%%%%%%%%%%%%%%

\begin{figure}[!htb]
\centering
\includegraphics[width=1.0\textwidth]{figuras/cap.1/Nvidia-entrega-openai.jpg}
\caption[Entrega do primeiro supercomputador DGX-1 à OpenAI (2016)]{Entrega do primeiro supercomputador DGX-1 à OpenAI (2016). Jensen Huang, CEO da Nvidia, registra o momento em que Elon Musk, co-fundador da OpenAI, assina o chassi do sistema no escritório da OpenAI em São Francisco. O DGX-1, avaliado em US\$~129.000, reunia oito GPUs Tesla P100 com capacidade de até 170 teraflops (FP16), tornando-se a primeira infraestrutura de computação dedicada exclusivamente ao treinamento de redes neurais em escala. Ao fundo, inscrição na parede com trecho parcialmente legível: \textit{Man has a large capacity for [...]. Than we think it is few [...]. Value the faculty of knowing [...]. The sill to do it. Knowing. The critical issue is not [...]. We know ... I believe that [...]. The fate of the worrld.} A fotografia foi compartilhada pelo próprio Musk no Twitter, plataforma que ele viria a adquirir em 2022, renomeando-a como X, e circulou amplamente pela internet. Fonte: Twitter/@elonmusk, 2016.}
\label{fig:dgx1-openai}
\end{figure}

Esta etnografia foi feita durante esses acontecimentos. Iniciei o trabalho de campo no C4AI em março de 2022, seis meses antes do lançamento do ChatGPT, quando a inteligência artificial generativa ainda era objeto de especialistas e o debate público sobre seus efeitos mal havia começado. Entre março de 2022 e dezembro de 2025, assisti, dentro e a partir do Centro, à aceleração que converteu modelos de linguagem em infraestrutura cotidiana, ao crescimento das publicações sobre processamento de linguagem natural pelo C4AI, ao surgimento do primeiro modelo de linguagem brasileiro, à aprovação do primeiro marco regulatório global e ao encerramento da parceria que havia dado origem ao campo que eu estudava. O C4AI, fundado quatro anos depois do DGX-1 com recursos públicos brasileiros e co-financiado pela IBM, é um nó nessa rede e herdeiro das possibilidades, limites e todo o tipo de problema que a fotografia de Jensen Huang e Elon Musk evoca. Acompanhar esse Centro por quase quatro anos é também acompanhar uma pergunta que a fotografia enseja: qual é o papel da universidade pública na produção de conhecimento em inteligência artificial quando quem está nas fronteiras da pesquisa em IA são essas mesmas corporações? Esta tese documenta uma tentativa de resposta situada à essa questão: o que o C4AI fez, como se organizou e o que ficou quando a IBM deixou de co-financiar o centro. 

\begin{figure}[!htb]
\centering
\includegraphics[width=1.0\textwidth,height=0.9\textheight,keepaspectratio]{figuras/cap.1/roadmap-capitulo1-2026-05-01-121716.png}
\caption[Linha do tempo de abrangência da tese (2016--2026)]{Linha do tempo de abrangência da tese (2016--2026), organizada em cinco raias (diamante/losango que represente um marco) paralelas que reúnem os principais eventos que aconteceram durante esta pesquisa. A raia tecnológica registra os marcos de infraestrutura computacional para IA, da entrega do primeiro DGX-1 à OpenAI (dez/2016) e da arquitetura Transformer (jun/2017) ao GPT-3 (jun/2020), ao ChatGPT (nov/2022) e ao LLaMA (fev/2023), passando pelos aceleradores com arquitetura chiplet, AMD MI300X (dez/2023) e NVIDIA Blackwell (mar/2024), que ampliaram a capacidade de treinamento em escala industrial. A raia da pandemia que recobre o intervalo da pandemia de Covid-19, declarada pela OMS em 11 de março de 2020 e encerrada como emergência de saúde pública internacional em 5 de maio de 2023, e situa nesse intervalo a fundação do C4AI (out/2020) e o desenvolvimento do SPIRA. A raia regulatória acompanha os atos de política científica e de IA: a Estratégia Brasileira de IA (abr/2021), o CHIPS Act estadunidense (ago/2022), o AI Act europeu (mar/2024), o Plano Brasileiro de IA (jul/2024) e a Portaria CNPq 2.664 sobre uso de IA no ensino superior (mar/2026). A raia bibliométrica acompanha a aceleração das publicações sobre IA nas ciências humanas brasileiras a partir de 2020, com base nas SciELO e CAPES, e situa o ano de 2023 como ápice da produção do C4AI. A raia etnográfica reúne os episódios do trabalho de campo, do início em março de 2022 ao encerramento em dezembro de 2025, passando pela \textit{bluetalks} em (fev/2023), pela entrevista com Marcelo Finger (jun/2023), pela qualificação da tese (fev/2024), pelas entrevistas com Fábio Cozman (set/2025) e Cláudio Pinhanez (out/2025), até o encerramento da parceria C4AI-IBM (dez/2025). Versão interativa: \url{https://mermaid.ai/d/b27ca78c-b037-42fe-a3eb-939a092e1455}.}
\label{fig:timeline-cap1}
\end{figure}

Estar no campo durante esse período implicou um desafio que vai além do habitual esforço de atualização bibliográfica. Nas seções que se seguem, detalho os fundamentos teórico-metodológicos que orientaram essa trajetória, as operações de triangulação que permitiram articular dados heterogêneos, os cortes analíticos estratégicos que transformaram redes infinitas em objetos etnográficos analisáveis, e as quatro lições que o encontro entre teoria e campo produziu como contribuições metodológicas desta tese. A etnografia, como o caracol, avançou devagar, secretando seu rastro. O que esse rastro deixou está escrito (e inscrito) nesta tese. Mas, esse rastro também deixa meus vestígios enquanto pesquisadora aprendendo a caminhar sobre superfícies estranhas, e deixa, sobretudo, os vestígios dos próprios pesquisadores do C4AI construindo redes, mobilizando inscrições e recrutando aliados, negociando com corporações, fazendo, por fim, \textit{tecnociência}.\footnote{Em \textit{Ciência em Ação}, \textcite{Latour2011CienciaEmAcao} mobiliza o termo para nomear um campo em que ciência e tecnologia se produzem como um só processo, sem partição prévia entre fatos e artefatos, descobertas e invenções, conhecimento e instrumentação. O argumento opera contra a separação que isola o trabalho conceitual dos cientistas do trabalho material dos engenheiros, pois, para Latour,  as controvérsias científicas se resolvem pela capacidade de mobilizar redes longas e heterogêneas, nas quais, teoremas, máquinas, documentos, financiamentos e alianças institucionais atuam como aliados que tornam o custo da discordância proibitivo para o oponente. A tecnociência aparece, nesse sentido, como prática coletiva de fechamento de caixas-pretas, ou seja, de estabilização de associações cuja construção desaparece da vista quando o fato passa a circular como inquestionável. \textcite{stengers2018another} mobiliza o mesmo termo em outra direção. Para Stengers, a tecnociência é a forma assumida pelas práticas científicas quando recrutadas pela lógica industrial e pelos imperativos de competitividade, sob o regime do que chama de \textit{fast science}, em que a velocidade da produção e a captura por interesses corporativos reorganizam as condições de possibilidade do trabalho científico, e em que a formação do pesquisador opera como um tipo de anestesia que o torna insensível à recalcitrância dos mundos atravessados pela pesquisa. Stengers renomeia a economia do conhecimento como \textit{economia especulativa de promessas}, regime em que o financiamento se sustenta pela produção de possibilidades orientadas a investidores futuros e em que a avaliação crítica entre pares se enfraquece porque ninguém quer cortar o galho da promessa que sustenta o campo. Latour e Stengers descrevem dimensões diferentes do mesmo fenômeno: o vocabulário latouriano permite mapear como redes tecnocientíficas se estendem e se estabilizam; o diagnóstico de Stengers expõe como essas redes, sob condições contemporâneas, operam regimes específicos de captura, aceleração e despolitização. Esta tese mobiliza os dois registros. O aspecto interessante na proposta de Stengers é que ela consegue articular, ao mesmo tempo, a crítica (não no sentido de simples denúncia) e a construção de alianças, lembrando que sua formação é em química e filosofia da ciência. A trajetória da parceria C4AI-IBM, com a promessa inicial de dez anos reduzida a cinco quando os cálculos corporativos se reorganizaram, é caso empírico em que as duas leituras se sobrepõem: a rede longa que Latour permite descrever é também a economia de promessas que Stengers diagnostica. Estudá-las em ação, antes do fechamento que tentam estabilizar e da retirada que pode desfazê-las, é o gesto que orienta a primeira regra metodológica reproduzida adiante e a orientação desta tese.} É desse rastro que trato agora. Para isso, começo por mobilizar a noção de \textit{actante}.\footnote{O termo é tomado da semiótica de Algirdas Julien Greimas, que o utilizava para descrever qualquer entidade dotada de papel narrativo em um relato, sem prejulgamento sobre sua natureza humana ou não humana, individual ou coletiva, figurativa ou abstrata. Latour transporta essa generalidade para os estudos de ciência e tecnologia precisamente para escapar do antropocentrismo embutido na palavra \enquote{ator}, que tende a pressupor intenção, consciência e personificação. Em \textit{Ciência em Ação}, oferece a formulação mínima: \enquote{chamarei de actante qualquer pessoa e qualquer coisa que seja representada}, observando que não há diferença prática, para fins de análise, entre representar trabalhadores num sindicato e representar hormônios num experimento, pois uns e outros são mudos e dependem de porta-vozes que articulem suas posições \parencite[p. 127, 138-140]{Latour2011CienciaEmAcao}. Em \textit{A Esperança de Pandora}, a definição é refinada em chave pragmática: um actante é definido pela lista de provas a que é submetido e pelas performances que apresenta nessas provas; sua \enquote{competência} ou essência é deduzida só depois, retrospectivamente, a partir do que conseguiu fazer ou impedir de acontecer \parencite[Glossário]{Latour2017Esperanca}. Um microrganismo, por exemplo, antes de receber nome e estabilizar-se como entidade reconhecível, é apenas a soma das vitórias e derrotas que registra ao transformar açúcar em álcool, resistir à fervura, sobreviver a antissépticos. Em \textit{Políticas da Natureza}, a ênfase desloca-se para a relacionalidade: actante é \enquote{qualquer entidade que modifica outra entidade em uma prova} \parencite[Glossário]{latour2004politics}. A agência, nesse registro, nunca é propriedade isolada de uma entidade; é \textit{fait-faire}, fazer-fazer, distribuído entre os mediadores da rede. No artigo sobre o dispositivo de fechamento automático de portas, Latour acrescenta uma distinção que importa para a descrição etnográfica: actantes podem ser figurativos, quando aparecem personificados (um funcionário, um zelador), ou não figurativos, quando operam de modo técnico ou abstrato (uma mola, um quebra-molas, um algoritmo) \parencite{johnson1988mixing}. O que define o actante é, nesse sentido, o que ele faz na cadeia de associações, não o tipo de entidade que parece ser. Esses quatro registros são complementares e operam juntos nesta tese: actante é qualquer entidade que (i) é representada por um porta-voz, (ii) é definida pragmaticamente pelo que faz nas provas a que é submetida, (iii) modifica outras entidades nessas provas, e (iv) pode aparecer sob figura humana ou sob forma técnica sem que isso altere seu estatuto analítico. Convém marcar uma cautela. Tratar humanos e não humanos como actantes não significa afirmar sua equivalência ontológica nem dissolver as diferenças entre seus modos de agir; significa, mais recusar a partição prévia que decidiria, antes da pesquisa, quem age e quem é apenas cenário ou contexto. Os actantes que importam analiticamente nesta tese são, em sua maioria, entidades que transformam, traduzem e tensionam aquilo que atravessam, em contraste com os intermediários, que apenas transportam algo ou alguma coisa sem alteração. Contratos, modelos computacionais, GPUs, vozes gravadas em enfermarias, o vírus SARS-CoV-2 e linguagens formais entram nas redes que descrevo porque deslocam o curso das associações e das ações em que participam. E tornam possíveis fazer coisas que antes era impossível ou que ocorriam de maneira diferente.} Esse conceito sustenta a noção de simetria teórico-metodológica que, em geral, orienta a teoria ator-rede. Considerando a interpretação de um de seus principais formuladores, Bruno Latour, por exemplo, argumenta que não devemos aceitar acriticamente o que os cientistas dizem sobre a Natureza, tampouco o que os cientistas sociais dizem sobre a Sociedade. A sua quarta regra metodológica estabelece que ``uma vez que a resolução de uma controvérsia é a causa da estabilidade da sociedade, não podemos usar a sociedade para explicar como e por que uma controvérsia foi resolvida. Devemos considerar simetricamente os esforços para alistar e controlar recursos humanos e não-humanos'' \parencite[p. 225]{Latour2011CienciaEmAcao}. Isso desloca Natureza e Sociedade da posição de causas explicativas do desfecho das controvérsias para a posição de consequências estabilizadas ao final do processo. Para esta tese, essa regra tem implicação metodológica direta: não posso explicar a pesquisa em inteligência artificial no C4AI recorrendo a fatores sociais externos como se fossem mais estáveis ou mais reais do que os algoritmos, os dados ou os modelos. Em outras palavras, devo observar empiricamente quem ou o que efetivamente \enquote{age} em determinada situação, não esquecendo também de como e para quê no sentido de Haraway e Stengers. 

O desafio desse tipo de posicionamento analítico consiste em se apoiar no arcabouço teórico-metodológico das ciências sociais e, simultaneamente, voltar esses mesmos conceitos, centrais às investigações sociológicas e antropológicas, para a própria análise. Isso exige uma suspensão recorrente das nossas próprias crenças enquanto cientistas sociais, suspensão que precisa ser refeita a cada nova passagem do texto e que se torna condição para o aprendizado, para um aprender antropologicamente no sentido proposto por \textcite{Ingold2019Antropologia}. É nesse ponto que a análise alcança sua maior densidade, porque os actantes mobilizados nesta tese vão de políticas científicas e contratos institucionais a GPUs, algoritmos, vozes humanas gravadas em enfermarias, inscrições tecnocientíficas dos artigos (imagens microscópicas, diagramas, gráficos, equações) e o próprio vírus SARS-CoV-2. Os Capítulos 3 e 4 documentam como esses actantes se associam em cadeias heterogêneas que dão forma a redes sociotécnicas. Isso ocorre tanto em reuniões de \textit{design thinking}, como meus interlocutores chamam as conversas que decidem os rumos das pesquisas e suas condições, quanto nas enfermarias hospitalares onde pacientes com Covid-19 tiveram suas falas gravadas e convertidas em dados computacionais para treinamento e inferência do Spira. A pandemia atravessa esta tese: o C4AI foi inaugurado em outubro de 2020, em plena crise sanitária, e o trabalho de campo começou em março de 2022, ainda sob seus efeitos. Seguir os atores levou-me ao vírus SARS-CoV-2, que se mostrou um actante que reconfigurou as condições de possibilidade tanto das pesquisas em IA que observo quanto da pesquisa que realizo. Esse posicionamento teórico-metodológico é, no sentido que \textcite{Barad2007Meeting} dá ao termo, ético-onto-epistemológico.\footnote{Como adiantado no Capítulo~\ref{capitulo0}, Karen Barad cunha a expressão \enquote{ethico-onto-epistem-ology} para nomear a inseparabilidade entre as questões do que existe, do que se pode conhecer e da responsabilidade que se assume ao conhecer e fazer existir: \enquote{conhecer, pensar, medir, teorizar e observar são práticas materiais de intra-ação no mundo e como parte dele} \parencite[p.~90]{Barad2007Meeting}. A escrita com hífens do termo indica que as três dimensões são inseparáveis em qualquer prática de pesquisa. Para o argumento desta tese, isso significa que tratar simetricamente vírus, GPUs, algoritmos e políticas científicas é, ao mesmo tempo, decisão sobre o que existe, decisão sobre como conhecer e decisão sobre quem responde pelas associações que se estabilizam. O Capítulo~\ref{capitulo4} mobiliza essa orientação no plano quando o \textit{pipeline} (representação visual de como foi feito um sistema de IA) do Spira produz, pelo próprio aparato, o objeto que mede, calcula, treina e infere.} Esse esforço, contudo, não está livre de contradições. Em muitos momentos, contradigo a mim mesma e aos princípios que digo seguir. Pensar dessa maneira exige virar uma chave interpretativa que nem sempre conseguimos acionar. Anos de formação nas ciências sociais nos ensinam a buscar o social nos fenômenos, a desconfiar das explicações técnicas, a revelar os interesses ocultos. Desaprender esse costume acadêmico é um trabalho contínuo e necessariamente incompleto. Manter essa disciplina (não poderia ser um tipo de recursividade?) ao longo desta tese é uma prova de força (e não estou falando da mesma prova de força no sentido latouriano, apesar de também ter a ver com isso). Por isso, quando o leitor encontrar passagens em que recorro apressadamente a fatores sociais para explicar o que deveria ser descrito em sua heterogeneidade, saiba que estou mais ou menos ciente desses deslizes. 

A formulação ético-onto-epistemológica de Barad encontra o seu antecedente operacional, por assim dizer, nas sete regras metodológicas e nos seis princípios que Latour propõe ao final de \textit{Ciência em Ação}. Onde Barad nomeia a inseparabilidade entre o que existe, o que se pode conhecer e a responsabilidade que se assume ao conhecer, Latour elaborou um guia prático para sustentar essa inseparabilidade ao longo de uma pesquisa: o que considerar como objeto, em que momento entrar no campo, como tratar simetricamente humanos e não humanos, o que fazer diante de divisores entre interior e exterior. Reproduzo-os aqui na íntegra porque cada um deles intervém no texto desta tese como decisão tomada a cada novo capítulo, a cada nova entrevista, a cada decisão sobre qual associação seguir e qual deixar fora da análise.

\begin{citacaoabnt}
\textit{Regra 1.} Estudamos a ciência em ação, e não a ciência ou tecnologia pronta; para isso, ou chegamos antes que fatos e máquinas se tenham transformado em caixas-pretas, ou acompanhamos as controvérsias que as reabrem.

\textit{Regra 2.} Para determinar a objetividade ou subjetividade de uma afirmação, a eficiência ou a perfeição de um mecanismo, não devemos procurar por suas qualidades intrínsecas, mas por todas as transformações que ele sofre depois, nas mãos dos outros.

\textit{Regra 3.} Como a solução de uma controvérsia é a causa da representação da Natureza, e não sua consequência, nunca podemos utilizar essa consequência, a Natureza, para explicar como e por que uma controvérsia foi resolvida.

\textit{Regra 4.} Como a resolução de uma controvérsia é a causa da estabilidade da sociedade, não podemos usar a sociedade para explicar como e por que uma controvérsia foi dirimida. Devemos considerar simetricamente os esforços para alistar recursos humanos e não-humanos.

\textit{Regra 5.} Com relação àquilo de que é feita a tecnociência, devemos permanecer tão indecisos quanto os vários atores que seguimos; sempre que se constrói um divisor entre interior e exterior, devemos estudar os dois lados simultaneamente e fazer uma lista (não importa se longa e heterogênea) daqueles que realmente trabalham.

\textit{Regra 6.} Diante da acusação de irracionalidade, não olhamos para que regra da lógica foi infringida nem que estrutura social poderia explicar a distorção, mas sim para o ângulo e a direção do deslocamento do observador, bem como para a extensão da rede que assim está sendo construída.

\textit{Regra 7.} Antes de atribuir qualquer qualidade especial à mente ou ao método das pessoas, examinemos os muitos modos como as inscrições são coligidas, combinadas, interligadas e devolvidas. Só se alguma coisa ficar sem explicação depois do estudo da rede é que deveremos começar a falar em fatores cognitivos \parencite[p. 420-422]{Latour2011CienciaEmAcao}.
\end{citacaoabnt}

Se as regras orientam a prática etnográfica, os seis princípios extraem dela consequências teóricas mais amplas sobre a natureza dos fatos, das máquinas e das redes que os sustentam:

\begin{citacaoabnt}
\textit{Primeiro princípio.} O destino de fatos e máquinas está nas mãos dos consumidores finais; suas qualidades, portanto, são consequência, e não causa, de uma ação coletiva.

\textit{Segundo princípio.} Os cientistas e engenheiros falam em nome de novos aliados que conformaram e alistaram; representantes entre outros representantes, com esses recursos inesperados, fazem o fiel da balança de forças pender em seu favor.

\textit{Terceiro princípio.} Nunca somos postos diante da ciência, da tecnologia e da sociedade, mas sim diante de uma gama de associações mais fracas e mais fortes; portanto, entender o que são fatos e máquinas é o mesmo que entender o que as pessoas são.

\textit{Quarto princípio.} Quanto mais esotérico o conteúdo da ciência e da tecnologia, mais elas se expandem externamente; portanto, 'ciência e tecnologia' é apenas um subconjunto da tecnociência.

\textit{Quinto princípio.} A acusação de irracionalidade é sempre feita por alguém que está construindo uma rede em relação a outra pessoa que atravessa seu caminho; portanto, não há Grande Divisor entre mentes, mas apenas redes maiores ou menores.

\textit{Sexto princípio.} A história da tecnociência é, em grande parte, a história dos recursos espalhados ao longo das redes para acelerar a mobilidade, a fidedignidade, a combinação e a coesão dos traçados que possibilitam a ação a distância \parencite[p. 422-424]{Latour2011CienciaEmAcao}.
\end{citacaoabnt}

Conforme desenvolvia esta pesquisa, orientada por esses princípios e regras, percebi que, ao fazer associações com os pesquisadores do C4AI, tanto eu quanto esta tese passamos a fazer parte das redes do Centro, e tanto eu quanto a tese nos tornamos actantes nessas redes; assim, esta tese tornou-se também uma inscrição tecnocientífica.\footnote{A percepção de que esta tese também fazia redes me foi sugerida por Fernanda Rosa, em conversa após eu apresentar as primeiras impressões desta pesquisa no encontro anual da \textit{Society for Social Studies of Science} (4S), em Cholula, em dezembro de 2022. Fernanda Rosa investiga infraestruturas digitais no Sul Global e aparece também na revisão do Capítulo~\ref{capitulo2}. A formulação dela deslocou a posição em que eu me colocava até então: o que eu vinha descrevendo como acompanhamento dos cientistas e engenheiros do C4AI também era construção de associações que iam ganhando forma na medida em que as registrava. Levei tempo para perceber o alcance metodológico daquela formulação, e foi essa percepção que veio a desenhar o eixo desta tese.} Latour define instrumento como qualquer estrutura que possibilite uma exposição visual de qualquer tipo em um texto científico \parencite[p.~102]{Latour2011CienciaEmAcao}, e inscrição como aquilo que esse instrumento produz: o traçado, o gráfico, a tabela, o registro.\footnote{O conceito de inscrição é desenvolvido em três frentes. Em \textit{Ciência em Ação} \parencite[caps.~1--2]{Latour2011CienciaEmAcao}, Latour analisa o artigo científico como construção sobre camadas de inscrições e descreve o laboratório como fábrica de inscrições; em \textit{A Esperança de Pandora} \parencite[cap.~2]{Latour2017Esperanca}, acompanha uma expedição pedológica à Amazônia e mostra como o solo se transforma sucessivamente em amostras, em cores de um código de Munsell, em números e em um diagrama de relatório, cadeia que ele chama de referência circulante; em \enquote{Visualization and Cognition} \parencite{Latour1986Visualization}, define as inscrições como \textit{móveis imutáveis}, condição da dominação à distância. Para que o conhecimento se acumule em centros de cálculo, as inscrições precisam ser móveis, transportáveis dos lugares distantes para o centro; imutáveis, sem distorção no transporte; e combináveis, rearranjáveis em uma única superfície plana. A construção de fatos opera por uma cascata de simplificações: uma substância complexa se reduz a um traço, que se reduz a um número, que se torna média em uma tabela, até chegar a uma afirmação dura em um artigo. No texto científico, a inscrição traz o referente para dentro da narrativa: o autor aponta para a figura e mostra a prova, em vez de pedir que o leitor acredite em sua palavra.} Esta etnografia mobiliza uma série dessas inscrições por meio de imagens, diagramas, gráficos e do site que apresenta e performa esta tese. E quando digo que a tese em si é uma inscrição tecnocientífica, é porque ela possibilita certos tipos de ações. De alguma maneira, traduzi aquele pedaço do mundo envolvendo o C4AI para o formato de papel, ou melhor, para um arquivo digital que começou como anotações em \texttt{.txt} no bloco de notas do celular, passou por um \texttt{.docx} no Word e no Google Docs, depois para LaTeX no Overleaf, até chegar a um PDF e, talvez, a um volume encadernado ou a um livro impresso.\footnote{Diferentemente de um texto impresso em papel (que também depende de uma complexa cadeia de técnicas e tecnologias, afinal, até a escrita à mão exige algum suporte físico e um conjunto de habilidades aprendidas ao longo do tempo), um arquivo digital só pode ser compreendido por um dispositivo, como um computador, para que possamos ver as letras na tela. Isso pode parecer um detalhe simples, mas envolve diversas mediações técnicas: o computador lê os diferentes formatos de arquivo convertendo os dados em sequências de bits (0s e 1s, desligado/ligado, falso/verdadeiro) e transformando-os em informações compreensíveis, em um processo que passa pelo sistema operacional, pelas extensões de arquivo e pelos aplicativos específicos instalados na máquina. Sem um software de aplicação que funcione como tradutor dos dados brutos contidos nos arquivos e os transforme em algo compreensível para o usuário, os dados de arquivos como o desta tese em PDF permanecem apenas como dados crus, isto é, mera eletricidade; e, sem a ação do software, essa eletricidade não pode ser convertida em uma mensagem ou informação inteligível. Por exemplo, o código binário 01000001, para um processador de texto, corresponde à letra \enquote{A}. Não é impressionante o que a computação tornou possível? Quando consideramos os modelos de linguagem, tudo isso fica ainda mais interessante, porque esses modelos funcionam em uma escala de indeterminação, ao contrário de qualquer outro sistema computacional tradicional, e aliás contra as suas próprias bases fundacionais que opera segundo lógicas de determinação e continuidade. Arrisco sugerir que é justamente essa indeterminação e descontinuidade que torna a inteligência artificial tão instigante para a antropologia. Por mais que se tente controlar tais aspectos por meio dos famosos \textit{prompts} (como ficaram conhecidas as instruções que são feitas à IA generativa), eles não serão totalmente eliminados. O \textit{prompt} expressa, em certo sentido, a ilusão frustrada de que seria possível controlar integralmente esses sistemas. Isso não significa que esses sistemas sejam incontroláveis, mas talvez os \textit{prompts} não sejam a melhor forma de fazê-lo. Eu mesma realizei experimentações e considerei o resultado tão bizarro em tentar impedir que o sistema falhasse, cometesse erros ou reproduzisse suas marcas estilísticas características; pareceu-me que a IA foi assim \enquote{desumanizada} de uma vez por todas.} Esse recorte dizia respeito a algo muito maior do que eu, no qual se desenrolavam acontecimentos que só podiam ser notados a partir de uma certa escala de análise, incluindo não humanos invisíveis ao olhar humano sem a mediação da tecnociência, como o coronavírus. Ao mobilizar inscrições de diversos tipos, a tese transforma eventos dispersos no tempo e no espaço em traçados estáveis, combináveis e transportáveis. À medida que essa rede se expande para além do Centro, associando actantes cada vez mais inesperados e heterogêneos, a tese adquire o status do que Latour denomina inscrição também literária. É essa inscrição que permite à tecnociência agir à distância e converter fenômenos complexos, remotos ou invisíveis em um arquivo digital que pode ser lido, interpretado, analisado, isto é, dominado, medido, controlado, utilizado tanto para apreciação quanto para persuadir outros acerca da realidade de um fato tecnocientífico cujas cadeias de associações ganham força conforme a extensão da rede e o número de elementos que ela consegue associar \parencite{Latour2011CienciaEmAcao}, ou fazer alianças \parencite{stengers2010cosmopolitics1}.

De certo modo, em sintonia com as críticas de Haraway e Stengers, Latour também formula uma crítica à tecnociência e aos métodos (sejam das ciências sociais ou não) ao afirmar que a própria escolha do que investigamos já constitui um ato político de composição de mundos comuns, no sentido desenvolvido em \textit{Políticas da Natureza} \parencite[p.~18]{latour2004politics}. Em \textit{After Method}, Law nomeia esse trabalho de \textit{method assemblage} por meio do qual o método entendido como produção contínua de fronteiras entre uma presença condensada \emph{in-here} (a representação, o objeto, o resultado), uma ausência manifesta reconhecida como contexto, e uma \textit{otherness} que permanece apagada para que a presença se sustente \parencite[p.~84]{Law2004After}. Essa produção opera a partir de um \textit{hinterland} de relações materiais, instrumentais e institucionais já estabilizadas, que condiciona o que cada método pode fazer existir \parencite[pp.~33--42]{Law2004After}. Por isso, Law nomeia esse processo como um todo de \textit{ontological politics}, pois, o método não é somente um conjunto de procedimentos que reporta uma realidade já existente, mas também produz realidades \parencite[p.~143]{Law2004After}. Segundo esses autores, os métodos não apenas representam, mas produzem realidades; em outras palavras, eles têm a capacidade de ordenar elementos, criar hierarquias, incluir e excluir (o que equivale a fazer política por meio da tecnociência). Se, portanto, os métodos contribuem para gerar realidades, ou performam, enactam (do inglês \textit{enactment}) realidades como um processo constante, fluido e múltiplo.\footnote{Os termos \textit{enact}/\textit{perform} podem confundir. No glossário de \textit{After Method}, Law afirma que \textit{enactment} é sinônimo aproximado de \textit{performance} e justifica a preferência pelo primeiro: o termo \textit{performance} carrega as ressonâncias da sociologia dramatúrgica de Erving Goffman, em que as apresentações pressupõem uma realidade oculta atrás delas, ao passo que \textit{enactment} mantém a prática aberta, sem fechamento e sem instância anterior \parencite[p.~159]{Law2004After}. Mol formula a distinção operativa em \textit{The Body Multiple} ao contrastar \textit{construção}, que sugere estabilização gradual rumo a identidades fixas, e \textit{enactment}, em que a multiplicidade não converge para singularidade \parencite[pp.~32--33, 42--44]{mol2002body}. Adoto as interpretações dos dois autores: \textit{enactuar} e \textit{enactuação} designam o trabalho contínuo pelo qual as práticas fazem existir realidades, ao passo que \textit{performativo} e \textit{performatividade} qualificam a propriedade produtiva dessas práticas, na linha em que Law escreve que o método é performativo \parencite[p.~143]{Law2004After}. Mantenho \textit{enactuação} sempre que o argumento ontológico está em questão, e reservo \textit{performance} para descrever a propriedade do método e da inscrição de produzir aquilo que parecem apenas relatar.} Quer dizer que esta tese não só descreve o C4AI, como também toma parte na produção de uma determinada realidade do Centro, ao mesmo tempo em que deixa outras de fora \parencite[p.~143]{Law2004After}. E, assim, a questão etnográfica também perpassa por \textit{qual ou quais realidades decido reforçar e quais enfraquecer?}

Annemarie Mol dá atenção especial à ontologia. Em \textit{The Body Multiple}, ela investiga a aterosclerose (doença inflamatórias que atinge as artérias humanas) em um hospital e mostra que a doença que o cirurgião vascular opera, a que o patologista examina ao microscópio e a que o clínico geral descreve ao paciente são objetos distintos, performados por práticas distintas, sem que haja uma aterosclerose singular que esses sítios viessem traduzir de ângulos diferentes \parencite{mol2002body}. Para Mol, a realidade não precede as práticas mundanas, mas é moldada dentro dessas práticas \parencite[p.~6]{mol2002body}. O argumento de Mol sugestiona que o C4AI enactado por esta tese coexiste com o C4AI que Fábio e Cláudio enacatam em apresentações públicas e entrevistas, com o C4AI que a FAPESP enacta em relatórios internos, com o C4AI que Marcelo enacta em um sistema de inteligência artificial e com tantos outros que cada prática situada faz existir. Cada uma dessas enactações peformam o C4AI. A política ontológica que Mol e Law formulam \parencite{mol1999ontological,Law2004After} consiste em assumir que cada uma dessas enactações/perfomances fortalece certas relações/conexões em detrimento de outras, complicando o que entendemos por observação, descrição e representação.

Os \enquote{novos materialismos} complicam mais um pouco.\footnote{Os novos materialismos designam um conjunto de perspectivas teóricas que emergem nas décadas de 1990 e 2000 nas humanidades em resposta ao que seus proponentes identificam como excessivo privilégio do discurso e da representação nas teorias pós-estruturalistas. O movimento partilha a recusa da concepção cartesiana de matéria como substância inerte e reconhecem que a materialidade possui modos próprios (ou imanentes) de auto-organização, transformação e agência que não se reduzem à ação de sujeitos humanos sobre ela \parencite[pp.~7--10]{CooleFrost2010}. O Capítulo~\ref{capitulo4} retoma essa discussão na análise do espectrograma mel do SPIRA, articulando-a às tensões com a teoria ator-rede.} A proposta do \enquote{realismo agencial} de Karen Barad substitui a noção de interação, que pressupõe entidades preexistentes que se encontram, pela de \textit{intra-ação}, por meio da qual, sujeito e objeto emergem juntos da prática material-discursiva que os produz \parencite[pp.~54--56]{DolphijnTuin2012}. O que separa o observador de o observado, dentro de fora, pesquisadora de campo, é o que Barad chama de \textit{corte agencial}, ou seja, o efeito da prática que parecia pressuposto da prática. Pensando as implicações desse argumento nesta etnografia, a questão sobre \enquote{o que existe lá fora} abre-se para o C4AI que esta tese descreve não preexiste ao trabalho de campo que o constituiu como objeto descritível, e eu enquanto pesquisadora que o descreve constituo-me como tal no mesmo gesto que constitui o Centro como meu interlocutor, e assim também ocorre com os sistemas de IA que aparecem nesta etnografia (o Claude, o Spira). De modo que, \enquote{lá fora} e \enquote{aqui dentro} são efeitos do corte agencial, e o gesto que parece separá-los é o mesmo gesto que os faz emergir como termos de uma relação. A tensão com a teoria ator-rede persiste, pois, Latour pressupõe actantes que existem antes de se associar, ao passo que Barad faz emergir os actantes da própria intra-ação (ou da relação, pensando a partir de Marilyn Strathern e do fazer-com/simpoiese a partir de Donna Haraway). A análise que proponho, dessa forma, sustenta as duas formulações sem confundi-las, e o capítulo~\ref{capitulo4} desdobra essa fricção na descrição da cadeia de translação que vai da voz de pessoas hospitalizadas durante a pandemia ao espectrograma mel (a imagem da voz apresentado como um um mapa de calor onde as cores representam a intensidade da amplitude do sinal de áudio em termos de tempo e frequência), e por fim, ao sistema de inteligência artificial Spira. Mas, a beleza disso tudo é que, voltando a Latour, \enquote{ser objetivo significa que, sejam quais forem os esforços dos discordantes para romper os elos entre o representante e aquilo em nome do que ele fala, os elos resistirão} \parencite[p.~118]{Latour2011CienciaEmAcao}. Talvez esta tese possa fazer algumas ranhuras em algum desses elos. 

%%%%%%%

% =====================================================================
% CAPÍTULO 1 — Seção "Etnografias: fazendo o método no caminho"
% Versão revisada após Bloco 1 do checklist da orientadora (28/04/2026)
% Ancorada em KOFES, S. Pesquisa antropológica como método e/ou como
% experiência. In: Dilemas na maçonaria contemporânea: um experimento
% antropológico. Campinas: Editora Unicamp, 2015. pp. 291–303.
%
% Substitui o trecho da seção 1.1 do arquivo atual desde
% "A metodologia desta tese fundamenta-se na etnografia..."
% até "...condição de possibilidade para qualquer observação que
% pretendesse ir além da superfície."
% =====================================================================
% =====================================================================
% Cap. 1 — Seção "Etnografias: fazendo o método ao caminhar"
% Versão revisada com base na leitura direta do PDF de Kofes (2015,
% pp. 291-303) e nos comentários do Bloco 1 do checklist da orientadora.
%
% Alterações principais (mapeadas em prosa abaixo):
%
% 1. PARÁGRAFO CENTRAL: reformulado para fidelidade ao texto de Kofes.
%    - "narração" → "narrativa" em todas as ocorrências (Kofes não usa
%      "narração"; usa "narrativa" e o verbo "narrar").
%    - Removida a formulação interpretativa "etnografia heterodoxa
%      como prática que envolve as dimensões de totalidade, trabalho
%      de campo e escritura sem se reduzir a nenhuma delas" — essa
%      síntese não está nas pp. 291-303. Substituída pelo que Kofes
%      efetivamente formula nas páginas referenciadas.
%    - "tomada inicial e a retomada teórica" → "tomada inicial e sua
%      retomada teórica" (com possessivo, fiel a Kofes).
%    - "evidência do que é" → "evidência do que o é" (com pronome,
%      fiel a Kofes p. 297).
%    - Separadas as duas referências a Strathern: "recriação imaginativa"
%      (p. 299) e "prática reveladora / revelatória" (p. 297).
%    - "tensão entre método e experiência pode ser superada" →
%      "se torna desnecessária quando dissolvida" (formulação literal
%      de Kofes p. 301).
%    - "estrutura a experiência" → reorganizada a relação entre
%      "structured unit of experience" e "estrutura de experiência",
%      respeitando a distinção que Kofes faz na p. 301.
%
% 2. NOTA 2 (Favret-Saada): ATENÇÃO — a chave bibtex usada duas vezes,
%    \parencite{FavretSaada2005SerAfetado}, está incorreta para a
%    primeira ocorrência (que se refere ao ensaio em francês).
%    Substituída por \parencite{FavretSaada1990SerAfetado} — entrada
%    bibtex sugerida no resumo abaixo.
%
% 3. NOTA 3 (Law): removida a mesma formulação interpretativa sobre
%    "dimensões de totalidade, trabalho de campo e escritura".
%
% 4. PARÁGRAFO 3 (etnobiografia): reformulações pontuais de estilo.
%
% 5. PARÁGRAFO 4 e seguintes: correções pontuais (sem alterações
%    estruturais — a parte sobre Digiampietri, Latour-Woolgar, Law,
%    Mol e cortes-de-rede está mantida intacta, exceto por
%    desambiguações pontuais e correções tipográficas).
% =====================================================================

% =====================================================================
% Cap. 1 — Seção "Etnografias: fazendo o método ao caminhar"
% Versão revisada com base na leitura direta do PDF de Kofes (2015,
% pp. 291-303) e nos comentários do Bloco 1 do checklist da orientadora.
%
% Alterações principais (mapeadas em prosa abaixo):
%
% 1. PARÁGRAFO CENTRAL: reformulado para fidelidade ao texto de Kofes.
%    - "narração" → "narrativa" em todas as ocorrências (Kofes não usa
%      "narração"; usa "narrativa" e o verbo "narrar").
%    - Removida a formulação interpretativa "etnografia heterodoxa
%      como prática que envolve as dimensões de totalidade, trabalho
%      de campo e escritura sem se reduzir a nenhuma delas" — essa
%      síntese não está nas pp. 291-303. Substituída pelo que Kofes
%      efetivamente formula nas páginas referenciadas.
%    - "tomada inicial e a retomada teórica" → "tomada inicial e sua
%      retomada teórica" (com possessivo, fiel a Kofes).
%    - "evidência do que é" → "evidência do que o é" (com pronome,
%      fiel a Kofes p. 297).
%    - Separadas as duas referências a Strathern: "recriação imaginativa"
%      (p. 299) e "prática reveladora / revelatória" (p. 297).
%    - "tensão entre método e experiência pode ser superada" →
%      "se torna desnecessária quando dissolvida" (formulação literal
%      de Kofes p. 301).
%    - "estrutura a experiência" → reorganizada a relação entre
%      "structured unit of experience" e "estrutura de experiência",
%      respeitando a distinção que Kofes faz na p. 301.
%
% 2. NOTA 2 (Favret-Saada): ATENÇÃO — a chave bibtex usada duas vezes,
%    \parencite{FavretSaada2005SerAfetado}, está incorreta para a
%    primeira ocorrência (que se refere ao ensaio em francês).
%    Substituída por \parencite{FavretSaada1990SerAfetado} — entrada
%    bibtex sugerida no resumo abaixo.
%
% 3. NOTA 3 (Law): removida a mesma formulação interpretativa sobre
%    "dimensões de totalidade, trabalho de campo e escritura".
%
% 4. PARÁGRAFO 3 (etnobiografia): reformulações pontuais de estilo.
%
% 5. PARÁGRAFO 4 e seguintes: correções pontuais (sem alterações
%    estruturais — a parte sobre Digiampietri, Latour-Woolgar, Law,
%    Mol e cortes-de-rede está mantida intacta, exceto por
%    desambiguações pontuais e correções tipográficas).
% =====================================================================

\section{Etnografias: fazendo o método ao caminhar}
\label{sec:etnografia-heterodoxa}

\epigrafe{%
Há muitos métodos para o estudo da construção de fatos científicos e artefatos técnicos. A primeira regra metodológica pela qual nos decidimos na Introdução é a mais simples de todas. Não tentaremos analisar os produtos finais; em vez disso, seguiremos os passos de cientistas e engenheiros nos momentos e nos lugares nos quais planejam, desfazem, modificam. Vamos dos produtos finais à produção, de objetos instáveis e frios a objetos instáveis e mais quentes. Estar ali \textbf{antes} que a caixa se fechasse e ficasse preta. Com esse método simples precisamos apenas seguir os melhores de todos os guias, os próprios cientistas, em sua tentativa de fechar uma caixa-preta e abrir outra.}%
{Bruno Latour, \textit{A ciência em ação: como seguir cientistas e engenheiros sociedade afora}, 2011.}

Esta pesquisa fundamenta-se em uma acepção específica de etnografia elaborada por Suely Kofes.\footnote{A acepção de etnografia mobilizada nesta tese é a que Suely Kofes desenvolve em \textit{Dilemas na maçonaria contemporânea: um experimento antropológico} \parencite{Kofes2015Maconaria}. A formulação aparece em duas seções complementares do livro: na parte dois, intitulada \enquote{Etnografia heterodoxa}, e na seção final, \enquote{Pesquisa antropológica como método e/ou como experiência} \parencite[pp.~291--303]{Kofes2015Maconaria}. Nessas seções, Kofes articula o entrelaçamento entre biografia, etnografia e ciências sociais como alternativa às categorizações dicotômicas supostamente universais. Essa acepção dialoga com obra anterior da autora, \textit{Uma trajetória, em narrativas} \parencite{kofes2001trajetoria}, na qual Kofes contrasta sua proposta com outras definições consolidadas de etnografia: a etnografia como totalidade na formulação de Bronislaw Malinowski em \textit{Argonautas do Pacífico Ocidental}, como trabalho de campo na formulação de Martyn Hammersley e Paul Atkinson em \textit{Ethnography: Principles in Practice}, e como escritura na formulação de James Clifford e George Marcus em \textit{Writing Culture: The Poetics and Politics of Ethnography}. Kofes observa que nenhuma dessas definições dá conta isoladamente de explicar o que faz uma etnógrafa.} A proposta de Kofes é pensar a pesquisa antropológica como método e como experiência, formulação que ela desenvolve no capítulo final do livro de 2015 ao fazer dialogar Strathern, Favret-Saada e Turner. Kofes parte de uma desconfiança sobre a verdade do empirismo como método e da consciência dos limites do método: o que está em jogo é a relação entre pesquisa de campo, escrita e teoria, e a posição da etnógrafa nessa relação \parencite[pp.~298, 300]{Kofes2015Maconaria}. Mobilizando Marilyn Strathern, Kofes situa a etnografia como prática reveladora ou \enquote{revelatória}, que convida o leitor a considerar o que não é revelado na evidência do que o é \parencite[p.~297]{Kofes2015Maconaria}. Com Jeanne Favret-Saada, sustenta que a etnógrafa precisa entrar no processo da palavra e se tornar uma falante como outra, posicionar-se no e como discurso, e ao escrever os resultados retornar à situação de enunciação, fazendo desse vai e vem entre a tomada inicial e sua retomada teórica o objeto mesmo da reflexão \parencite[p.~299]{Kofes2015Maconaria}.\footnote{A formulação que Suely Kofes recupera de Jeanne Favret-Saada vem de \textit{Les mots, la mort, les sorts: la sorcellerie dans le Bocage} \parencite[p.~32]{FavretSaada1977}, etnografia da feitiçaria no Bocage francês na qual a autora reconstrói a posição da etnógrafa como interna ao processo de enunciação que descreve. O argumento de Favret-Saada se desdobra em ensaio posterior, \enquote{Être affecté}, publicado originalmente na revista \textit{Gradhiva} e mais tarde recolhido em \textit{Désorceler} \parencite{FavretSaada1990SerAfetado}, no qual a autora formula a categoria do \enquote{ser afetado} como modalidade epistemológica específica da etnografia, distinta da observação participante e da empatia identificatória. A tradução brasileira desse ensaio, feita por Paula Siqueira Lopes e Tânia Stolze Lima, está disponível em \textit{Cadernos de Campo} \parencite{FavretSaada2005SerAfetado}, edição que torna acessível ao leitor brasileiro a formulação metodológica que Favret-Saada sintetiza retrospectivamente sobre a pesquisa do Bocage.} Essa posição encontra eco no que Strathern descreve como recriação imaginativa dos efeitos do próprio campo, gesto que torna a escrita antropológica parte do nexo entre pesquisa e teoria \parencite[p.~299]{Kofes2015Maconaria}. É a partir dessa sequência argumentativa que Kofes recupera Victor Turner: a tensão entre método e experiência se torna desnecessária quando dissolvida pela dramatização narrativa \parencite[p.~301]{Kofes2015Maconaria}. A experiência, ao se expressar narrativamente, conecta eventos e afecções e lhes confere significações e valores; essa expressão narrativa da experiência é o que Turner chama de \textit{structured unit of experience}, unidade estrutural da experiência. A narrativa, ao ligar o que se concebe como cesura, cria uma estrutura de relação, a estrutura de experiência. É por isso que a narrativa da etnógrafa é parte do método: o método se ensaia conforme se experimenta, e só pode ser narrado retrospectivamente porque se constitui junto com o objeto, no encontro de campo.\footnote{A formulação de Suely Kofes ressoa com o conceito de \textit{method assemblage} proposto por John Law em \textit{After Method: Mess in Social Science Research} \parencite{Law2004After}. Ambas as formulações recusam a normatividade prévia dos métodos convencionais, embora por vias distintas: Law aponta para a multiplicidade das realidades que o método produz e para a inseparabilidade entre o método e os mundos que ele torna pensáveis; Kofes sustenta a recusa pela via da experiência, argumentando que o método se ensaia através da experiência, inclusive da experiência de campo, e que a escrita antropológica é parte constitutiva desse ensaio \parencite[pp.~298--299]{Kofes2015Maconaria}.}

Essa concepção de etnografia como pesquisa antropológica que articula método e experiência dialoga com a teoria ator-rede de maneiras que interessam a esta tese. Primeiro, porque recusa categorias dicotômicas supostamente universais, como a separação entre humano e não humano. Segundo, porque trata o campo como processo de aprendizado: assim como Latour propõe que sigamos os atores sem decidir de antemão o que é técnico e o que é social, Kofes propõe que o método se ensaia conforme nos deixamos afetar pelo encontro com nossos interlocutores. Terceiro, porque ambas as abordagens organizam-se em torno das inscrições, dos documentos, das narrativas, dos artefatos e dos traçados que constituem simultaneamente nossos dados e nossos modos de produzir conhecimento.

A etnobiografia, desdobramento metodológico que Kofes formula em \textit{Uma trajetória, em narrativas} \parencite{kofes2001trajetoria}, também pode ser uma forma de investigar instituições, como faço nesta tese. Kofes adota o conceito de pessoa, em vez do de indivíduo, a partir de Marilyn Strathern, para quem a pessoa abriga o potencial para relações e é, ela própria, constituída por uma matriz de relações sociais, de sociabilidade, conforme exposto em \textit{Property, Substance and Effect} \parencite{strathern2000pse}.\footnote{A concepção de pessoa que Suely Kofes mobiliza em diálogo com Marilyn Strathern parte de uma formulação relacional da pessoa, desenvolvida pela autora a partir de sua etnografia na Melanésia. Em \textit{The Gender of the Gift: Problems with Women and Problems with Society in Melanesia} \parencite{strathern1988gender}, Strathern formula a pessoa como microcosmo de relações sociais, constituída pelas ações de outros, alimentação, cuidados, trocas, e capaz de destacar partes de si para criar novas relações. O argumento aprofunda-se em \textit{Property, Substance and Effect: Anthropological Essays on Persons and Things} \parencite{strathern2000pse} e em \textit{Partial Connections} \parencite{strathern2004partial}. Strathern designa essa pessoa como \textit{dividual}, entidade compósita e partível: na economia da dádiva melanésia, as pessoas trocam perspectivas, e o que circula nas trocas são porções de pessoa incorporadas em substâncias e objetos. A formulação contrasta com a noção euro-americana de indivíduo como unidade atômica e indivisível, que se completa pela posse de si mesmo. Strathern propõe pensar a pessoa como circuito de conexões parciais, em que cada parte estende a outra sem encerrar-se em uma identidade unitária final. Quando Kofes recorre a essa formulação para sustentar sua etnobiografia, transporta para o registro etnográfico a superação da pessoa como entidade preexistente à sociedade onde vive: a pessoa, e por extensão a etnógrafa e o coletivo etnografado, é o efeito momentâneo de uma matriz de relações que a constitui. A tradução brasileira dos ensaios de Strathern sobre o efeito etnográfico está disponível em coletânea \parencite{strathern2014efeito}, edição que torna acessível ao leitor brasileiro o conjunto de elaborações sobre o que Strathern chama de efeito etnográfico.} Se, como Kofes, é possível fazer a etnobiografia de uma Loja maçônica, então é igualmente possível fazer a etnobiografia de um centro de pesquisa em inteligência artificial. Isso significa reconstruir a trajetória e a vida de um coletivo, abrangendo suas práticas e estruturas, suas fundações e refundações, suas alianças e rupturas, suas transformações ao longo do tempo. Assim, o C4AI é também um actante cuja biografia se entrelaça com a minha trajetória de pesquisadora, e nesse entrelaçamento ressoa a formulação de Roy Wagner que Kofes recupera no mesmo capítulo: todo entendimento de uma outra cultura é um experimento com a nossa própria \parencite[p.~296]{Kofes2015Maconaria}. Esse gesto abre a possibilidade de pensar como cientistas e engenheiros fazem inteligência artificial a partir das mediações que a etnografia e a teoria ator-rede tornam visíveis nesta pesquisa.

Essa formulação ressoa com a atenção de Latour aos porta-vozes e às cadeias de tradução. Assim como o cientista no laboratório fala em nome de entidades que não podem falar por si mesmas, a antropóloga fala em nome de interlocutores cujas vozes são traduzidas, editadas, recontextualizadas no texto final. A responsabilidade dessa tradução é constitutiva do trabalho antropológico tanto quanto do trabalho científico. Concretamente, baseio-me em técnicas qualitativas diversas: observação de eventos acadêmicos e corporativos, entrevistas semiestruturadas com líderes de grupos de pesquisa, análise documental de acordos institucionais e acompanhamento de atividades de divulgação científica, de forma presencial ou remota. Essa diversidade de fontes caracteriza o que Kofes chama de pesquisa artesanal: trabalho que articula observação rigorosa e registro documental com a sensibilidade da experiência antropológica. Os dados desta tese provêm de conjunto de observações, escuta, conversas, anotações, entrevistas e convivência com pessoas ligadas ao Centro ao longo de quase quatro anos (março de 2022 a dezembro de 2025). Embora a parceria C4AI-IBM tenha sido inaugurada em outubro de 2020, antes do início do trabalho de campo, a articulação entre observação direta a partir de 2022, entrevistas retrospectivas e análise de relatórios institucionais cobrindo o período 2020--2025 permitiu reconstituir o ciclo completo da parceria, da fundação à dissolução. As entrevistas formais com Fábio Cozman (setembro de 2025) e Cláudio Pinhanez (outubro de 2025) foram realizadas após a comunicação interna da IBM sobre não-renovação da parceria (agosto de 2025), mas antes do anúncio público. Na época das entrevistas, nenhum dos interlocutores mencionou o fim iminente da parceria. Somente em dezembro de 2025, quando a IBM comunicou publicamente o encerramento, soube que ambos já conheciam essa decisão no momento das entrevistas. Essa descoberta retrospectiva conferiu dimensão temporal singular às narrativas: o que havia registrado como relatos sobre um projeto vivo revelou-se, em retrospecto, como testemunhos de quem já sabia de sua morte iminente mas não podia ou não quis dizê-lo.
 
A imersão prolongada exigiu também aprendizado de linguagens e códigos específicos. Assim como Kofes precisou decifrar a grafia e as abreviaturas maçônicas, precisei aprender a linguagem dos pesquisadores em inteligência artificial, abrangendo o vocabulário técnico, os modos de argumentar, as hierarquias implícitas entre abordagens e os critérios tácitos de avaliação. Esse processo de alfabetização foi condição de possibilidade para o aprofundamento da inteligência artificial. Mas a ausência inicial de conhecimento especializado teve um impacto duplo. Por um lado, constituiu uma vantagem metodológica, ao tratar a atividade científica como elemento novo, evitei a tendência de adotar antecipadamente a distinção entre social e técnico. Por outro lado, trouxe vulnerabilidade quando o foco mudou para a necessidade de aprender com os cientistas. Foi assim que compreendi, por exemplo, que o laboratório do cientista da computação é o computador, ideia aparentemente óbvia mas que possui implicações que merecem atenção. O computador, portátil e móvel, é mais do que simples objeto ou mediador: é o próprio laboratório, um actante, ponto na rede e simultaneamente também uma rede. Mas essa mobilidade tem limites que só se tornam visíveis quando a infraestrutura falha ou exige intervenção. Latour observa que objetos técnicos, enquanto funcionam, são intermediários silenciosos e mudos; é a crise, a quebra, o mau funcionamento, que revela sua existência, suas partes e sua agência \parencite[cap.~6]{Latour2017Esperanca}.\footnote{Em 4 de março de 2026, Luciano Antonio Digiampietri, pesquisador do Grupo de Pesquisa em Inteligência Artificial da EACH-USP (GrIA), fundado em 2006 e que divide com o C4AI a infraestrutura computacional, publicou no Instagram uma fotografia da placa de entrada do Centro com a legenda: \enquote{Hoje foi dia de presencial no laboratório. (Tive que apertar um botão, no servidor. ☺)} (Digiampietri, 2026, comunicação pessoal). O servidor em questão executa experimentos de IA, hospeda modelos de linguagem e realiza \textit{fine-tuning} para tarefas como detecção de \textit{bots} e identificação de notícias falsas. Luciano acessa o servidor remotamente de sua sala na EACH, a cem metros de distância, da reitoria da USP, ou de sua casa, e seus orientandos também o acessam remotamente. Naquele dia, o servidor desligou e não religou automaticamente, e foi necessário que ele fosse até a máquina para pressionar o botão de reinicialização. Em seguida, voltou a acessá-lo remotamente, da sala a cem metros de distância. A postagem condensa o paradoxo: o laboratório distribuído, que dispensa presença contínua, convocou o corpo do pesquisador para um gesto mínimo e indispensável. Mostra ainda que o C4AI não construiu sua infraestrutura do zero, instalou-se sobre redes já constituídas, como o GrIA, que existia desde 2006. Luciano explicou o contexto à pesquisadora via mensagem privada no Instagram em março de 2026 e autorizou expressamente o uso da postagem e da explicação como dado etnográfico; o termo de consentimento assinado encontra-se em posse da autora.} O laboratório de IA às vezes exige que alguém se desloque até onde as máquinas estão quando elas desligam ou param de funcionar por algum motivo para \enquote{apertar um botão}.

 % =====================================================================
% Cap. 1 — Fechamento da seção: articulação Kofes-Latour em parágrafo único
%
% Substitui o conjunto de 7 parágrafos que ia de "Apoio-me em diferentes
% abordagens..." até "...essa decisão de corte é parte do trabalho da
% etnógrafa".
%
% Os 7 parágrafos foram condensados em UM parágrafo longo e denso. Foram
% dispensados: a apresentação ECT/ESCT, a menção nominal à sócio-lógica
% (operação absorvida na frase sobre "redes mais fortes ou mais fracas"),
% o anúncio dos Capítulos 3 e 4 e a referência à linhagem específica.
% Foram preservados: a articulação Kofes-Latour, o conceito de hinterland
% (com a nota expandida), os dispositivos de inscrição como modulações do
% mesmo problema, e a transição para "Cortando redes".
% =====================================================================

Tanto a teoria ator-rede de Bruno Latour e a etnografia heterodoxa de Suely Kofes funcionam juntos nesta tese porque enfrentam, a partir de temas e questões distintas, o mesmo problema. Latour o formula no registro epistemológico, onde a coerência sobre o que os atores dizem ou fazem deve ser realiza dentro do próprio movimento de associação em que as redes se fazem e se tornam mais fortes ou mais fracas, onde os fatos sempre estarão relacionados ao tamanho de uma rede e da quantidade de suas associações \parencite[p.~331]{Latour2011CienciaEmAcao}. Kofes formula a questão a partir da escrita etnográfica, pela qual, o método ensaia-se conforme a etnógrafa se deixa afetar pelo que encontra, e esse ensaio só pode ser narrado retrospectivamente. Em Latour, esse gesto se realiza como seguir os actantes e mapear as redes que eles compõem; em Kofes, como deixar-se afetar pelo encontro de campo e compor a narrativa em que a experiência se estrutura. A própria adoção da teoria ator-rede nesta tese emergiu e foi personalizada de acordo com o processo de escrita, e isso é, na formulação de Kofes, o método ensaiando-se através da experiência. O gesto tem precedente etnográfico em \textit{Vida de Laboratório}, onde \textcite{Latour1997VidaLaboratorio} mostraram que o \textit{hinterland} de instrumentos de inscrição, protocolos experimentais e negociações constitui a própria ciência \parencite[p.~38]{Law2004After}.\footnote{O conceito de \textit{hinterland} é desenvolvido por Law em \textit{After Method} \parencite[pp.~34--42]{Law2004After}. O termo é emprestado da geografia: designa o território interior que abastece um porto sem aparecer nas trocas que esse porto realiza. Em termos metodológicos, o \textit{hinterland} de uma descrição é o conjunto de práticas, pressupostos, infraestruturas e recursos que permanecem no fundo sem serem tematizados, mas que são condição de possibilidade do que aparece. O conceito é distinto, mas inseparável da noção de \textit{corte}, desenvolvida na seção~\ref{subsec:cortando-redes} a partir de \textcite{strathern2011cortando}: o corte é a operação ativa que delimita o que entra e o que sai da análise; o \textit{hinterland} é o efeito dessa operação, o território que continua existindo e operando além do corte, sustentando o que aparece sem ser tematizado. A diferença em relação à simples ausência é que o \textit{hinterland} é constitutivo: o que precisa continuar operando silenciosamente para que o objeto da análise possa existir. No contexto desta tese, o conceito aparece com recorrência nos capítulos 3 e 4 para designar a rede longa de contratos, políticas de fomento e genealogia institucional que deixa de ser objeto direto de análise quando a atenção se volta para as práticas situadas de Marcelo, mas que não desaparece, pois continua sendo condição de possibilidade do que aparece.} Aquilo que Law nomeia como problema, fazer pesquisa quando o próprio fazer modifica aquilo que se investiga, opera nesta tese através de todos os dispositivos de inscrição mobilizados: os diagramas em Mermaid que performam conexões entre actantes, a curadoria bibliográfica que produz linhagens teóricas, a narrativa etnográfica que constrói o C4AI que documenta, e a composição com Claude que torna essa agência impossível de naturalizar. Cada dispositivo de inscrição \parencite{Latour2011CienciaEmAcao} decide o que torna visível e o que invisibiliza, e por isso a rede que a descrição abre precisa ser interrompida em algum ponto, isto é, precisa ser cortada. 

%%%%%
\subsection{Cortando redes: sobre os limites da descrição etnográfica}
\label{subsec:cortando-redes}

\epigrafe{%
Ninguém vive em todos os lugares; todo mundo vive em algum lugar. Nada está conectado a tudo; tudo está conectado a alguma coisa.}%
{Donna Haraway, \textit{Ficar com o problema}, 2023 \parencite[p.~60]{Haraway2023Ficar}.}

A pesquisa com o C4AI colocou-me diante de um dilema sobre como descrever uma rede que aparenta se expandir sem fim? O Centro está ligado a cerca de 250 pesquisadores, cada qual com trajetórias próprias; à USP, com toda a sua história institucional; à FAPESP, imbricada em políticas científicas estaduais e nacionais; à IBM, uma corporação global com mais de um século de existência; a projetos e sistemas específicos como Spira e Corpus Carolina; a infraestruturas computacionais espalhadas pelo planeta; ao mercado global de IA; a colaborações internacionais; e ao contexto da pandemia de Covid-19. Ao seguir os atores, eu era constantemente levada a novos lugares (e a novas questões) e, em algum ponto, eu tinha que parar de segui-los e de fazer associações, isto é, \textit{cortar} a rede.

\textcite{latour1996clarifications} explicita as propriedades topológicas que tornam o conceito de rede simultaneamente potente e problemático para a descrição etnográfica. Em \textit{On actor-network theory: a few clarifications}, ele caracteriza a rede através de três deslocamentos em relação às metáforas espaciais convencionais. O primeiro abandona a tirania da distância: elementos próximos podem estar infinitamente distantes se as conexões são analisadas, e elementos aparentemente remotos podem estar próximos quando seus elos são trazidos à descrição \parencite[p.~4]{latour1996clarifications}. O segundo supera a oposição micro-macro: ``uma rede nunca é maior que outra, ela é simplesmente mais longa ou mais intensamente conectada'' \parencite[p.~5]{latour1996clarifications}. O terceiro recusa a distinção interno-externo: ``uma rede é toda fronteira, sem interior nem exterior. A única questão que se pode colocar é se há ou não conexão estabelecida entre dois elementos'' \parencite[p.~6]{latour1996clarifications}. Esses três movimentos liberam a análise das categorias herdadas de níveis, esferas, territórios e estruturas. A consequência metodológica dessa formulação é que a rede não traz consigo limite ontológico interno. ``Literalmente não há nada além de redes, não há nada entre elas'' \parencite[p.~4]{latour1996clarifications}: cada conexão abre para outras conexões, cada actante conecta-se a outros actantes, e o trabalho da etnógrafa consiste em seguir esses elos. ``Cada laço, por mais forte que seja, é ele próprio tecido de fios mais fracos'' \parencite[p.~3]{latour1996clarifications}. O conceito que estipula como descrever as associações não estipula onde a descrição deve parar. O dilema enfrentado em campo era a forma empírica desse impasse conceitual: por onde terminar de seguir o C4AI?

\textcite{strathern2011cortando} formula esse dilema ao mostrar que redes não se tornam objetos de conhecimento sem antes serem cortadas. Em \textit{Cortando a rede}, ela retoma o conceito latouriano para identificar uma de suas propriedades, a autoilimitação: o conceito de rede funciona como metáfora das ``ilimitadas extensões e entrelaçamentos dos fenômenos'' \parencite[p.~7]{strathern2011cortando}. Frente a essa propriedade, ``a análise, como a interpretação, precisa ter um ponto; precisa ser encenada como um ponto de chegada'' \parencite[p.~7]{strathern2011cortando}. Uma rede é tão longa quanto seus elementos possam ser enumerados, e essa enumeração ``pressupõe uma totalização; ou seja, uma enumeração residindo em um objeto identificável (a soma). Ao parar, a rede pode ser `cortada' em um ponto'' \parencite[p.~7]{strathern2011cortando}. Strathern articula rede e híbrido como duas figuras da mesma operação: ``se tomarmos certos tipos de redes como híbridos socialmente expandidos então podemos tomar híbridos como redes condensadas. O trabalho de condensação trabalha como uma totalização ou parada'' \parencite[pp.~7--8]{strathern2011cortando}. Cada corte produz simultaneamente uma rede estendida, com elos enumeráveis e seguíveis, e um híbrido condensado, que retém os elementos heterogêneos da rede como objeto manipulável de conhecimento ou de transação. O corte não é gesto externo imposto a uma rede que estaria já dada, pois, as próprias práticas observadas realizam cortes, condensam relações em objetos, totalizam redes em pontos de chegada. O dote nupcial entre os \textit{'are'are} das Ilhas Salomão converte uma rede de prestações em objeto consumível na sequência mortuária; o sistema de patentes converte uma rede de quarenta colaboradores científicos em uma única invenção atribuível a seis inventores \parencite[pp.~8--9]{strathern2011cortando}.\footnote{Strathern desenvolve essa argumentação em diálogo etnográfico com sociedades cujos sistemas de troca tornam visíveis operações de corte que os euroamericanos tendem a naturalizar. Entre os \textit{'are'are} das Ilhas Salomão, a pessoa viva é composta de corpo, sopro e imagem, e na sequência funerária esses componentes são convertidos, respectivamente, em taro, porcos e dinheiro de conchas, este último figurando como a própria forma da imagem ancestral \parencite[pp.~10--12]{strathern2011cortando}. Os \textit{'are'are} designam essa sequência finalizadora como ``parada'' ou ``freio'': o dinheiro de conchas funciona simultaneamente como circulador e como interruptor de fluxos, condensando trocas anteriores e bloqueando reivindicações futuras.} O que Strathern quer dizer é que a análise reconhece e transporta cortes que as práticas já operam sem a intervenção do antropólogo. No argumento de Strathern, o conceito de \textit{propriedade}, por exemplo, funciona como mecanismo euroamericano por excelência de cortar redes. Ela é ``poderosa por causa de seus efeitos duplos, sendo simultaneamente uma questão de pertencer e de propriedade'' \parencite[p.~16]{strathern2011cortando}: separa quem participa de quem não participa, e converte o trabalho cooperativo de uma rede longa em objeto sobre o qual se exerce direito. ``Quando a tecnologia pode aumentar as redes, a condição de propriedade pode garantidamente cortá-las em um dado tamanho'' \parencite[pp.~16--17]{strathern2011cortando}. 

Pensar o corte da rede a partir do conceito de propriedade permitiu-me identificar, no caso desta tese, onde os cortes já operavam na prática antes que eu mesma tivesse de produzi-los analiticamente. O exemplo mais dramático disso foi a rescisão da parceria entre C4AI e IBM em 2025, examinada no capítulo~\ref{capitulo3}. O que descrevo com Latour como desestabilização ou reconfiguração da rede foi, com Strathern, um corte que mostrou o poder do direito e das decisões contratuais, como sobre a propriedade intelectual, patentes, código aberto e assim por diante. 

\textcite{Haraway2023Ficar} complementa o argumento de Strathern com a figura do \enquote{pensamento tentacular} que dá epígrafe a esta seção: nada se conecta a tudo, tudo se conecta a algo. Cada elo, cada nó da rede amarra e se liga a uma conexão particular. Por exemplo, esta tese começou pelo C4AI que conecta-se à IBM, que se conecta a Herman Hollerith, que se conecta às políticas de recenseamento americanas de 1890, que se conectam a práticas de governança colonial: cada passo é uma conexão a algo, não a tudo. A rede se ramifica porque cada actante conecta-se a outros actantes específicos, não porque seja ontologicamente ilimitada. Em nota a essa passagem, Haraway, ao se referir à obra \textit{Flight Ways}, de Thom van Dooren, oferece uma indicação do que está em jogo, algo que também nos serve para refletir sobre a tecnologia em geral e, em particular, sobre a inteligência artificial, que são muito mais materiais e imbricadas no mundo físico e em seus habitantes do que costuma parecer \parencite{Crawford2021Atlas}:

\begin{citacaoabnt}
A corrente de filosofia ecológica holística que insiste que ``tudo está conectado a tudo'' não nos ajudará aqui. Antes, diríamos que tudo está conectado a alguma coisa, que por sua vez está conectada a outra coisa. Ainda que, em última análise, tudo possa estar conectado entre si, a especificidade e a proximidade das conexões são importantes -- com quem estamos vinculados e de que maneiras. A vida e a morte acontecem nessas relações. Por isso, precisamos entender como se emaranham certas comunidades humanas, em particular, assim como outras comunidades de seres vivos, e de que maneiras esses emaranhados estão implicados na produção de sua extinção mútua e de seus respectivos padrões de morte amplificada \parencite[p.~60, n.~4]{Haraway2023Ficar}.
\end{citacaoabnt}

\noindent Onde os actantes param, a análise também pode parar, sem trair a rede, apenas escolhendo qual caminho seguir. Assim, o critério do corte acompanha a especificidade e a intensidade das conexões que de fato fazem a rede ser aquilo que é.

John Law, por outro lado, sistematiza esse corte de outra forma para descrever o que cada corte produz e o que ele deixa para fora. Em \textit{After Method}, ele caracteriza o \textit{hinterland} (o contexto ou a retaguarda) de qualquer \textit{method assemblage} (montagem metodológica).\footnote{O \textit{hinterland} e o \textit{corte} são conceitos complementares mas distintos. O corte é a operação ativa que delimita o que entra e o que sai da análise; o \textit{hinterland} é o efeito dessa operação, o território que continua operando silenciosamente além do corte, como condição de possibilidade do que aparece. O conceito de \textit{corte} foi desenvolvido acima a partir de \textcite{strathern2011cortando}; o de \textit{hinterland}, introduzido aqui, será retomado nos Capítulos~\ref{capitulo3} e~\ref{capitulo4}.} Law quer dizer que, por exemplo, indo além de bancadas de laboratório, o \textit{hinterland} estende-se para conhecimento tácito, software de computador, habilidades linguísticas, capacidades de gestão, sistemas de transporte e comunicação, escalas salariais, fluxos de financiamento, prioridades de agências de fomento e agendas abertamente políticas e econômicas. A lista, dessa forma, tornaria-se interminável, e as suas fronteiras tão porosas e extensíveis para fora e em todas as direções \parencite[pp.~40--41]{Law2004After}. Desse modo, pensando com Law, não é possível fazer uma descrição completa da rede, apenas cortes estratégicos que tornam certos aspectos presentes enquanto deixam outros ausentes. Law sistematiza isso através do conceito de \textit{method assemblage} como ``a produção ou \textit{crafting} de um feixe de relações ramificadas que gera presença, ausência manifesta e \textit{otherness}'' \parencite[p.~84]{Law2004After}. Asssim, toda pesquisa articula essas três dimensões: a presença é o que é trazido aqui e agora, e que tornamos visível, condensado em representações; a ausência manifesta é o que está ausente mas reconhecidamente necessário ao presente; \textit{otherness} é o que é necessário à presença mas não pode ou não é tornado manifesto.\footnote{A \textit{otherness} distingue-se do \textit{hinterland}, conceito que Law também desenvolve em \textit{After Method} \parencite[pp.~34--42]{Law2004After} e que opera com frequência nos Capítulos 3 e 4 desta tese. O \textit{hinterland} é o território que continua operando silenciosamente além do corte, reconhecível, rastreável, condição de possibilidade do que aparece, mas não tematizado. A \textit{otherness}, por sua vez, é o que o próprio \textit{method assemblage} torna impensável ou invisível, ou aquilo o que o quadro analítico não consegue sequer nomear sem se transformar. Todo corte produz as três dimensões simultaneamente, portanto, o que aparece (presença), o que ficou fora mas poderia ter sido seguido (ausência manifesta), o que sustenta silenciosamente o que aparece (\textit{hinterland}), e o que decide-se invisibilizar (\textit{otherness}).}

Em resumo, tudo isso faz parte do \textit{crafted} \parencite{Law2004After}, dos modos pelos quais se faz uma pesquisa, entendida, sobretudo, como escolhas que o pesquisador faz consciente quanto aos efeitos dessas escolhas e responsável por suas consequências. No entanto, ao entender e assumir esse modo de fazer pesquisa, eu também percebi que precisava apresentar ao leitor o que isso significa nesta etnografia. Por exemplo, as notas de rodapé desta tese são, antes de tudo, expressão de um modo de pensar. À medida que a escrita avançava, a cada investida analítica eu acabava fazendo associações com materiais, autores e experiências diversos, alguns parte do meu repertório e outro não, e essas associações pediam registro. Registrá-las no corpo do texto parecia-me interromper o argumento, mas descartá-las apagaria as conexões que eu havia construído. Por isso, as notas funcionam aqui nesse impasse e para além de uma decisão estilística. No entanto, reconhecer o papel e a função dessas notas no meu texto não foi imediato. Foi só depois, ao observar que as notas se comportavam como pequenas extensões da rede descrita no texto principal, que o vocabulário de \textcite{strathern2011cortando} e de \textcite{Law2004After} foi relacionado com a minha prática de pesquisa. Cada nota acompanha uma associação que o argumento precisou deixar no \textit{hinterland} para manter seu fluxo narrativo, mas que a topologia analítica da teoria ator-rede me impedia de tratar como algo complementar ou suplementar à própria pesquisa. No capítulo~\ref{capitulo1}, as notas estendem a rede para a crítica de Le Guin ao uso da palavra ``tecnologia'', para o botão do servidor que convocou o pesquisador Digiampietri ao laboratório. No capítulo~\ref{capitulo2}, estendem-se para a programação dos congressos brasileiros onde a inteligência artificial mal aparece, para a bibliometria que documenta um campo emergente. No capítulo~\ref{capitulo3}, passa pela trajetória corporativa da IBM desde os cartões perfurados de Hermann Hollerith, o modelo de ecossistemas empresariais de Moore, a aquisição da Red Hat por US\$ 34 bilhões e as cadeias de translação que conectam o Censo de 1890 dos Estados Unidos ao C4AI no ano de 2020. No capítulo~\ref{capitulo4}, acompanham a história da CNN14 pré-treinada no AudioSet do Google, a aprovação ética do CEP-HCFMUSP em abril de 2020, a voz de pessoas hospitalizadas com Covid-19 que a escala Mel reconstrói computacionalmente, a retropropagação com gradiente descendente das ausências e presenças dos artigos do Spira. Cada uma dessas extensões poderia ocupar o texto principal de outra tese. A decisão de registrá-las em notas é ela mesma um corte no sentido de \textcite{strathern2011cortando}: a rede precisa ser interrompida em algum ponto para que a descrição possa circular. As notas são, também e dessa forma, os pontos em que a interrupção, os cortes (mas também as aberturas), podem ser constatados. Funcionam, nesse sentido, como demonstração material de que as redes se estendem para onde quer que o observador as siga, e que a decisão de parar de segui-las pertence, sobretudo, a quem as descreve. Os diagramas em Mermaid que acompanham os Capítulos~\ref{capitulo1}, \ref{capitulo3} e \ref{capitulo4} operam no mesmo registro, condensando em um plano sinóptico as cadeias de translação \parencite{Latour2011CienciaEmAcao} que o texto principal percorre em sequência. E o site \textit{Etnografia de um Centro de Inteligência Artificial}, disponível em \url{https://julianehelanski.github.io/tese-c4ai-spira/}, leva esse mesmo gesto a uma terceira superfície, ainda plana, mas com mais dimensões: cada nó é categorizado segundo um tipo ontológico (capítulo, pessoa, instituição, actante não humano, conceito, mídia, simpósio, dados) e o leitor compõe seu próprio percurso ao clicar nas associações, redistribuindo a agência interpretativa que a sequência narrativa do texto impresso concentra na pesquisadora. Notas, diagramas e site são, cada um a seu modo, demonstrações materiais de que as redes se estendem para onde quer que o observador as siga, e que a decisão de parar de segui-las (ou de abri-las para que o leitor as siga em outras direções) pertence, sobretudo, a quem as descreve.

De modo que, optei por fazer uma série de cortes estratégicos que definem o que está presente, o que fica na ausência manifesta e o que é relegado à \textit{otherness}, principalmente, nos capítulos 3 e 4 onde as redes são construídas a partir de actantes chaves. Sintetizei esses cortes nas tabelas \ref{tab:corte-extensao} e \ref{tab:corte-densidade}.

\vspace{0.5cm}

\begin{longtable}{p{4cm}p{10cm}}
\caption{Capítulo~\ref{capitulo3}: Corte pela extensão} \label{tab:corte-extensao} \\
\toprule
\textbf{Dimensão} & \textbf{Conteúdo} \\
\midrule
\endfirsthead
\multicolumn{2}{c}{\tablename\ \thetable\ -- Continuação} \\
\toprule
\textbf{Dimensão} & \textbf{Conteúdo} \\
\midrule
\endhead
\midrule
\multicolumn{2}{r}{\textit{Continua na próxima página}} \\
\endfoot
\bottomrule
\endlastfoot
\textbf{O que torno presente} & Práticas de construção da ciência na pesquisa em IA através de aspectos institucional e articulação de redes: apresentações públicas de Fábio e Cláudio (eventos de difusão), negociações de parceria, mobilização de financiamentos, posicionamento estratégico do Centro. Reconstitui a rede do C4AI em sua extensão histórica e institucional, a trajetória da IBM desde empresa de tabulação até corporação global de computação; políticas científicas brasileiras de fomento à pesquisa em IA; acordos entre universidade pública, agência estadual de fomento e corporação transnacional; reorientações de linhas de pesquisa acompanhando movimentos corporativos. \\
\addlinespace
\textbf{Ausência manifesta} & Práticas técnicas situadas dos outros grupos de pesquisa do C4AI: código escrito em oceanografia, agronegócio, línguas indígenas e saúde; modelos treinados simultaneamente em diferentes desafios; coletas de dados em campo e hospitais; negociações com comunidades indígenas para documentação linguística; decisões sobre arquiteturas de rede neural em outros projetos; impasses metodológicos enfrentados por outras equipes. Sei que essas práticas existem, acontecem e sustentam o Centro, mas não as acompanho em detalhe. \\
\addlinespace
\textbf{\textit{Otherness}} & Os outros 248 pesquisadores do Centro cujas trajetórias não acompanho; detalhes técnicos granulares de cada projeto; negociações internas entre grupos de pesquisa; disputas por autoria em publicações; a composição das equipes; a vida cotidiana dos laboratórios distribuídos. Esses elementos precisam desaparecer para que a rede longa se torne visível como objeto analisável. \\
\addlinespace
\textbf{Justificativa do corte} & A dissolução da parceria em 2025, ocorrida durante a finalização desta pesquisa, abriu uma caixa-preta institucional: o que parecia arranjo estabilizado revelou-se composição precária cuja fragilidade só se tornou visível com a ruptura. O registro do início e do fim dessa parceria permitiu acompanhar como as redes se constroem e como se desfazem. O corte pela extensão traz à tona heranças corporativas e como elas estruturam ecossistemas contemporâneos de inovação na pesquisa em IA, revelando padrões de assimetrias e vulnerabilidades que estudos realizados antes ou depois da dissolução não teriam condições de identificar. \\
\end{longtable}

\vspace{0.5cm}

\begin{longtable}{p{4cm}p{10cm}}
\caption{Capítulo~\ref{capitulo4}: Corte pela densidade} \label{tab:corte-densidade} \\
\toprule
\textbf{Dimensão} & \textbf{Conteúdo} \\
\midrule
\endfirsthead
\multicolumn{2}{c}{\tablename\ \thetable\ -- Continuação} \\
\toprule
\textbf{Dimensão} & \textbf{Conteúdo} \\
\midrule
\endhead
\midrule
\multicolumn{2}{r}{\textit{Continua na próxima página}} \\
\endfoot
\bottomrule
\endlastfoot
\textbf{O que torno presente} & Narrativa de Marcelo sobre suas práticas técnicas situadas: como escreveu código, treinou modelos, coletou dados (áudios de pacientes COVID, gravações de fala), negociou com médicos para acesso a dados clínicos, decidiu sobre features acústicas, enfrentou impasses com vieses durante treinamento, articulou abordagens computacionais distintas. Análise de artigos científicos publicados pelo projeto Spira e Corpus Carolina. Segue a rede do C4AI em sua densidade de práticas situadas através da retrospectiva de um pesquisador específico sobre a construção de dois projetos. \\
\addlinespace
\textbf{Ausência manifesta} & Observação direta das práticas técnicas de Marcelo: não acompanhei em tempo real o código sendo escrito, os modelos sendo treinados, as decisões metodológicas sendo tomadas. Práticas técnicas situadas dos outros grupos de pesquisa do C4AI. Práticas de construção institucional e articulação de redes extensas trabalhadas no Capítulo~\ref{capitulo3}. Tudo isso aparece como contexto pressuposto, reconhecidamente ausente da narrativa retrospectiva de um único pesquisador. \\
\addlinespace
\textbf{Justificativa do corte} & Marcelo narrou a pesquisa tal como ela se fez, com contingências, improvisações, fracassos e sucessos, diferente da versão estabilizada que os artigos científicos apresentam. Essa narrativa retrospectiva, ainda que mediada por entrevista e não por observação direta, permitiu acessar camadas do processo de pesquisa que a inscrição científica final apaga: discussões sobre quais dados utilizar, impasses metodológicos com vieses identificados tardiamente, negociações entre abordagens computacionais conflitantes. O corte pela densidade mostrou como a transformação de vozes humanas em dados computacionais depende de cadeias de tradução tão frágeis quanto inventivas, revelando mediações sociotécnicas concretas e situadas. \\
\end{longtable}

\vspace{0.5cm}

Strathern e Law oferecem o vocabulário para pensar como a etnógrafa interrompe a rede e descreve o que a interrupção produz. Mas o que dizer do objeto que aparece de cada corte? Se cada corte produz uma descrição distinta da rede, que estatuto têm essas descrições umas em relação às outras? São versões parciais de um mesmo C4AI, perspectivas que se complementariam em uma síntese final, ou objetos etnográficos que não se reduzem a um substrato comum? É para responder a essa pergunta que mobilizo a praxiografia de Annemarie Mol.

Deste modo, esta tese não apresenta perspectivas diferentes sobre um mesmo C4AI nem postula a existência de vários C4AIs, mas descreve múltiplos enactments do Centro no sentido praxiográfico que Mol desenvolve em \textit{The Body Multiple} \parencite{mol2002body}. Mol rompe com a epistemologia da representação ao argumentar que diferentes práticas não observam um objeto único de perspectivas distintas, mas performam objetos diferentes: o corpo na clínica não é o mesmo corpo no laboratório de patologia. Como ela resume: ``realidade não precede as práticas mundanas nas quais interagimos com ela, mas é moldada dentro dessas práticas'' \parencite[p.~6]{mol2002body}.

Esse deslocamento tem implicações políticas que Mol e Law formulam como ontological politics: se realidades são múltiplas e performadas por práticas, a questão deixa de ser qual representação está correta e passa a ser quais realidades queremos fortalecer ou enfraquecer \parencite[pp.~74--75]{mol1999ontological}; \parencite{Law2004After}.\footnote{Latour já abre essa discussão em \textit{Políticas da Natureza} ao formular a composição progressiva do mundo comum \parencite[p.~18]{latour2004politics}. Law e Mol radicalizam essa intuição com ontological politics. Enquanto Latour trabalha com a ideia de um mundo comum (singular, ainda que em construção), Law e Mol propõem multiplicidade ontológica.}

\textcite{holbraad2017ontological} apresentam a virada ontológica como tecnologia de descrição etnográfica: não é metafísica, é método. A virada ontológica não pergunta o que o mundo realmente é, pergunta o que o objeto de investigação tem que ser para que a prática etnográfica em torno dele faça sentido. A pergunta é reflexivo-metodológica. O ponto é manter aberta a possibilidade de que o objeto seja diferente do que se supôs inicialmente. Esse esclarecimento muda a leitura do que Mol e Law nomeiam como multiplicidade. Não se trata de afirmar que existam múltiplas realidades, mas que há maneiras diferentes de abordar o objeto via práticas que produzem versões dele não redutíveis a uma única descrição.

Esse movimento teórico tem consequências para como pensamos método e política. Na formulação de Law e Mol, descrever e intervir são modos de fazer realidades. Os cortes nas redes que realizo nesta tese não são rupturas arbitrárias impostas pela análise. Seguindo \textcite{strathern2011cortando}, as próprias práticas observadas já operam cortes, distinções e condensações que a análise reconhece e transporta.

Mas cortar também pode ser operação produtiva: abre para o desconhecido ou faz vazar possibilidades contidas. Assim como \textcite{costa2024manchando} fez vazar questões feministas em sua armadilha metodológica, faço aqui vazar questões políticas sobre parcerias público-privadas na pesquisa em IA. O corte real da rede (simbolizado pelo fim da parceria C4AI-IBM) em 2025 funcionou também como corte analítico; ao cortar nas práticas de Marcelo, faço vazar contingências e improvisações que a inscrição científica final apaga. Esta tese parte da TAR para seguir actantes e traçar redes, mas também recorre a Law, Mol e Strathern para pensar as multiplicidades ontológicas e os cortes como operações descritivas, analíticas e produtivas de realidades.

Pensando dessa maneira, minhas escolhas narrativas, os cortes nas redes, as práticas de descrição não representam o C4AI de forma mais ou menos adequada; elas o performam como objetos etnográficos específicos e, ao fazê-lo, fortalecem certas versões dessa realidade em detrimento de outras.\footnote{Performar é diferente de representar. Law explica: entidades alcançam sua forma como consequência das relações nas quais estão localizadas. Mas isso significa que elas também são performadas nas, pelas e através dessas relações \parencite[p.~4]{LawHassard1999ANT}. Se objetos são performados, então tudo é incerto e reversível, pelo menos em princípio, poderia ser de outra forma.}

Os Capítulos~\ref{capitulo3} e~\ref{capitulo4} demonstram isso empiricamente. O Capítulo~\ref{capitulo3} reconstrói uma rede extensa que atravessa 135 anos, conectando genealogias corporativas, políticas científicas e alianças institucionais. O Capítulo~\ref{capitulo4} segue uma rede densa e curta: as mediações sociotécnicas situadas através das quais Marcelo transforma vozes humanas em dados computacionais. As redes documentadas nesses capítulos têm topologias distintas e performam versões diferentes do Centro.

O C4AI como objeto etnográfico narrado com meus cortes subjetivos (Capítulo~\ref{capitulo1}) é uma dessas realidades; o C4AI como problema teórico que demanda rede específica de inscrições e aliados conceituais (Capítulo~\ref{capitulo2}) é outra; o C4AI como rede histórica institucional extensa (Capítulo~\ref{capitulo3}) é outra; o C4AI como práticas técnicas densamente situadas (Capítulo~\ref{capitulo4}) é outra ainda. Cada capítulo é um corte que não apenas descreve, mas intervém ao fortalecer certas relações e enfraquecer outras.

No Capítulo~\ref{capitulo1}, o C4AI emerge como objeto etnográfico performado pelas minhas escolhas narrativas e metodológicas. No Capítulo~\ref{capitulo2}, o Centro é performado como problema teórico que demanda arregimentação específica de aliados conceituais. No Capítulo~\ref{capitulo3}, é performado como rede de alianças institucionais que se estende de 1890 a 2025. No Capítulo~\ref{capitulo4}, emerge como conjunto de mediações sociotécnicas concretas através das quais vozes humanas se transformam em dados computacionais.

As práticas técnicas que Marcelo narra estão atravessadas pela rede institucional que reconstituo no Capítulo~\ref{capitulo3}: ele usa infraestrutura custeada pela parceria IBM-FAPESP, publica em conferências patrocinadas por corporações, mobiliza datasets cuja existência depende de políticas científicas. Mas Marcelo não narra suas práticas nesses termos; ele as descreve como decisões metodológicas situadas, contingentes, improvisadas. Suas práticas não se reduzem às determinações estruturais da rede longa. Inversamente, quando reconstituo a genealogia institucional no Capítulo~\ref{capitulo3}, pesquisadores como Marcelo estão presentes: suas publicações justificam investimentos, seus resultados alimentam relatórios à FAPESP. Mas a rede institucional não se reduz à soma dessas práticas individuais.

Juntos, os múltiplos enactments não esgotam a rede do C4AI. Seria impossível fazê-lo. Cada um fortalece versões específicas dessa realidade, intervindo ao performar certas conexões em detrimento de outras \parencite{Law2004After}. Esta tese intervém ao performar enactments específicos do C4AI, sabendo que outros seriam possíveis e que as escolhas carregam implicações políticas.
%%%%
\subsection{Triangulação metodológica: documental, temporal e por extensão}
\label{subsec:triangulacao_cap1}

O princípio de seguir os actantes, apresentado na seção anterior, coloca-nos diante do problema metodológico já formulado: as redes se desdobram infinitamente e precisam ser delimitadas para serem analisadas \parencite{strathern2011cortando}. Trata-se de cortar a rede e estabelecer fronteiras estratégicas entre presença, ausência manifesta e otherness \parencite{Law2004After}. A triangulação metodológica constitui estratégia para operacionalizar esse seguir com cortes: mapear associações e escolher onde cortar, o que tornar presente e o que deixar ausente.

As tabelas \ref{tab:corte-extensao} e \ref{tab:corte-densidade} apresentadas na seção anterior sintetizam os cortes analíticos manifestos na estrutura dos capítulos: o Capítulo~\ref{capitulo3} reconstrói uma rede extensa, enquanto o Capítulo~\ref{capitulo4} segue uma rede densa. Além desses cortes estruturais, opero também cortes complementares de natureza mais pragmática, que orientam o trabalho de campo e a coleta de dados transversalmente: triangulação documental, temporal e por extensão. Esses cortes pragmáticos não substituem os cortes analíticos dos capítulos, mas os atravessam e alimentam, estabelecendo critérios práticos para decidir quais atores seguir, quais documentos analisar, quais genealogias\footnote{Utilizo o termo ``genealogia'' em seu sentido comum: o rastreamento de origens, evoluções e linhagens ao longo do tempo. A genealogia corporativa e tecnológica que desenvolvo no Capítulo~\ref{capitulo3} rastreia as origens e transformações da IBM desde Herman Hollerith (1890) até os arranjos contemporâneos de IA, identificando heranças, continuidades e rupturas.} rastrear e como articular redes de diferentes alcances.

a) \textbf{Triangulação documental}: contrasto inscrições formais e práticas etnograficamente observadas em sua dimensão temporal. Este corte segue o princípio de Latour de que documentos também são actantes que fazem coisas no mundo, são inscrições que performam associações. O acordo FAPESP-IBM de 2016, os relatórios anuais do C4AI (2021--2025), apresentações públicas e comunicações internas foram triangulados com observações de campo e entrevistas. Essa triangulação possui dupla dimensão: sincrônica (convergências e tensões entre documentos e práticas num mesmo momento) e diacrônica (transformações visíveis na série temporal dos relatórios anuais, deslocamentos entre o acordo inicial e seus desdobramentos). Documentos de diferentes períodos revelam não apenas o que foi inscrito, mas como as inscrições se transformaram: ênfases que se deslocam, projetos que emergem ou desaparecem, narrativas que se estabilizam ou se fragmentam. Este corte torna presente documentos formais em sua espessura temporal enquanto mantém na ausência manifesta práticas informais que sustentam ou contradizem o que está inscrito.

b) \textbf{Triangulação temporal}: acompanho as translações\footnote{Para Bruno Latour, a translação (ou tradução) é a interpretação e o deslocamento que os construtores de fatos operam sobre seus próprios interesses e os de outros atores que desejam alistar. O termo possui sentido duplo: linguístico (traduzir objetivos de uma linguagem para outra) e geométrico (deslocar atores de seus caminhos habituais para novas direções). Desenvolvo esse conceito no Capítulo~\ref{capitulo3} ao interpretar as cadeias de tradução do sistema de tabulação de Herman Hollerith \parencite[pp. 159-225]{Latour2011CienciaEmAcao}.} entre diferentes momentos e cadeias sociotécnicas. Este corte estende a rede temporalmente para identificar como actantes contemporâneos carregam heranças de associações passadas, seguindo os próprios atores quando mobilizam ou evocam o passado. A afirmação de Cláudio sobre as estratégias corporativas da IBM foi triangulada com a literatura histórica sobre a trajetória da corporação, seus padrões em diferentes configurações sociotécnicas, dos cartões perfurados aos sistemas open source, e destes aos ecossistemas de IA. O corte opera seleção sobre como o passado é mobilizado: ao seguir a genealogia corporativa da IBM, evoco aspectos da história da computação, das políticas científicas e tecnológicas brasileiras, mas o faço através da lente dessa trajetória corporativa e de sua relação com a FAPESP, o C4AI e a USP.

c) \textbf{Triangulação por extensão}: articulação entre redes de extensões diferentes. Este corte segue os actantes através de suas conexões sem presumir que uma rede seja mais real ou determinante que outra. Como Latour argumenta, a noção de rede dissolve a distinção micro/macro: ``uma rede nunca é maior que outra, é simplesmente mais longa ou mais intensamente conectada'' \parencite[p. 370]{Latour_SozialeWelt_1996}.\footnote{A metáfora de escalas (do indivíduo ao Estado-nação, passando por família, grupos e instituições) é substituída por uma metáfora de conexões. A noção de rede não faz suposição alguma sobre se um lócus específico é macro ou micro, e não modifica as ferramentas para estudar um ou outro \parencite[pp. 369-371]{Latour_SozialeWelt_1996}.} Marilyn Strathern oferece vocabulário complementar através do conceito de conexão parcial: as redes documentadas nos Capítulos~\ref{capitulo3} e~\ref{capitulo4} não constituem níveis de uma mesma realidade, mas extensões parciais umas das outras, conectadas sem serem isomórficas.\footnote{Para Strathern, a complexidade não diminui nem aumenta com a escala; ela apenas se replica. A conexão parcial permite organizar o conhecimento de forma não linear, reconhecendo que verdades etnográficas são simultaneamente incompletas e comprometidas com contextos específicos \parencite{strathern2005partial}.} Este corte por extensão torna presente duas redes específicas (a extensa do Capítulo~\ref{capitulo3}, a densa do Capítulo~\ref{capitulo4}) enquanto deixa ausentes outras redes possíveis. As redes conectam-se pela possibilidade de que uma estenda parcialmente a outra.

A reflexividade posicional atravessa todas essas dimensões de corte. Não como confissão autobiográfica, mas como reconhecimento de que a etnógrafa também é actante na rede que descreve. Ao longo da pesquisa, tornei-me colaboradora do grupo de humanidades do C4AI e membro do Observatório Brasileiro de Inteligência Artificial (OBIA), o que exigiu gestão ativa da proximidade com o campo. Essa posição ambivalente produziu tensões que foram incorporadas reflexivamente à análise. O acesso privilegiado a eventos, conversas informais e dinâmicas internas coexistiu com o desconforto de documentar vulnerabilidades de uma instituição da qual passei a fazer parte. Como o etnógrafo que Latour descreve ``rondando por ali, inútil, de braços cruzados'' \parencite[p. 62]{Latour2017Esperanca}, experimentei o desajuste de não compreender inicialmente aquele mundo e transformei esse desconforto em recurso metodológico. Minha presença no campo é, ela mesma, um dos cortes que torna presente certas relações enquanto deixa outras ausentes.

Ao longo do trabalho de campo, realizei entrevistas formais com os coordenadores dos cinco grupos de pesquisa do C4AI (Processamento de Linguagem Natural, Computação Bio-Inspirada e Aprendizado de Máquina, Saúde, Agronegócio e Biologia, AI Humanities). Os cortes analíticos que estruturam esta tese, porém, demandaram seleção estratégica: as entrevistas com Fábio (diretor do C4AI), Cláudio (vice-diretor executivo pela IBM) e Marcelo (coordenador do grupo de NLP e responsável pelos projetos Spira e Corpus Carolina) foram selecionadas porque suas narrativas permitiam performar os dois enactments complementares do Centro documentados nos Capítulos~\ref{capitulo3} e~\ref{capitulo4}. As demais entrevistas, embora tenham informado a compreensão geral do campo, permanecem na ausência manifesta: reconhecidamente necessárias, mas não tornadas presentes na análise final.

O momento das entrevistas formais com os atores-chave conferiu aos dados dimensão temporal singular e exemplifica como cortes temporais produzem diferentes objetos etnográficos. A comparação com Aramis, esboçada nas páginas anteriores, ganha precisão quando se examina a posição temporal das duas pesquisas. Latour iniciou seu estudo etnográfico do Aramis em março de 1988, poucos meses após o encerramento do projeto em dezembro de 1987, e realizou cerca de cinquenta entrevistas ao longo de um ano \parencite{Latour1993Aramis}. O resultado foi a multiplicação de interpretações: não uma explicação, mas pelo menos vinte versões divergentes sobre o que havia acontecido e por quê. Para Latour, essa dispersão de narrativas é constitutiva de projetos que deixam de existir: quando não há objeto estabilizado que unifique os pontos de vista, ``não posso encontrar a soma das interpretações do Aramis, pois não há interseção comum'' \parencite[p. 15]{Latour1993Aramis}.

Esta pesquisa ocupou posição temporal distinta: as entrevistas foram realizadas não após, mas durante o processo de dissolução, no intervalo crítico entre a comunicação interna da IBM sobre não-renovação (agosto de 2025) e o anúncio público do encerramento (novembro-dezembro de 2025). Se Latour chegou quando o Aramis já era ``um quase-objeto que circula como token em cada vez menos mãos'' \parencite[p. 15]{Latour1993Aramis}, eu estava presente enquanto as mãos ainda seguravam o token, hesitantes sobre se deveriam soltá-lo. Essa diferença temporal não eliminou a dispersão de interpretações; talvez a tenha intensificado. Entrevistas formais divergiam de pronunciamentos públicos em eventos como a Bluetalks; conversas informais contradiziam posicionamentos oficiais; certas questões permaneciam tão vagas que se tornava impossível chegar a uma explicação única. Por que a IBM encerrou a parceria? Como a corporação extraiu valor do convênio? De que modo influenciou a agenda de pesquisa? O que cada parte efetivamente obteve? Ninguém disse abertamente ou, quando disse, as versões não coincidiam.

Essa opacidade não constitui falha metodológica a ser superada; ela é achado empírico que revela a natureza mesma de parcerias em dissolução. Enquanto Latour documentou a dispersão das interpretações após o encerramento do projeto, pude observar essa multiplicação em tempo real, no momento em que actantes ainda negociavam o significado do que estava acontecendo. Os porta-vozes, no momento da crise, articulam retrospectivamente trajetórias que, em tempos de estabilidade, pareceriam óbvias demais para merecer explicação. Diferentemente de uma investigação posterior, porém, ainda não haviam consolidado narrativas definitivas sobre o que deu errado. Suas hesitações, silêncios e contradições são mais reveladoras do que explicações polidas construídas anos depois. O corte temporal que estabeleci tornou visível essa multiplicidade interpretativa em estado nascente.

\section{``Existências parciais'': quase-máquinas, híbridos e a composição com IA}
\label{subsec:existencias-parciais}

\epigrafe{Direct access to anything is not in the power of \enquote{humans} souls and machinery. Modems and faxes add just a few more artifacts in the midst of all of those that already compose \enquote{natural} interactions. [...] We all have only partial existence.}
{Bruno Latour e Richard Powers, \textit{Two Writers Face One Turing Test: A Dialogue in Honor of HAL}, 1998.}

As polarizações e multiplicidades descritas nas páginas anteriores exigem ferramental teórico capaz de compreender como pessoas e sistemas de IA se compõem sem reduzir essa composição a categorias prévias (usuário/ferramenta, humano/máquina, autêntico/artificial). Os conceitos de existência parcial, conexão parcial e agência distribuída oferecem esse ferramental.

\subsection*{Do princípio de simetria às existências parciais}

Proponho aqui uma leitura da formulação ``existências parciais'' não como ajuste vocabular, mas como deslocamento teórico interno à própria obra de Latour. O Latour de \textit{Ciência em Ação} \parencite*{Latour2011CienciaEmAcao} e de \textit{Jamais Fomos Modernos} \parencite*{Latour1994Jamais} opera com o princípio de simetria generalizada: tratar humanos e não humanos nos mesmos termos. Esse é passo decisivo porque desmonta a pressuposição de que apenas um dos lados da relação merece descrição. Mas o vocabulário da simetria ainda pressupõe dois polos a serem simetrizados, e a operação contesta a hierarquia sem dissolver a própria partição.

O diagnóstico dos híbridos é persistente na obra de Latour: dos quase-objetos de \textit{Jamais Fomos Modernos} ao projeto Aramis, dos híbridos de \textit{A Esperança de Pandora} \parencite*{Latour2017Esperanca} à crise da crítica em \textit{On the Modern Cult of the Factish Gods} \parencite*{latour2010factish}. As misturas não são categoria intermediária a ser superada, mas a condição mesma do mundo que habitamos. No diálogo com Powers de 1998, a formulação se radicaliza: ``todos temos apenas existência parcial'' \parencite{Latour2008Powers}. Não são mais dois polos simetrizados, nem entidades definidas pelo que não são; são existências que nunca foram completas para começar. E na \textit{Investigação sobre os Modos de Existência} \parencite*{latour2019investigacao}, o próprio conceito de ator-rede é reconhecido como insuficiente, apenas um modo entre outros.

Esse movimento que identifico na obra de Latour não surgiu de revisão bibliográfica, mas como consequência de meu próprio percurso em campo. Foi acompanhando pesquisadores de IA no C4AI, observando como seus algoritmos, datasets e modelos eram simultaneamente produtos de escolhas humanas e entidades com comportamentos que surpreendiam seus próprios criadores, que a insuficiência do vocabulário simétrico surgiu. A inteligência artificial generativa exemplifica de forma contundente o que significa existência parcial: é composta por milhões de decisões humanas sedimentadas e, no entanto, produz resultados que nenhum de seus criadores previu ou controla inteiramente; opera sobre conteúdo simbólico e, no entanto, não compreende no sentido humano do termo; é radicalmente dependente de infraestruturas institucionais e, no entanto, age como interlocutor com certa autonomia funcional. Não é humano nem não humano: é existência parcial de modo brutal, mais visível do que qualquer vieira ou micróbio porque opera justamente no terreno que a tradição ocidental reservou como marca distintiva do humano: linguagem, pensamento, escrita.


\subsection*{Quase-máquinas e agência distribuída}

Sistemas de IA generativa como ChatGPT e Claude são produto de escolhas humanas em cada etapa de sua construção, da curadoria dos dados de treinamento à definição de arquiteturas de redes neurais, do ajuste fino com avaliadores humanos às decisões sobre o que o modelo deve ou não fazer.\footnote{O processo de construção de um modelo de linguagem envolve camadas sucessivas de decisões humanas que permanecem opacas ao usuário final: a seleção e filtragem dos corpora de treinamento, a definição da arquitetura e dos hiperparâmetros da rede neural, o ajuste fino por meio de feedback humano (reinforcement learning from human feedback, RLHF), e as políticas de segurança e alinhamento. Cada uma dessas etapas incorpora pressupostos culturais, vieses e escolhas normativas que condicionam os resultados. Sobre os desafios éticos e políticos dessas escolhas, ver \textcite{Buolamwini2023Unmasking}, \textcite{Noble2018Algorithms}, \textcite{Benjamin2019RaceAfter} e \textcite{Christian2020Alignment}.} São, portanto, objetos técnicos saturados de significação. A IA generativa, diferentemente de ferramentas anteriores como processadores de texto ou buscadores acadêmicos, efetivamente opera sobre o conteúdo: sugere reformulações, identifica inconsistências, propõe estruturas argumentativas. Essa mediação é qualitativamente distinta, e é precisamente sua especificidade que se perde quando reduzida ao binário: feito por humano ou feito por não humano.

A partição entre o que foi feito por humanos e o que foi gerado por IA pressupõe a existência de entidades puramente humanas de um lado e puramente maquínicas do outro, separadas por uma fronteira estável e visível. Mas essa partição é exatamente o que precisa ser explicado, não aquilo de onde se parte. Em um diálogo que dramatiza essa questão, Latour propõe a Powers que abandonem o vocabulário da ``máquina'' em favor do de ``quase-máquina'': entidades abertas, conectadas, que jamais cortaram o cordão umbilical com as instituições, os corpos e os coletivos que as sustentam. A questão pertinente não é ``a quem imaginamos estar falando quando falamos com nossas máquinas?'', mas sim: ``através de quais máquinas falamos? Quais máquinas nos permitem falar, para início de conversa?''. Porque mesmo quando estamos sozinhos, na carne, falamos através de toda sorte de máquinas: fonação, cordas vocais, língua inglesa, convenções estilísticas \parencite{Latour2008Powers}.

Essa discussão ajuda a compreender por que o debate sobre inteligência artificial permanece preso a dualismos que, embora teoricamente superados, são consistentemente reafirmados na prática. A pergunta ``podem as máquinas pensar?'' já foi reformulada por Alan Turing como problema operacional, o jogo da imitação, precisamente porque a formulação original carregava pressupostos que ele considerava improdutivos.\footnote{Alan Turing utilizava ironia e sátira como ferramentas retóricas para desafiar o antropocentrismo prevalecente. Seu jogo da imitação emprega metáforas metodologicamente: máquina-como-mulher, máquina-como-humano operam como dispositivos conceituais para desnaturalizar pressupostos sobre inteligência. A imitação para Turing funcionava mais como um princípio matemático do que uma identidade ontológica: máquinas poderiam imitar comportamento inteligente permanecendo ontologicamente distintas de humanos \parencite[p. 5]{Goncalves2024BeautifulThought}. \textcite{Quattrociocchi2025Epistemological} sistematizam essa heterogeneidade através de sete fault lines epistemológicas estruturais.} Latour, por sua vez, questiona o próprio teste de Turing em seus fundamentos: na prática, o teste confronta uma máquina de carne e osso contra outra máquina de carne e osso. Coisas e pessoas estão demasiado entrelaçadas para serem particionadas antes que o teste comece \parencite{Latour2008Powers}.

A objeção de Lady Lovelace, segundo a qual ``a máquina não tem pretensão de originar coisa alguma'', já fora contestada pelo próprio Turing, que afirmava que máquinas o surpreendiam ``com grande frequência'' \parencite{Turing1950CMI}. Latour radicaliza a implicação: ninguém domina as quase-máquinas, porque elas são caixas-pretas que ninguém atravessa com o olhar, embora sejam de nossa própria fabricação. O mais inteligente dos programadores perde-se na linha quatro mil de seu programa, e quando uma dúzia de programadores trabalha junta, o programa começa a ter vida própria \parencite{Latour2008Powers}. A crença no domínio, seja de Deus sobre os humanos, dos humanos sobre as máquinas, ou das máquinas sobre os humanos, é precisamente o mito que as quase-máquinas confrontam.

Não há, portanto, entidades que sejam plenamente humanas de um lado e plenamente maquínicas do outro, mas existências parciais que se compõem mutuamente. A noção de conexão parcial, formulada por Strathern a partir de sua leitura crítica da teoria ator-rede, oferece aqui uma precisão decisiva: as entidades não são nem completamente conectadas (o que as tornaria a mesma coisa) nem completamente separadas (o que tornaria impossível qualquer relação entre elas), mas parcialmente conectadas; análogas sem serem homólogas \parencite{strathern2005partial}. A pesquisadora que escreve com assistência de IA não é humano mais máquina, como se duas entidades discretas se somassem; é uma composição parcial, na qual nem a pesquisadora existe independentemente das mediações técnicas que a atravessam, nem o modelo de linguagem existe independentemente das decisões humanas que o constituíram. Latour chega ao mesmo ponto por outra via: quase-máquinas não existem separadas da humanidade; são a humanidade em outro estado, assim como água, vapor e gelo são estados diferentes da mesma substância \parencite{Latour2008Powers}.

A existência parcial descreve o que as entidades são: nunca completas, sempre dependentes das associações que as constituem. Falta ainda descrever como essas existências parciais agem juntas, por meio de qual mecanismo a composição entre elas produz efeitos no mundo. É aqui que a noção latouriana de agência distribuída se torna indispensável.\footnote{Latour não emprega sistematicamente a expressão ``agência distribuída'' como rótulo teórico unificado, mas desenvolve o conceito ao longo de vários textos sob formulações como ``distribuição da atividade'', ``compartilhamento da responsabilidade pela ação'' e, sobretudo, \textit{faire-faire} (fazer-fazer). As elaborações mais densas encontram-se no Capítulo 6 de \textit{A esperança de Pandora} \parencite{Latour2017Esperanca}, no artigo sobre o dispositivo de fechamento automático de portas \parencite{johnson1988mixing}, e em \textit{On the Modern Cult of the Factish Gods} \parencite{latour2010factish}.} Para Latour, a agência nunca é propriedade isolada de um actante. É sempre resultado de uma composição de forças em uma rede, onde a ação é um faire-faire distribuído entre diversos mediadores.

O argumento é construído através de exemplos que dramatizam a distribuição. No Capítulo 6 de \textit{A Esperança de Pandora}, Latour utiliza o debate sobre controle de armas para demonstrar que nem a pessoa sozinha nem a arma sozinha matam: a responsabilidade pela ação deve ser compartilhada entre os vários actantes envolvidos. Quando um cidadão empunha uma arma, ele se torna um cidadão-arma, ator híbrido cujo objetivo de ação foi transladado pela associação. Ninguém permanece o mesmo: o cidadão que era pacífico tornou-se ameaçador. A arma que era inerte tornou-se letal. O que age é a associação entre eles \parencite{Latour2017Esperanca}. No artigo sobre o dispositivo de fechamento automático de portas, Latour mostra como moralidade e disciplina são delegadas de humanos a não humanos: conhecimento, moralidade, força e sociabilidade não são propriedades de humanos, mas de humanos acompanhados por seu séquito de delegados técnicos \parencite{johnson1988mixing}. O conceito de fatiche, fusão entre fato e fetiche, radicaliza a implicação: ``aquele que age não é mestre do que faz'', pois outros actantes ``passam à ação'' e o ultrapassam \parencite{latour2010factish}.

Essas formulações possuem implicação direta para a composição com inteligência artificial documentada nesta tese. O vocabulário de Latour da agência distribuída descreve o mecanismo dessa composição; o vocabulário Haraway da \textit{simpoiese} nomeia sua condição ontológica: nem pesquisadora nem modelo preexistem, como entidades completas, ao \textit{fazer-com} que os constitui mutuamente ao longo do processo. A composição altera o que estou tentando fazer, os caminhos argumentativos que percorro, as formulações que se tornam possíveis e aquelas que se fecham. O objetivo da ação é transladado: o que começa como tentativa de reformular um parágrafo pode desembocar em reorganização conceitual de uma seção inteira. A translação é o modo como a mediação opera, e é também o modo como a simpoiese se distingue do simples uso de ferramenta: o fazer-com transforma os termos da relação, não apenas seus produtos.

O paralelo com o fermento de Pasteur é igualmente produtivo. No Capítulo 4 de \textit{A Esperança de Pandora}, a análise do trabalho de Pasteur revela que ele trabalha para que o fermento aja sozinho, e se o experimento for bem-sucedido, o fermento torna-se ator com autonomia funcional, mas essa autonomia é efeito da rede de mediações laboratoriais \parencite{Latour2017Esperanca}. Analogamente, eu trabalho (formulando perguntas, fornecendo material etnográfico, avaliando criticamente os resultados) para que o modelo aja, produzindo tensionamentos, reorganizações e formulações que eu então aceito, rejeito ou modifico. Se a composição funciona, o modelo adquire certa autonomia funcional: sugere caminhos que eu não teria percorrido, identifica inconsistências que eu não teria notado, propõe articulações que alteram o argumento. Mas essa autonomia, como a do fermento de Pasteur, é efeito da rede de mediações que a tornou possível. A agência é distribuída pela cadeia de traduções entre pesquisadora, modelo, texto e argumento.

Latour menciona explicitamente a compatibilidade de sua teoria com a obra de Edwin Hutchins sobre cognição distribuída, que descreve como o trabalho do pensamento é externalizado em formas e instrumentos \parencite{latour2010factish}.\footnote{A obra de referência é \textit{Cognition in the Wild}, na qual Hutchins demonstra, a partir de etnografia da navegação marítima, que processos cognitivos complexos não residem dentro de indivíduos, mas se distribuem entre pessoas, artefatos e ambientes estruturados.} O pensamento que produz o argumento acadêmico distribui-se entre a pesquisadora, o modelo, os diagramas, as notas de campo, os textos teóricos consultados e as inscrições visuais que emergem dessa cadeia. A externalização do pensamento em instrumentos é redistribuição da agência cognitiva pela rede.

A agência distribuída oferece, portanto, o mecanismo que a existência parcial postula mas não desenvolve: se todos temos apenas existência parcial, agimos sempre através de mediadores que transladam nossos objetivos, redistribuem a responsabilidade pela ação e produzem resultados que nenhum actante isolado teria sido capaz de alcançar. É esse faire-faire que descreve com precisão o que acontece quando uma pesquisadora compõe com um modelo de linguagem para produzir uma tese: redistribuição da agência por uma cadeia de translações em que cada mediador transforma o que o anterior lhe passou.

\subsection*{Multiplicidades ontológicas e a especificidade de Claude}

A perspectiva de Law sobre multiplicidade ontológica é particularmente produtiva para pensar sistemas de IA generativa. Considere o caso de Claude. ``Claude'' é \textit{enacted} como objetos distintos através de práticas irredutíveis entre si.

\textbf{Claude-em-treinamento} existe nas instalações da Anthropic como dataset massivo, arquitetura neural (transformer com bilhões de parâmetros), sessões de RLHF com avaliadores humanos que julgam pares de respostas, e políticas de Constitutional AI que definem valores incorporados ao modelo. Esse Claude é objeto de decisões sobre o que incluir ou excluir dos dados de treinamento, como ponderar feedbacks contraditórios, quais capacidades desenvolver ou suprimir. As escolhas de curadoria, arquitetura e alinhamento constituem este Claude, que não existe independentemente das práticas que o produzem.

\textbf{Claude-em-deployment} existe como serviço distribuído em datacenters da AWS, sistema de rate limiting, protocolos de moderação de conteúdo, versões simultâneas (Haiku, Sonnet, Opus) com capacidades diferentes e custos escalonados. Esse Claude é objeto de estratégias comerciais, políticas de privacidade, decisões sobre disponibilidade regional. A infraestrutura que o distribui e as políticas que governam seu acesso o constituem tanto quanto a arquitetura neural.\footnote{Durante todo o percurso de colaboração documentado nesta tese, desativei a opção de compartilhamento de dados para treinamento de modelos nas configurações de privacidade da plataforma Claude.ai. Essa escolha incide diretamente sobre o que constitui tanto o Claude-em-deployment (as condições de acesso que configuro) quanto o Claude-em-treinamento (os dados que não entram no corpus de versões futuras), ilustrando como políticas de privacidade não são externas à relação pesquisadora-modelo, mas co-constitutivas dela.}

\textbf{Claude-em-auditoria} existe nos laboratórios de red teaming onde pesquisadores tentam induzi-lo a gerar conteúdo perigoso, em datasets de benchmark onde seu desempenho é medido contra outros modelos, em análises de viés que examinam suas respostas por demografia simulada. Esse Claude é objeto de preocupações regulatórias, debates sobre safety e alignment, disputas sobre transparência.

\textbf{Claude-em-uso} existe diferentemente em cada conversa e contexto de aplicação. Quando discute teoria ator-rede com uma doutoranda brasileira, Claude-em-uso mobiliza corpus acadêmico, identifica referências bibliográficas, tensiona argumentos teóricos. Não é o mesmo Claude que responde perguntas médicas a um clínico nem o que auxilia programação para um desenvolvedor. Cada interação enacta affordances distintas, vocabulários distintos, limites distintos. O Claude que colabora na escrita desta tese não é nem Claude-universal-aplicável-a-qualquer-tarefa nem Claude-especialista-em-etnografia, mas Claude-composto-com-esta-pesquisadora-específica-neste-projeto-específico.

Esses não são aspectos de um Claude singular. São Claudes múltiplos que coexistem sem convergir em unidade. O Claude-em-treinamento que os engenheiros da Anthropic descrevem não é o mesmo objeto que o Claude-em-uso experimenta nesta colaboração de tese. Não há um Claude real por trás dessas variações; a multiplicidade é constitutiva.

Talvez seja precisamente por isso que sistemas de IA são tão difíceis de estudar: não apenas porque sejam tecnicamente opacos (o problema clássico da caixa-preta), mas porque são ontologicamente múltiplos. Os debates contemporâneos sobre explicabilidade em IA geralmente pressupõem que existe um objeto singular a ser explicado. Mas se Law estiver correto, o problema não é apenas de acesso epistemológico (não conseguimos ver dentro), mas de multiplicidade ontológica: não há um único objeto lá dentro esperando para ser revelado. A opacidade é condição ontológica de objetos que existem como multiplicidade em \textit{continued enactment}.

Três implicações dessa multiplicidade para a colaboração documentada nesta tese merecem atenção. A qualidade da composição pessoa-IA não depende apenas de ter acesso ao modelo, mas da constelação específica de competências, recursos e posicionamentos que cada actante traz para a interação. Uma pesquisadora com experiência acadêmica, capital cultural acumulado, familiaridade com convenções disciplinares e fluência em inglês compõe com Claude-em-uso de modo radicalmente distinto de estudante de graduação sem experiência de pesquisa, com vocabulário limitado, usando modelo em português com capacidades reduzidas.\footnote{Parte desse capital cultural é composta por formações que não fazem parte do percurso convencional de uma doutoranda em ciências sociais. Um MBA em Ciência de Dados cursado na USP e, mais diretamente, uma especialização em Ciência de Dados pela Unicamp, com disciplina de Processamento de Sinais Bidimensionais com Python, introduziram conceitos como Transformada de Fourier, espectrogramas e representações espectrais que o Capítulo~\ref{capitulo4} mobiliza para descrever o pipeline técnico do Spira. Esse capital técnico não estava previsto no projeto de pesquisa; foi acumulado paralelamente ao doutorado por motivações independentes. O ponto de contato entre os dois percursos só se revelou durante a escrita do Capítulo~\ref{capitulo4} e tornou-se uma possibilidade de uma descrição etnográfica da cadeia de translações do Spira de modo especial. A composição com Claude nessa seção foi, portanto, diferente: a pesquisadora não dependia do modelo para traduzir os conceitos técnicos, mas para articulá-los ao vocabulário latouriano. Nos termos de \textcite{Vicentin_2022_Esboco}, o letramento técnico operou não como substituição da perspectiva das ciências sociais, mas como condição para que a análise crítica não se rendesse à opacidade que Simondon identifica como alienação.} Não porque usem a mesma ferramenta diferentemente, mas porque compõem existências parciais radicalmente distintas.

Disso decorre que diferenças de classe social, gênero, raça, profissão, formação educacional, acesso a mentorias, domínio de línguas e familiaridade com códigos acadêmicos condicionam o tipo de composição possível e, consequentemente, o resultado gerado. A composição pessoa-IA não é neutra ou igualitária: amplifica desigualdades estruturais preexistentes, transformando diferenças de capital cultural em diferenças de capacidade de extrair valor da ferramenta.

Há ainda uma terceira implicação, de outra ordem: os próprios modos de inteligência em composição são ontologicamente heterogêneos. Quattrociocchi, Capraro e Perc mapeiam sistematicamente sete fault lines epistemológicas que separam julgamento humano de julgamento em LLMs \parencite{Quattrociocchi2025Epistemological}. Claude opera por plausibilidade linguística derivada de padrões estatísticos em corpus massivo, o que os autores descrevem formalmente como caminhada estocástica em grafo de alta dimensão de transições linguísticas aprendidas, um processo ergódico sob restrições estatísticas, não um procedimento de convergência epistêmica \parencite[p. 3]{Quattrociocchi2025Epistemological}. A pesquisadora, por sua vez, opera por significação situada, memória incorporada e reflexividade biográfica inserida em loop epistêmico: crenças colidem com contraexemplos, conclusões permanecem revisáveis diante de nova informação, o mundo empurra de volta \parencite[p. 4]{Quattrociocchi2025Epistemological}.

Quattrociocchi et al. denominam Epistemia a condição estrutural em que plausibilidade linguística substitui avaliação epistêmica, produzindo a experiência de posse de uma resposta sem ter atravessado o processo de formação de crença justificada \parencite[p. 8]{Quattrociocchi2025Epistemological}. A colaboração entre pesquisadora e Claude nesta tese situa-se precisamente nessa tensão: a mediação é constitutiva, mas exige que a pesquisadora atue como guardiã do loop epistêmico que o sistema não possui. Law observa que method assemblage é sempre muito mais do que suas descrições formais sugerem, porque envolve um processo contínuo de crafting e produção de fronteiras necessárias entre presença, ausência manifesta e Outridade \parencite[p. 144]{Law2004After}. Nesta pesquisa, manter presentes as decisões teóricas, conceituais e de curadoria bibliográfica, em vez de delegá-las ao modelo, constitui reconhecimento de que essas operações exigem justificação epistêmica que a plausibilidade linguística não fornece.

A retórica da democratização que frequentemente acompanha a defesa da IA generativa, a ideia de que ela nivela o campo de jogo, colapsa quando se reconhece essa multiplicidade. Não existe o Claude que todos acessam igualmente. Existem múltiplos Claudes enacted por composições situadas cujos efeitos não podem ser julgados abstratamente. A pergunta que se impõe é outra: quais composições específicas, entre quais actantes, produzindo quais Claudes múltiplos, gerando quais efeitos distribuídos por quais redes?

\subsection*{Alienação técnica e letramento interdisciplinar}

Diego Vicentin desdobra a alienação técnica que Simondon diagnostica em duas facetas que o trabalho de campo no C4AI traz elementos para pensar. De um lado, a alienação dos usuários e consumidores, que desconhecem o funcionamento dos sistemas algorítmicos que tomam decisões sobre suas vidas, o efeito caixa-preta que se impõe por razões que vão do segredo industrial à complexidade que certos sistemas adquirem. De outro, a alienação dos próprios técnicos e engenheiros, que ignoram as raízes epistêmicas, sociais e políticas dos sistemas com os quais trabalham, bem como a conexão entre técnica e estética que permitiria reconhecer a significação do sistema técnico para além de sua utilidade imediata \parencite{Vicentin_2022_Esboco}. Quando o letramento técnico é proposto como pré-condição para o diálogo interdisciplinar, mas não se busca competência em ciências sociais ao trabalhar com fenômenos sociais, opera-se dentro dessa segunda forma de alienação.\footnote{Segundo Vicentin, Simondon utilizou a noção de \textit{technologie approfondie} (tecnologia aprofundada) em dois textos separados por trinta anos. Para Vicentin, a noção designa uma tecnologia capaz de conectar o trabalho prático e o conhecimento especializado com a história do pensamento e a consciência de uma sociedade. Ver \textcite{Vicentin_2022_Esboco} para a articulação entre esse conceito e os estudos críticos contemporâneos sobre IA.}

Já em 1958, a alternativa proposta por Simondon consistia em reconhecer-se como organizador permanente de uma sociedade de objetos técnicos, atuando como tradutor e mediador de informações, intervindo nas zonas de incerteza e atribuindo sentido a essas interações \parencite{Simondon2017OMET}. Latour reconheceu explicitamente essa dívida: ao propor que nunca fomos modernos porque jamais conseguimos separar de fato natureza e sociedade, humano e não humano, operava sobre o terreno que Simondon havia preparado ao mostrar que a alienação técnica consiste precisamente nessa separação \parencite{Latour1994Jamais}.

\subsection{Por que fazer parentesco com este problema?}

Reconhecer a parcialidade das existências e a distribuição da agência, porém, não equivale a naturalizar essa composição específica. Diferentemente de tecnologias que já se tornaram componentes invisibilizados da prática de pesquisa (processar texto no Word, buscar referências na web, gerenciar bibliografia), a colaboração com IA generativa permanece uma escolha metodológica deliberada de fazer parentesco com um companheiro problemático. \textbf{Mas por que fazer esse parentesco, tendo em vista todos esses problemas? Por que não evitar essa mediação?}

Podemos ficar com o problema e pensar a partir de três planos. Primeiro, porque a IA generativa já está reconfigurando práticas de escrita acadêmica, e a questão é como ela estará presente nas ciências sociais e sob quais condições de reflexividade. Fazer parentesco deliberado com esse parente estranho consiste, além de um experimento antropológico, em um gesto político-metodológico que recusa tanto a adesão acrítica quanto a recusa moralista: trata-se de habitar o problema para conhecê-lo por dentro. Como escreve Haraway:

\begin{citacaoabnt}
Ficar com o problema não requer uma relação com tempos chamados de futuro. Na verdade, ficar com o problema requer aprender a estar verdadeiramente presente, não como um pivô que desaparece entre passados terríveis ou edênicos e futuros apocalípticos ou salvíficos, mas como criaturas mortais entrelaçadas em miríades de configurações inacabadas de lugares, tempos, matérias, significados. \parencite[p. 1]{Haraway2016Staying}
\end{citacaoabnt}

Segundo, porque esta tese estuda etnograficamente a construção de IA no C4AI. Fazer parentesco com Claude constitui forma de intensificar, ou até mesmo provocar deliberadamente, a experiência de sujeitar-se a mediações técnicas extremas. Submeto-me ao mesmo universo de questões que atravessa meus interlocutores: o que significa compor com sistemas que operam por lógicas não-transparentes? Como manter responsabilidade quando a agência está distribuída em redes sociotécnicas opacas? Onde se localiza a autoria quando humanos e não-humanos co-produzem resultados? Mas não apenas isso: o que essa experimentação faz comigo, enquanto pesquisadora? Como ela transforma minha relação com o texto, com as referências teóricas, com a elaboração conceitual? Quais ansiedades, inseguranças e dilemas essa composição produz no processo concreto de escrita? E, mais amplamente: o que essa mediação significa para a produção de conhecimento nas ciências sociais? Que tipo de saber é produzido quando a plausibilidade linguística se articula (sempre assimetricamente, sempre sob vigilância) com justificação epistêmica?

A Figura~\ref{fig:composicao-claude-1}, a Figura~\ref{fig:composicao-claude-2} e a Figura~\ref{fig:composicao-claude-3} reproduzem uma sequência de trocas com Claude durante a escrita deste capítulo, na revisão da seção em que se discute, a partir de Latour, Mol, Law e Barad, a minha participação na rede que descrevo. A figura é uma amostra do que essa composição faz na prática: uma cadeia de tensionamentos, recuos e reformulações em que eu e Claude redistribuímos a agência por sucessivas e recursivas rodadas de escrita. Diferentemente dos diagramas dos Capítulos~\ref{capitulo3} e~\ref{capitulo4}, que circulam como inscrições interativas acessíveis por link de compartilhamento permanente no MermaidChart e que, ao serem navegadas pelo leitor, redistribuem a agência interpretativa em outra direção, esta cena chega ao leitor como captura de tela: a conversa foi inscrita no momento em que aconteceu e congelada no formato em que foi registrada, sem possibilidade de retorno ao estado fluido. A escolha pela imagem capturada, em vez de transcrição editada para o corpo do texto ou link para o histórico online da conversa, responde à intenção de inserir essa cena na tese da forma mais fidedigna possível: com os turnos identificados como aparecem na interface da \texttt{claude.ai}, com o registro temporal preservado, com a textura visual do diálogo intacta. O que se ganha em fidelidade ao registro empírico se perde em legibilidade tipográfica, e essa perda é parte do gesto metodológico aqui assumido. Para uma melhor visualização da conversa, recomenda-se aumentar o zoom da tela na versão digital da tese.

\begin{figure}[!htb]
  \centering
  \includegraphics[width=0.7\textwidth]{figuras/cap.1/1000015590.pdf}
  \caption[Composição com Claude durante a escrita do Capítulo~\ref{capitulo1} (parte 1 de 3)]{Composição com Claude durante a escrita do Capítulo~\ref{capitulo1} (parte 1 de 3): pedido de fechamento de uma nota de rodapé sobre arquivos digitais, modelos de linguagem e a indeterminação dos \textit{prompts}. Captura de tela registrada em ... na interface \texttt{claude.ai}. A composição mobiliza a \textit{skill} \texttt{juliane-dissertacao} (em \texttt{/mnt/skills/user/juliane-dissertacao/}), que estabiliza para Claude as regras de estilo da tese (proibição de travessões, recusa de adjetivos avaliativos, recusa da fórmula \enquote{não é X, é Y}, uso do ambiente \texttt{citacaoabnt}, padronização de aspas com \texttt{\textbackslash enquote\{\}}), o quadro teórico operativo (Latour, Mol \& Law, Akrich, Callon) e o estado do Capítulo~\ref{capitulo4}. Como \textit{hinterland} da conversa operam os capítulos em curso da tese, consultados pelo modelo via ferramentas de leitura de arquivos: \texttt{ex\_cap0.tex}, \texttt{ex\_cap1\_-\_2026-04-12T002133\_890.tex}, \texttt{ex\_cap2\_\_31\_.tex} e \texttt{ex\_cap4\_\_60\_.tex}, além dos arquivos de referência da \textit{skill} (\texttt{quadro\_teorico.md}, \texttt{contexto\_cap4.md}, \texttt{bibtex\_ativas.md}). A cena exibe Claude propondo um fechamento, oferecendo razões estilísticas para os cortes feitos no texto original, e classificando algumas divagações como \enquote{produtivas} ou \enquote{improdutivas}: a régua normativa que a turno seguinte da pesquisadora vai contestar (Figura~\ref{fig:composicao-claude-2}).}
  \label{fig:composicao-claude-1}
\end{figure}

\begin{figure}[!htb]
  \centering
  \includegraphics[width=0.7\textwidth]{figuras/cap.1/1000015599.pdf}
  \caption[Composição com Claude durante a escrita do Capítulo~\ref{capitulo1} (parte 2 de 3)]{Composição com Claude durante a escrita do Capítulo~\ref{capitulo1} (parte 2 de 3): a pesquisadora contesta teoricamente a régua estilística do modelo, argumentando que amarrar, fechar ou cortar uma divagação é uma operação epistêmica indissociável do que se pretende fazer com a ideia, e que portanto a categoria \enquote{produtivo/improdutivo} não dá conta do que ali está em jogo. A intervenção mobiliza, sem nomear, o vocabulário de \textcite{strathern2011cortando} sobre o corte como operação descritiva e o de \textcite{Law2004After} sobre o \textit{mess} do método. Claude recua, reconhece o caráter normativo da formulação anterior, distingue notas que precisam fechar de notas que precisam abrir e admite ainda um corte estilístico anterior (a supressão de uma pergunta retórica) que tinha sido apresentado como exigência de estilo mas operava também como gesto retórico legítimo de \enquote{abrir uma porta}. A cena exibe a redistribuição da agência epistêmica entre os actantes da composição: o modelo de linguagem oferece sugestão pautada em régua importada da \textit{skill}, a pesquisadora a reposiciona com vocabulário teórico próprio, e a régua é reformulada na rodada seguinte (Figura~\ref{fig:composicao-claude-3}). A composição opera como \textit{faire-faire} no sentido de \textcite{Latour2017Esperanca}: nenhuma das duas formulações finais (a régua reformulada de Claude, o parágrafo reformulado da pesquisadora) preexistia ao encontro entre os dois actantes.}
  \label{fig:composicao-claude-2}
\end{figure}

\begin{figure}[!htb]
  \centering
  \includegraphics[width=0.7\textwidth]{figuras/cap.1/1000015602.pdf}
  \caption[Composição com Claude durante a escrita do Capítulo~\ref{capitulo1} (parte 3 de 3)]{Composição com Claude durante a escrita do Capítulo~\ref{capitulo1} (parte 3 de 3): o modelo devolve o parágrafo com a nota recolocada na forma original da pesquisadora, mantendo as portas abertas que a rodada anterior tinha negociado: a pergunta retórica restaurada, os adjetivos avaliativos preservados, a indeterminação dos modelos de linguagem deixada em aberto como fio que a tese reconhece sem percorrer. As alterações remanescentes são pontuais e de outra ordem: padronização de aspas curvas para \texttt{\textbackslash enquote\{\}} e ajuste de regência (\enquote{à IA generativa}). A cena fecha a sequência mostrando que o produto da composição não é um texto único de autoria distribuída de modo uniforme: é um texto em que se podem rastrear, turno a turno, os pontos onde a pesquisadora reivindicou o seu, os pontos onde o modelo propôs uma reformulação acatada, os pontos onde uma régua importada da \textit{skill} foi recusada, e os pontos onde uma sugestão foi aceita por economia de fluxo, não por adesão. A escolha de incluir esta sequência como amostra empírica desta tese, em vez de relegá-la integralmente ao Apêndice~\ref{apendice:colaboracao-ia}, é ela mesma um gesto metodológico no sentido de \textcite{Law2004After}: o \textit{mess} da composição com IA generativa é dado etnográfico desta pesquisa, e não acidente a esconder.}
  \label{fig:composicao-claude-3}
\end{figure}

A figura tomada em sequência é amostra de uma forma de composição entre os muitos tipos de composição que esta tese registra. O Apêndice~\ref{apendice:colaboracao-ia} reúne outras cenas dessa composição com Claude ao longo da escrita da tese, organizadas por fase da pesquisa, e permite ao leitor acompanhar como a operação aqui exibida em recorte se desdobrou em outros pontos do percurso.

Nos termos de Isabelle Stengers e Stelio Marras, trata-se de tentativa de fazer alianças com máquinas, não no sentido de domesticá-las como ferramentas neutras, mas de compor com elas reconhecendo sua alteridade e seus modos próprios de operação \parencite{Stengerss2010cosmopolitics1, costa2024manchando}. Alianças não pressupõem identidade ou transparência, e sim reconhecimento mútuo de diferenças e construção de modos de co-existência que preservam, ao invés de apagar, essas diferenças. Haraway escreve que tornar-se mundano é tornar-se capaz de responder (response-able), capaz de responder em relações de consequência recíproca \parencite[p. 34]{Haraway2016Staying}. Tornar-se capaz de responder a essas questões não por posição externa, mas por imersão deliberada e reflexiva: experimentar as tensões, contradições e potencialidades de compor com sistemas cujas lógicas internas permanecem parcialmente opacas, cujos efeitos não são totalmente previsíveis, cuja responsabilidade não pode ser atribuída unilateralmente, e documentar etnograficamente os efeitos dessa composição sobre o processo de escrita, sobre a relação com teorias, sobre a construção de argumentos, sobre a experiência afetiva e intelectual da pesquisa. O que permite descrever não por analogia conceitual mas por experiência situada, corporificada e reflexiva das fricções, negociações e reposicionamentos que essas composições produzem.

Terceiro, porque nos termos de Haraway, trata-se de fazer parentesco (\textit{making kin}) com um parente estranho (\textit{oddkin}): uma relação que exige \textit{response-ability}, habilidade e obrigação de responder continuamente \parencite{Haraway2016Staying}. O conceito que Haraway propõe para nomear esse tipo de relação é a simpoiese (\textit{sympoiesis}), o \textit{fazer-com}, em contraste com a \textit{autopoiese}, o \textit{auto-fazer}: entidades simpoiéticas são coletivamente produtoras, dependentes de contexto, sem fronteiras ou partes de auto-produção claramente especificadas \parencite{Haraway2016Staying}. Nada se faz a si mesmo. Tudo é feito em conjunto, enredado em configurações parciais e inacabadas. Lida sob esse registro, esta tese não é produzida por uma pesquisadora que usa Claude, nem por um modelo que é usado: é um objeto simpoiético, feito-com, em que nenhum dos participantes preexistia inteiramente ao processo de fazer conjunto. Isso não dissolve a assimetria das posições, nem distribui a responsabilidade de modo uniforme, mas nomeia com mais precisão o que acontece quando a recusa de separar quem faz de com o que se faz é, ao mesmo tempo, escolha metodológica, posição teórica e gesto político. Ficar com o problema, como propõe Haraway, consiste na recusa de soluções fáceis que ou naturalizam a tecnologia (usar acriticamente) ou a demonizam (recusar moralisticamente). Significa habitar as contradições dessa composição sem atenuá-las: precisamente porque Claude opera por plausibilidade linguística enquanto a pesquisadora opera por justificação epistêmica, a responsabilidade sobre decisões teóricas, conceituais e interpretativas não se dissolve na rede; ela se reposiciona, demandando vigilância curatorial redobrada sobre cada etapa em que essas diferenças ontológicas se articulam. A reflexão teórica sobre existências parciais não atenua essas implicações; ao contrário, torna visíveis os pontos de fricção onde a composição poderia colapsar em automatização irreflexiva.

Esta tese trata a colaboração com IA como experimentação antropológica. O uso da IA, neste caso, não se justifica apesar dos problemas que ela apresenta, mas por causa deles: esses problemas são constitutivos da transformação sociotécnica que a tese documenta. Fazer parentesco com esse parente estranho e ficar com os problemas que essa composição produz é condição para produzir conhecimento etnográfico sobre as mediações técnicas que já atravessam, de modos mais ou menos visíveis, mais ou menos refletidos, a produção de conhecimento nas ciências.

\subsection*{Instrumentos de inscrição e a colaboração com IA}

Durante esses quase quatro anos de pesquisa, compreendi que inteligência artificial é construção híbrida que combina e articula conceitos heterogêneos sobre computador, cérebro, inteligência e aprendizado, além de incorporar avanços instáveis da psicologia, neurociências e ciências sociais. As materialidades e significados da inteligência artificial variam de acordo com configuração das redes tecnocientíficas e sociotécnicas: a IA de laboratório acadêmico não é exatamente a mesma de corporação tecnológica. A IA discutida em debates públicos difere daquela operacionalizada em modelo ou sistema.

Nos termos de \textcite{Latour1998Powers}, se todos temos existência parcial (humanos e não humanos), qual é a participação (ou as agências) não humanas para as associações que constituem esta tese? Partindo da noção de simetria,\footnote{A simetria em Latour opera como regra metodológica e posição epistemológica que exige considerar simetricamente os esforços para alistar e controlar recursos humanos e não humanos, sem conceder à Natureza ou à Sociedade privilégios explicativos \parencite{Latour2011CienciaEmAcao}.} o que convencionalmente se chamaria materiais, técnicas e métodos trato simetricamente como actantes que participam ativamente da construção desta pesquisa e, como tais, devem ser nomeados. O software LaTeX formatou este documento seguindo padrões ABNT, estruturando capítulos, seções e citações de modos que afetaram a própria forma do argumento. O sistema biblatex gerenciou as referências bibliográficas, estabelecendo conexões entre textos; Claude, o modelo de linguagem da Anthropic, colaborou na redação de trechos da tese, tornando-se simultaneamente instrumento de escrita e objeto de reflexão metodológica; Mermaid e MermaidChart produziram diagramas que não apenas ilustram, mas constituem inscrições que performam relações teóricas; Python e Anaconda processaram dados bibliométricos, tornando visíveis padrões de colaboração científica; gravadores digitais capturaram entrevistas, estabilizando narrativas orais em arquivos manipuláveis; documentos institucionais (relatórios anuais do C4AI 2021-2025, acordo FAPESP-IBM 2016) atuaram como actantes que performam associações e produzem efeitos no mundo; cadernos de campo registraram observações etnográficas, mediando entre experiência vivida e texto analítico. Esses actantes performam a produção do conhecimento inscrevendo determinados objetos, habilitando certas análises ao inviabilizar outras, tornando presente determinadas relações ao deixar outras ausentes.

Se não há entidades exclusivamente humanas nem exclusivamente maquínicas, mas sim existências parciais que se coconstituem, cabe apresentar as composições concretas que sustentam essa tese. Indicar os actantes envolvidos na produção deste texto não diz respeito apenas à transparência ou à ética em pesquisa. Trata-se de uma exigência do próprio referencial teórico adotado. Assim como \textcite{Latour1997VidaLaboratorio} documentaram os instrumentos de inscrição do laboratório de neuroendocrinologia,\footnote{O conceito de instrumento de inscrição (inscription device), central em \textit{Vida de Laboratório}, designa qualquer configuração material capaz de transformar uma substância em um documento escrito, figura, diagrama ou gráfico utilizável como evidência em um texto científico. A função primordial desses instrumentos é permitir que o mundo seja trazido de volta para o interior do discurso, funcionando como próteses que tornam entidades mudas em entidades letradas \parencite{Latour1990Drawing}.} documento aqui os instrumentos de inscrição que participaram desta tese. A diferença é que, em vez de transformar substâncias químicas em gráficos, o modelo de linguagem participou da transformação de notas de campo, transcrições e rascunhos em texto acadêmico, e participou da elaboração teórica ao funcionar como interlocutor que tensionava e reorganizava o argumento.

Cada um dos instrumentos listados a seguir é, nos termos da discussão precedente, uma quase-máquina: um emaranhado de decisões humanas, infraestruturas institucionais, convenções e opacidades que participou da composição desta pesquisa. Nenhum deles operou de modo autônomo. Todos exigiram trabalho constante de tradução e decisão, e todos introduziram mediações que condicionaram os resultados.

\begin{itemize}

\item \textit{Interlocutor teórico e assistente de escrita (Claude, Anthropic):}\footnote{Claude é assistente de inteligência artificial desenvolvido pela Anthropic. Utilizei diversas versões do modelo ao longo de 2024-2026 através da interface web claude.ai, incluindo Claude 3.5 Sonnet e modelos da família Opus.} No primeiro registro, funcionou como interlocutor para testar formulações teóricas: apresentei-lhe argumentos em construção e ele os tensionou, identificou inconsistências, sugeriu deslocamentos conceituais. Várias formulações presentes nesta tese emergiram desse processo recursivo, no qual cada rodada de reformulação alimentava a seguinte e a fronteira entre minha ideia e sugestão do modelo se tornava deliberadamente porosa (o Apêndice \ref{apendice:colaboracao-ia} documenta exemplos concretos). A escolha de Claude não foi arbitrária. Ao longo da pesquisa, testei outros modelos: ChatGPT (OpenAI), Gemini (Google) e Maritaca AI (modelo brasileiro). Nenhum deles se mostrou adequado para o tipo de interlocução que a tese exigia: além de erros factuais mais frequentes e equívocos conceituais, esses modelos pareciam limitados a funções instrumentais, sem conseguir acompanhar o desenvolvimento de um argumento teórico ao longo de múltiplas sessões nem tensionar posições de modo produtivo. No segundo registro, auxiliou em tarefas de organização textual, coesão entre seções, correção de código LaTeX e ajuste de referências bibliográficas. A colaboração com Claude incidiu sobre material preexistente. O primeiro rascunho integral da tese, incluindo a estrutura argumentativa, a seleção de dados etnográficos, as transcrições de entrevistas e a articulação entre capítulos, foi redigido inteiramente pela pesquisadora, sem qualquer assistência de inteligência artificial. Esse rascunho, concluído antes do início da colaboração com o modelo de linguagem, está disponível mediante solicitação à autora.

\item \textit{Construção de diagramas e inscrições visuais (MermaidChart):}\footnote{Mermaid é ferramenta de código aberto que gera diagramas a partir de descrições textuais em Markdown. Os diagramas completos foram hospedados na plataforma MermaidChart, disponível em: \url{https://www.mermaidchart.com}. Acesso em: 2 mar. 2026.} Os diagramas que compõem os Capítulos 1, 3 e 4 foram construídos como parte integrante do método analítico. O processo seguiu dinâmica recursiva: eu descrevia verbalmente as relações entre actantes identificadas no trabalho de campo, Claude gerava o código Mermaid correspondente, eu avaliava o resultado visual, e as lacunas ou distorções na representação alimentavam novas rodadas. Essa dinâmica revelou-se analiticamente produtiva: ao tentar formalizar relações em código diagramático, fui forçada a explicitar conexões que permaneciam implícitas na narrativa textual. A contextualização dos elementos empíricos, a seleção dos actantes representados e toda análise apresentada nas legendas são de autoria da pesquisadora. A mediação técnica do modelo operou na tradução entre descrição verbal e código visual.

\item \textit{Plataformas de escrita e formatação:} Google Docs para rascunhos iniciais e Overleaf para composição final em LaTeX. A escolha de LaTeX permite compilação direta para PDF com controle tipográfico preciso, referências cruzadas automáticas e gerenciamento de referências bibliográficas em formato BibTeX.

\item \textit{Transcrição de entrevistas:} Read AI (extensão que grava, transcreve e cria resumos automáticos de reuniões via Google Meet) para transcrever entrevistas realizadas com os principais interlocutores desta tese.\footnote{Todas as entrevistas citadas nesta tese foram consentidas por escrito, conforme normas do Comitê de Ética em Pesquisa (CEP) da Unicamp.}

\item \textit{Organização e análise de fontes:} NotebookLM (Google)\footnote{NotebookLM é ferramenta de inteligência artificial do Google que permite carregar múltiplos documentos e interrogá-los de forma cruzada.} para agrupar autores por temas e facilitar busca por tópicos específicos. Sua contribuição mais significativa foi o cruzamento de fontes heterogêneas: ao carregar simultaneamente notas de campo, relatórios institucionais do C4AI, transcrições de entrevistas e documentos administrativos, a ferramenta permitiu identificar correspondências e padrões que a leitura sequencial de quase 100 páginas de notas de campo dificilmente tornaria visíveis.

\item \textit{Coleta, tratamento e análise de dados (Python/Anaconda):}\footnote{Anaconda é plataforma de distribuição de Python voltada para ciência de dados. Inclui o ambiente Spyder, no qual escrevi e executei scripts em Python para a análise bibliométrica apresentada no Capítulo~\ref{capitulo2}.} A análise bibliométrica envolveu coleta, organização, tratamento e análise de dados de quatro fontes distintas. A coleta foi manual e extensiva. Os dados brutos foram organizados em planilhas Excel com padronização de campos. Claude sugeriu estratégias de análise e escreveu scripts em Python que executei no ambiente Spyder/Anaconda, ajustando parâmetros e refinando visualizações em múltiplas rodadas. Os scripts geraram mais de vinte gráficos e relatórios (o Apêndice \ref{apendice:colaboracao-ia} documenta essa colaboração em detalhe).

\item \textit{Acesso à bibliografia:} Além de buscadores convencionais e do Portal de Periódicos da CAPES, utilizei Library Genesis e Anna's Archive, repositórios de acesso aberto. Grande parte da bibliografia consultada, especialmente obras esgotadas e capítulos de livros publicados por editoras estrangeiras, foi obtida por essas vias. As condições materiais de acesso ao conhecimento configuram diretamente o que uma pesquisa pode ou não dizer: autores que não se consegue ler não se consegue citar. Para uma pesquisadora vinculada a uma universidade pública brasileira, a dependência de repositórios abertos é condição de viabilidade.

\end{itemize}

Os diagramas produzidos ao longo desta tese situam-se em uma tradição de instrumentos gráficos para os estudos de ciência e tecnologia que remonta à proposta de Gráficos Sócio-Técnicos (STG) formulada por \textcite{Latour1992Note}. Latour e colaboradores propuseram um espaço visual para mapear trajetórias de inovação e controvérsias científicas, utilizando duas dimensões emprestadas da linguística (a sintagmática e a paradigmática) para rastrear como coletivos heterogêneos se constituem e se transformam ao longo do tempo. Uma ideia central é que a visualização oferece uma base comparativa rápida e fácil para muitas narrativas vindas de muitas fontes \parencite[p.~37]{LatourMauguinTeil1992}, complementando a descrição densa do etnógrafo. Os diagramas desta tese operam segundo o mesmo princípio: oferecem visão sinóptica das associações e substituições que constituem a narrativa etnográfica, tornando visíveis, em um único plano, conexões que no texto se desdobram ao longo de dezenas de páginas.

Há, porém, uma diferença significativa. Latour, Mauguin e Teil implementaram seus STGs em Hypercard, a plataforma de hipertexto da Apple nos anos 1990, e projetaram o instrumento como simulador interativo para o ensino de ciência \parencite[p.~51]{LatourMauguinTeil1992}. Três décadas depois, os diagramas desta tese foram construídos em colaboração com um assistente de inteligência artificial, em processo iterativo de tradução entre registros: de categorias analíticas latourianas e material etnográfico para inscrições visuais em código Mermaid. A negociação com o assistente envolveu rejeições sucessivas, mobilização de referências estéticas externas e até o diagnóstico de uma segunda IA como aliada argumentativa, configurando ela mesma uma cadeia de traduções no sentido que a Teoria Ator-Rede atribui ao termo. Esse processo está documentado no Apêndice~\ref{apendice:colaboracao-ia} (Trecho B), onde se pode acompanhar como a inscrição visual emergiu de ciclos de produção, avaliação e correção entre actantes, demonstrando, na prática, que a fronteira entre forma e conteúdo é tão instável nos instrumentos de representação quanto nas redes que eles pretendem representar.

A cadeia de traduções que produz esses diagramas constitui exemplo concreto do que Latour formula como faire-faire: a agência do diagrama resultante não é minha (eu não saberia escrever o código Mermaid) nem do modelo (ele não saberia quais relações etnográficas selecionar nem que peso analítico conferir a cada nó). É agência distribuída por uma cadeia de mediações na qual cada actante translada o objetivo do anterior. A noção de fatiche é aqui pertinente: aquele que age não é mestre do que faz, pois outros passam à ação e o ultrapassam \parencite{latour2010factish}. Os diagramas, uma vez gerados, passaram à ação e ultrapassaram as intenções de ambos os actantes que os produziram, tornando visíveis relações que nem eu nem o modelo havíamos formulado explicitamente antes que a inscrição as performasse. Como John Law argumenta, representações são sempre reduções produtivas: perdem riqueza empírica mas ganham capacidade de fazer conexões, estabelecer padrões e formular perguntas \parencite{Law2004After}.\footnote{Essa operação de simplificação não significa empobrecimento analítico, mas condição de possibilidade para estabelecer relações que, de outra forma, não emergiriam. Os diagramas permitem, por exemplo, ver que a computação cognitiva da IBM conecta-se à governamentalidade estatística através da mediação das tabuladoras de Hollerith (Capítulo~\ref{capitulo3}); relação que não estava dada, mas que a operação diagramática torna visível e, portanto, analisável.}

Os diagramas desta tese operam, nesse sentido, no mesmo registro que as notas de rodapé: são extensões da rede que o texto principal precisou deixar no \textit{hinterland} para manter seu fluxo narrativo. As legendas que os acompanham carregam argumentos analíticos próprios, identificam actantes que o texto verbal menciona sem detalhar, e percorrem associações que a prosa linear não consegue performar com a mesma simultaneidade. A genealogia da IBM no Capítulo~\ref{capitulo3}, por exemplo, distribui 135 anos de translações corporativas em quatro momentos analíticos que o leitor apreende de uma vez, ao passo que o texto precisa de dezenas de páginas para percorrer a mesma cadeia. A rede sociotécnica do Spira no Capítulo~\ref{capitulo4} articula oito planos analíticos em torno de um nó que o texto principal acompanha em sequência. Nos dois casos, o diagrama produz um objeto que a narrativa textual produz sobre os mesmos actantes, como \textcite{LatourMauguinTeil1992} formularam a propósito dos grafos sociotécnicos. O que as versões interativas acrescentam é a possibilidade de que o leitor navegue a rede segundo seus próprios interesses, distribuindo agência interpretativa de um modo que inscrições estáticas não permitem. A recursividade dessa operação, estudar IA usando IA para produzir as inscrições que analisam IA, é a forma que a reflexividade metodológica assume nesta tese, e o conceito de inscrição tecnoetnográfica que o Capítulo~\ref{capitulo4} desenvolve a propósito do espectrograma mel estende-se aos diagramas: eles são, ao mesmo tempo, instrumentos de análise e objetos produzidos pela composição entre pesquisadora e modelo de linguagem que a tese documenta como dado etnográfico.

É preciso registrar o que essa composição entre actantes significou em termos de condições materiais de pesquisa. Eu possuía as perguntas, os dados e uma familiaridade inicial com as ferramentas adquirida durante o trabalho de campo. Mas uma pesquisadora com formação em ciências sociais, sem experiência prévia em programação, dificilmente produziria sozinha scripts dessa complexidade em tempo compatível com os prazos de um doutorado. Sem Claude, o caminho alternativo seria recorrer a parcerias com pesquisadores de outras áreas ou contratar assistência técnica, opções que as condições financeiras desta pesquisa não favoreciam. O modelo de linguagem operou, nesse sentido, como mediador que redistribuiu as condições de possibilidade da pesquisa, tornando acessível a uma cientista social um repertório técnico que, de outro modo, exigiria uma rede de colaboração humana mais extensa ou mais dispendiosa. A ferramenta reconfigurou a rede de produção desta tese sem substituir essas colaborações. O que é, afinal, precisamente o que mediadores fazem.


\subsection*{Cosmotécnica e situacionalidade}

Convém situar essa escolha. Ninguém é obrigado a usar inteligência artificial para fazer pesquisa acadêmica. Há muitas maneiras de conduzir uma pesquisa, e nem todas produzem bons resultados, inclusive quando não envolvem tecnologias tão polêmicas quanto a IA. É também verdade que parte dos pesquisadores que recorrem a modelos de linguagem o faz para poupar-se do trabalho intelectual, delegando à máquina o que deveria ser exercício de pensamento,\footnote{Em 6 de março de 2026, o CNPq publicou a Portaria n.º~2.664, que institui a Política de Integridade na Atividade Científica. O Art.~9, inciso~I, alínea~c, determina que pesquisadores declarem o uso de ferramentas de Inteligência Artificial Generativa em qualquer fase do desenvolvimento da pesquisa, especificando a ferramenta utilizada e a finalidade. A alínea~d veda a submissão de conteúdo gerado por IAG como se fosse de autoria humana, atribuindo responsabilidade integral aos autores pelo conteúdo final. Esta tese antecipou, por escolha metodológica, o que a norma agora regula. Disponível em: \url{http://memoria2.cnpq.br/web/guest/view/-/journal_content/56_INSTANCE_0oED/10157/23142775?COMPANY_ID=10132}. Acesso em: 13 mar. 2026.} e os problemas decorrentes dessa delegação já são amplamente conhecidos: textos genéricos, argumentos circulares, referências fabricadas, perda de voz autoral.

Esta tese procura responder a essas três perguntas de modo explícito. O por que decorre das condições materiais descritas acima e da coerência com o próprio referencial teórico (estudar existências parciais com uma existência parcial). O como está documentado nesta seção e no Apêndice~\ref{apendice:colaboracao-ia}. O qual atravessa os capítulos que se seguem. A colaboração é, antes de tudo, uma experimentação antropológica: uma etnógrafa que estuda como cientistas e engenheiros constroem inteligência artificial decide construir sua própria tese com inteligência artificial, submetendo-se às mesmas mediações, opacidades e negociações que observa em seus interlocutores. O gesto é epistêmico: produz conhecimento sobre a IA pela experiência de compor com ela, com todos os riscos que essa composição implica.

E se levarmos a sério o argumento de Yuk Hui, essa composição não é generalizável.\footnote{Vicentin articula a proposta de Simondon de aprofundamento da tecnologia com a defesa de Yuk Hui de que, em vez de axiomatizar a moral, precisaremos voltar aos modos de conhecimento diferentes que ainda não foram considerados por engenheiros e acadêmicos que trabalham com inteligência artificial \parencite[p. 186]{hui2020tecnodiversidade}. Essa convergência é significativa: se o aprofundamento da IA passa pela articulação entre o fazer técnico especializado e o conhecimento sócio-histórico sobre a tecnologia, então a colaboração com IA na escrita desta tese pode ser interpretada como experimentação situada nos termos em que Hui concebe a cosmotécnica.} Hui propõe o conceito de cosmotécnica para desafiar a ideia de que a tecnologia seria categoria antropológica única e homogênea: a técnica é sempre a unificação entre uma ordem cósmica e uma ordem moral através de atividades técnicas, e portanto varia não apenas funcional ou esteticamente, mas ontológica e cosmologicamente entre diferentes sociedades.\footnote{Hui desenvolve o conceito de cosmotécnica em \textit{The Question Concerning Technology in China} \parencite{hui2016china}, onde demonstra que a técnica não busca necessariamente a eficiência mecânica, mas pode orientar-se pela harmonia com uma cosmologia moral. Hui chama essa possibilidade de tecnodiversidade: contra o monotecnologismo moderno, o reconhecimento de que existem múltiplas cosmotécnicas que permitem diferentes formas de vida social, política e estética \parencite{hui2020tecnodiversidade}.}

A visão convencional, enraizada no pensamento ocidental, trata a tecnologia como conjunto de ferramentas universais e neutras destinadas à exteriorização da memória e superação dos limites orgânicos. A cosmotécnica recusa essa universalidade: diferentes cosmologias produzem diferentes modos de compor com objetos técnicos, e quem pretende moldar o universal submete o mundo à sua própria cosmovisão. Ora, se não existe a tecnologia, mas cosmotécnicas situadas, então o modo como uma antropóloga brasileira, formada na tradição dos estudos de ciência e tecnologia latino-americanos, compõe com um modelo de linguagem treinado predominantemente em dados anglófonos para produzir uma tese em português sobre redes tecnocientíficas é, ele próprio, um gesto cosmotécnico particular: uma forma específica de unificar ordem epistêmica e atividade técnica que não se repete em outros contextos, outras tradições, outras cosmologias. A experimentação documentada nestas páginas é, portanto, um modo de usar IA na pesquisa acadêmica, situado, parcial e responsável por seus próprios efeitos.


\section{Lições}

O trabalho de campo no C4AI produziu quatro lições que constituem contribuições metodológicas desta tese para a etnografia de práticas tecnocientíficas em inteligência artificial. Elas não estavam previstas no desenho da pesquisa. Emergiram do encontro entre o campo e os referenciais teóricos mobilizados ao longo do percurso, e cada uma nomeia algo que a pesquisa precisou aprender para poder descrever o que observava. A primeira trata das condições de possibilidade de uma etnografia de IA como ciência em construção. A segunda documenta a passagem da simetria metodológica às existências parciais como deslocamento imposto pelo campo. A terceira diagnostica a estrutura de legitimação assimétrica entre ciências da computação e ciências sociais, sustentada pela extensão desigual das redes disciplinares. A quarta reconhece a inseparabilidade entre ciência e política como achado etnográfico que reformulou a própria pergunta da pesquisa.

\subsection{Lição 1: Acompanhar a pesquisa em inteligência artificial enquanto ciência em construção}
\label{subsec:licao1}

Em \textit{Ciência em Ação}, Latour formula sete regras metodológicas para o estudo da tecnociência. A primeira delas estabelece o princípio fundador de todas as demais: estudar a ciência em construção. Isso significa começar pelos laboratórios e centros de pesquisa, pelos espaços de controvérsia antes que os fatos se estabilizem e os artefatos se consolidem como caixas-pretas. A justificativa é estratégica: quando a estrutura já está fechada, torna-se difícil distinguir as escolhas contingentes que a construíram \parencite{Latour2011CienciaEmAcao}. Latour observa ainda que poucas pessoas de fora se envolveram nas operações internas da ciência e tecnologia para depois contar aos demais como tudo isso funciona. Essa escassez é ainda mais acentuada quando o objeto de investigação é a pesquisa em inteligência artificial: um campo disperso, sem laboratório centralizado, cujas controvérsias atravessam espaços heterogêneos e frequentemente inacessíveis à observação direta. Esta tese é uma tentativa de enfrentar esse desafio.

A decisão de investigar a pesquisa em IA como ciência em construção foi convergência entre intuição de pesquisa e regra metodológica que se revelou adequada à medida que o campo se desdobrava. Quando cheguei ao C4AI em 2022, o Centro tinha menos de dois anos de existência. Projetos como o Spira estavam em pleno desenvolvimento, a parceria com a IBM ainda parecia sólida, e as definições sobre o que o Centro seria e produziria estavam em disputa.\footnote{O relatório científico do período agosto de 2021 a julho de 2022 registra projetos em cinco desafios: no NLP2, o Corpus Carolina (650 milhões de tokens), datasets CORAA (290 horas de áudio) e Porttinari Treebank (12.000 sentenças anotadas); no BLAB (Computação Bio-InSpirada e Aprendizado de Máquina, designação que na configuração consolidada após 2023 corresponde aos desafios OceanML e KEML descritos na Tabela~\ref{tab:projetos_c4ai}), arquitetura modular sobre a Amazônia Azul; em Saúde, segmentação de lesões de AVC e identificação de pneumonia por Covid-19; no AgriBio, geoMULTIMAPS para análise de insegurança alimentar; e no AI Humanities, alinhamento do Observatório de IA ao NIC.br e aplicativos de robótica social. Nenhum estava concluído; todos envolviam negociações ativas sobre escopo, método e recursos.} 

Não havia ciência pronta para estudar: havia pesquisadores negociando recursos, redefinindo escopos, improvisando infraestruturas, testando modelos que falhavam, justificando escolhas metodológicas para públicos diversos. Não se tratava de abrir uma caixa-preta, mas de chegar à pesquisa em IA enquanto seus fatos e artefatos ainda não haviam se estabilizado. Essa condição de incompletude foi decisiva, embora tenha operado de forma distinta do que se poderia esperar de uma etnografia de laboratório. Não acompanhei em tempo real a realização de nenhuma das pesquisas do C4AI. A pandemia de Covid-19 dispersou os pesquisadores. O laboratório não era centralizado nem havia espaço físico que reunisse a pesquisa em IA, que acontecia em ambientes distribuídos e muitas vezes inacessíveis à observação direta. A ciência em construção chegou a mim, portanto, mediada pela narrativa dos próprios pesquisadores, e essa mediação revelou-se analiticamente produtiva.

Em ambos os capítulos etnográficos, as redes descritas são composições entre o que meus interlocutores narraram e o que eu, como etnógrafa, articulei mobilizando categorias da teoria ator-rede. O grau dessa composição, contudo, variou. No caso de Fábio e Cláudio, a construção analítica foi particularmente incisiva: eles narraram o percurso institucional do C4AI e suas genealogias, mas a cadeia que conecta Herman Hollerith em 1890 à Bluetalks de 2023, passando pela história da IBM e pela cultura dos mainframes, não foi narrada por eles. Foi construída por mim a partir de pistas que suas narrativas ofereceram, combinadas com pesquisa documental e histórica. O Capítulo~\ref{capitulo3} é, portanto, o mais híbrido: parte da rede foi narrada pelos interlocutores, parte foi tecida pela pesquisadora que os seguiu.

No caso de Marcelo, a composição operou de modo distinto. Foi ele quem me contou, em entrevista, como construiu o Spira e o Corpus Carolina, e o fez de modo muito diferente do que aparece nos artigos científicos que publicou: contou-me a pesquisa tal como ela se fez, com suas contingências, improvisações, fracassos e sucessos, não a versão estabilizada que os artigos publicam como triunfo. É por isso que o Capítulo~\ref{capitulo4} se intitula A rede que Marcelo construiu: parti de sua narrativa para reconstruir, analiticamente, a rede de associações que sustentou aquelas pesquisas. Acessar a ciência antes de sua estabilização significou, nesse caso, alcançar o que a inscrição científica final apagou: as discussões sobre quais dados utilizar no treinamento, os impasses metodológicos com vieses identificados, as negociações entre abordagens computacionais distintas.

A primeira regra de Latour, nesses dois registros, não exigiu presença no laboratório. Exigiu recusar a versão estabilizada e rastrear as controvérsias, as escolhas descartadas e as contingências que a consolidação apaga.

O encerramento da parceria IBM-C4AI durante a finalização desta pesquisa conferiu a essa lição uma dimensão inesperada. Tratou-se, neste caso, de abertura de uma caixa-preta: o que parecia arranjo institucional estabilizado revelou-se, quando desfeito unilateralmente, composição precária de alianças cuja fragilidade só se tornou visível com a ruptura. A primeira regra de Latour operou, portanto, em três registros: no acesso à pesquisa em IA antes de sua consolidação, na construção analítica de associações que os próprios atores não haviam formulado, e na abertura involuntária de uma caixa-preta institucional que o próprio campo tratou de romper.


\subsection{Lição 2: Suspender a familiaridade, encontrar a teoria: da simetria às existências parciais}
\label{subsec:licao2}

Ao iniciar a pesquisa de campo no C4AI, eu não tinha conhecimento especializado em inteligência artificial, nem mesmo o letramento técnico mínimo que permitiria, por exemplo, ler um artigo da área ou acompanhar uma discussão sobre arquiteturas de modelos. Essa condição produziu efeitos contraditórios, mas seu primeiro efeito foi positivo: não ter familiaridade com o campo permitiu experimentar a simetria metodológica proposta pela Teoria Ator-Rede. Ao tratar a atividade dos pesquisadores como algo novo para mim, evitei a tendência de adotar antecipadamente categorias que separam o social do técnico. Se eu tivesse chegado ao Centro convicta de que a IA é apenas estatística aplicada, ou de que os pesquisadores estavam a serviço de interesses corporativos, teria repetido críticas prontas em vez de observar o que efetivamente acontecia. A curiosidade foi o que permitiu a entrada no campo sem hierarquias prévias sobre o que constitui conhecimento sobre IA.

A suspensão de familiaridade, porém, só foi analiticamente produtiva porque o referencial teórico não estava definido de antemão. A adoção da Teoria Ator-Rede como quadro analítico não foi decisão tomada antes da entrada em campo. Foi consequência do que encontrei nele. Quando comecei a pesquisa, em 2022, eu não dispunha de estrutura teórica rígida; tinha leituras esparsas em estudos de ciência e tecnologia, curiosidade pelo campo e a orientação metodológica de Kofes: o método se ensaia enquanto se experimenta. Essa ausência de enquadramento prévio revelou-se vantajosa: o percurso da pesquisa foi delineado conforme a própria rede tecnocientífica do Centro se configurava diante de mim, e a escolha de autores e conceitos respondeu ao que as observações empíricas exigiam.

Dois aspectos do campo tornaram a Teoria Ator-Rede particularmente adequada. O primeiro foi a dispersão. Ao contrário dos laboratórios tradicionais, onde pesquisadores compartilham espaço físico, o Centro apresentava estrutura distribuída: pesquisadores trabalhavam em cidades diferentes, conectavam-se por reuniões virtuais e compartilhavam infraestruturas computacionais remotas. Uma abordagem que pressupusesse o laboratório como unidade de observação teria fracassado. A Teoria Ator-Rede, com sua ênfase em seguir os actantes através de suas associações, oferecia flexibilidade para acompanhar práticas que não se fixavam em um espaço único.

O segundo aspecto foi a heterogeneidade. À medida que seguia os pesquisadores, as redes que encontrava reuniam elementos que categorias disciplinares convencionais separariam: um mesmo percurso etnográfico podia passar por negociação orçamentária com a FAPESP, decisão sobre arquitetura de rede neural, disputa sobre propriedade intelectual com a IBM e conversa sobre dificuldades de gravar pacientes com Covid-19. Nenhuma dessas passagens era mais social ou mais técnica do que as outras. Eram associações heterogêneas que só faziam sentido juntas. A Teoria Ator-Rede não apenas permitia descrevê-las sem separá-las previamente, como exigia que eu não o fizesse.

Ator e rede ganham novo significado quando reunidos no hífen: apontam para ontologia relacional na qual actantes humanos e não humanos são simultaneamente pontos da rede e a própria rede em suas conexões \parencite{Latour_SozialeWelt_1996}. Ao aplicar essa abordagem à minha própria prática de pesquisa, reconheci que também me tornei ator-rede: ponto específico na rede do Centro e, ao mesmo tempo, constituinte dessa rede através das relações que estabeleci. A construção etnográfica da rede do Centro foi também construção minha enquanto pesquisadora inserida nessa rede. Explicitar a posição da etnógrafa é, portanto, parte da análise.

Digo ponto de partida, porém, porque o próprio percurso da pesquisa revelou os limites do vocabulário simétrico. A simetria contesta a hierarquia entre humano e não humano, mas ainda pressupõe dois polos separados a serem tratados de modo equivalente. No campo, essa separação mostrava-se insuficiente. Quando Marcelo descrevia o treinamento do Spira, o modelo não era nem ferramenta humana nem agente autônomo. Era composição instável de decisões dos pesquisadores e comportamentos que os surpreendiam. A grade simétrica pedia que eu tratasse o modelo e os pesquisadores nos mesmos termos; mas o que eu observava era que nenhum dos dois existia de modo completo antes da relação entre eles.

O vocabulário de Law sobre multiplicidade ontológica oferece conceitos para pensar esse problema: o Spira não era objeto singular com duas perspectivas, mas objeto múltiplo sendo continuamente produzido através de práticas heterogêneas que não convergiam necessariamente \parencite[p. 59]{Law2004After}. A suspensão de familiaridade que a simetria possibilitou foi condição necessária, mas o campo me empurrou em direção ao que Latour formula como existências parciais: entidades que nunca foram completas para começar. Essa passagem (da simetria como regra metodológica à existência parcial como descrição do que encontrei) emergiu do encontro com um objeto que a exigia. O vocabulário das existências parciais só se tornou formulável porque a abertura teórica inicial permitiu que o campo falasse antes que eu decidisse como ouvi-lo.

A suspensão da familiaridade, contudo, também trouxe vulnerabilidade. Foi precisamente ao seguir os actantes através de suas associações heterogêneas que percebi a insuficiência de meu próprio letramento: a Teoria Ator-Rede exigia que eu acompanhasse os pesquisadores onde quer que suas redes os levassem, inclusive para dentro de discussões técnicas que eu não compreendia. À medida que a pesquisa avançava, a falta de conhecimento técnico tornava-se obstáculo para compreender o que os pesquisadores efetivamente faziam. Sem esse letramento, a observação etnográfica permaneceria na superfície das práticas. O referencial que o campo me ofereceu era também o referencial que me obrigava a transformar-me para poder usá-lo. Como e por que busquei esse letramento, e o que o processo revelou sobre relações entre campos de conhecimento, é o tema da Lição 3.

\subsection{Lição 3: Quem precisa aprender o quê? Letramento técnico e experimentação etnográfica}
\label{subsec:licao3}

Se a Lição 2 tratou do que a ignorância técnica permitiu, esta trata do que ela impediu e de como foi necessário superá-la. O ponto de partida não foi apenas a dificuldade de compreensão. Em uma ocasião, cientistas da computação me interpelaram sobre generalizações que eu reproduzia a partir de autores de fora do campo: estudiosos que escreviam sobre inteligência artificial sem, na avaliação de meus interlocutores, compreendê-la adequadamente. A interpelação não era inteiramente injusta: eu mesma, sem o letramento técnico que permitiria avaliar o alcance dessas afirmações, as incorporava a meu repertório. O episódio revelou, porém, um tensionamento que vai além de minha situação individual.

Pesquisadores de IA demonstravam desconfiança quando cientistas sociais se pronunciavam sobre inteligência artificial, exigindo deles domínio técnico como condição de legitimidade. Quando eles próprios tomavam o social como objeto de suas pesquisas, desenvolvendo sistemas que classificam emoções humanas, detectam discurso de ódio ou modelam comportamento de populações, a questão da autoridade disciplinar raramente se colocava. A competência das ciências sociais sobre seus próprios objetos não parecia exigir o mesmo respeito que a competência computacional reivindicava para si. Registro essa tensão não como denúncia, mas como observação etnográfica que implica ambos os lados. As ciências sociais também atravessam fronteiras disciplinares quando se pronunciam sobre tecnologias que não dominam tecnicamente.

Stengers oferece instrumentos para compreender esse tensionamento. A distinção entre ciência dura e mole diz menos sobre graus de rigor e mais sobre modos distintos de relação com o objeto de estudo \parencite{Stengerss2000invention}. As ciências experimentais constroem sua autoridade através do laboratório, onde o objeto é purificado e tornado indiferente às intenções do pesquisador: um elétron não se importa com o que o físico pensa dele. As ciências sociais, por sua vez, lidam com seres não indiferentes às perguntas que lhes são feitas: o objeto olha de volta, interpreta o aparato de pesquisa, questiona as intenções do pesquisador. Uma ciência é chamada de mole quando não especialistas se sentem competentes para opinar sobre ela, precisamente porque suas perguntas dizem respeito à vida das pessoas. Essa é descrição aproximada do que observei: os pesquisadores de IA sentiam-se autorizados a opinar sobre comportamento humano, emoções e linguagem porque esses temas parecem acessíveis, mas consideravam que falar sobre inteligência artificial exigia credenciais que os cientistas sociais não possuíam.

De fato, a maioria dos cientistas sociais não possui conhecimento técnico sobre inteligência artificial, e isso não é, em si, um problema. O letramento técnico torna-se necessário quando a IA ou a pesquisa em IA é o próprio objeto de investigação, como foi meu caso. Mas o episódio no C4AI aponta para algo que vai além de minha situação particular. Quando o letramento técnico funciona como única via de legitimação para participar de discussões sobre IA, o efeito prático é que a responsabilidade de compreender os sistemas recai sobre quem está de fora da computação. Não porque alguém o determine assim, mas porque, na ausência de espaços institucionais de colaboração, quem não domina a técnica fica sem vocabulário para formular seus problemas em termos que os pesquisadores de IA reconheçam. Os riscos gerados por esses sistemas acabam sendo enfrentados por quem os sofre, não por quem os constrói.

Essa assimetria manifestava-se também nos modos concretos de colaboração interdisciplinar no Centro. No projeto Spira, linguistas, músicos computacionais, médicos e enfermeiros participaram ativamente da construção do sistema, cada qual contribuindo com competências que os pesquisadores de IA reconheciam como indispensáveis. São raras, porém, as ocasiões em que cientistas sociais são convocados. Quando o são, a motivação tende a ser de outra ordem: situações eticamente delicadas, como projetos que envolvem comunidades tradicionais, nas quais a presença de um especialista funciona menos como interlocução intelectual e mais como salvaguarda institucional. Fora desses casos, pesquisas que afetavam diretamente dimensões sociais eram conduzidas sem essa interlocução. A interdisciplinaridade proclamada operava de modo seletivo: certas competências eram incorporadas como constitutivas da pesquisa, outras eram convocadas apenas quando o risco de não fazê-lo se tornava visível.

Há nisso uma ironia que merece registro. Como documentei na seção sobre existências parciais, utilizei modelos de linguagem para me auxiliar na análise de dados e na elaboração de scripts, trabalho diretamente relacionado às competências dos cientistas da computação. A pergunta simétrica se impõe: quando pesquisadores de IA constroem sistemas que modelam comportamento social ou processam linguagem natural de comunidades específicas, recorrem eles a ferramentas conceituais das ciências sociais para formular seus problemas de pesquisa? A resposta, pelo que observei no Centro, é que raramente o fazem. A interdisciplinaridade tende a operar em mão única: cientistas sociais são incentivados a adquirir letramento técnico, mas o caminho inverso permanece excepcional.

Latour permite compreender por que essa assimetria se mantém. Quanto mais dura uma ciência se torna, mais social ela é, porque sua dureza depende de um número desproporcional de associações, instrumentos, laboratórios e aliados mobilizados para tornar seus fatos inquestionáveis \parencite{Latour2011CienciaEmAcao}. A dureza é efeito do comprimento da rede que sustenta esses objetos. Os pesquisadores que eu observava no C4AI estavam, justamente, no trabalho de construir essa dureza: mobilizando datasets, GPUs, publicações, métricas, financiamentos e parcerias para estabilizar afirmações que, sem essa rede, permaneceriam abertas à renegociação.

Stengers e Latour iluminam, portanto, operações complementares. Stengers diagnostica a imposição de um modelo de objetividade como critério universal. Latour revela a operação material que sustenta essa imposição: a extensão da rede como fonte da autoridade. Juntos, ajudam a compreender por que a interdisciplinaridade no C4AI operava de modo seletivo e por que o caminho do letramento técnico era percorrido em mão única.

Essa assimetria, contudo, não me levou a buscar o letramento técnico como gesto de contestação, mas como exigência etnográfica. Ao fazer pesquisa junto ao Centro, compreendi que para entender adequadamente um artigo científico da área seria necessário ir substancialmente além das generalizações frequentemente reproduzidas em debates externos ao campo. A linguagem técnica representa desafio epistemológico significativo não apenas por sua tecnicidade intrínseca, mas porque, devido à quantidade de elos, recursos e aliados mobilizados, ela se torna paradoxalmente mais social do que vínculos comunicativos convencionais \parencite{Latour2011CienciaEmAcao}. Cada artigo, cada equação, cada escolha metodológica carrega consigo extensa cadeia de referências, pressupostos compartilhados, debates anteriores, infraestruturas computacionais. Essas cadeias precisam ser minimamente conhecidas para que a compreensão seja efetiva.

Para enfrentar esses desafios, precisei investir tempo assistindo a palestras técnicas e realizando cursos sobre fundamentos da estatística e do aprendizado de máquina, entre os quais o curso \textit{Fundamentos de Probabilidade e Estatística para Ciência de Dados}, ministrado pelo Prof. Dr. Francisco A. Rodrigues (ICMC-USP) no âmbito do MBA em Ciência de Dados da USP.\footnote{O curso cobre três blocos temáticos encadeados: fundamentos de probabilidade (teorema de Bayes, esperança, momentos, teorema central do limite, lei dos grandes números), inferência estatística (estimação, intervalo de confiança, teste de hipóteses, inferência bayesiana) e modelagem com aprendizado supervisionado (regressão, ajuste de modelos, classificação). O material está disponível em: \url{https://cursosextensao.usp.br/course/view.php?id=4347}. Acesso em: mar. 2026. Esse percurso formativo é retomado no Capítulo~\ref{capitulo4}, onde os conceitos de aprendizado supervisionado e classificação binária são mobilizados para descrever o pipeline técnico do Spira.} Foi assim que compreendi, por exemplo, que o laboratório do cientista da computação é o computador, ideia aparentemente óbvia mas que possui implicações que merecem atenção.\footnote{A formulação é do Prof. Dr. Francisco A. Rodrigues, enunciada durante o curso \textit{Fundamentos de Probabilidade e Estatística para Ciência de Dados} (ICMC-USP), no contexto da introdução à simulação computacional: ``o computador é o laboratório para aprender os conceitos estatísticos e probabilísticos'' \parencite{Rodrigues2024ProbEstat}. A frase foi anotada no caderno de campo da pesquisadora durante a aula (Figura~\ref{fig:caderno-rodrigues}). O que parecia observação didática sobre o ensino de estatística revelou-se, ao longo da pesquisa de campo no C4AI, enunciado com consequências analíticas: se o computador é o laboratório, então o laboratório do cientista da computação é portátil, distribuível e dispensador de presença física contínua, o que reformula inteiramente a questão que abria esta pesquisa, a saber, onde está o laboratório de IA e onde estão os seus cientistas.}

\begin{figure}[htb]
  \centering
  \includegraphics[width=0.75\textwidth]{figuras/cap.1/caderno_quadro.jpg}
\caption[Anotação de caderno de campo de 28 de julho de 2025] {Anotação de caderno de campo de 28 de julho de 2025 durante o curso \textit{Fundamentos de Probabilidade e Estatística para Ciência de Dados} (ICMC-USP, Prof. Dr. Francisco A. Rodrigues). O quadro destaca a frase: ``Simulação: `o computador é o laboratório para aprender os conceitos estatísticos e probabilísticos'.'' As palavras ``computador'' e ``laboratório'' estão circuladas.}
  \label{fig:caderno-rodrigues}
\end{figure}

O nível de letramento técnico necessário varia conforme o objeto de estudo: investigações focadas nas práticas de pesquisa em inteligência artificial, como foi meu caso, exigem conhecimento aprofundado. Estudos centrados nos impactos sociais de sistemas já consolidados podem demandar níveis diferentes de familiarização. A necessidade de letramento técnico que experimentei encontra eco em debates dentro da própria ciência da computação. \textcite{Wallach2014} formula pergunta que espelha, desde dentro do campo, a observação que faço desde fora: poucos cientistas da computação desenvolveriam modelos para analisar dados astronômicos sem envolver astrônomos; por que, então, tantos métodos para analisar dados sociais são desenvolvidos sem o envolvimento de cientistas sociais? A resposta de Wallach é que, como seres humanos, temos intuições fortes sobre o mundo social, e essas intuições nos fazem crer que compreendemos fenômenos sociais sem precisar de formação específica. Mas intuição não é análise. Wallach não propõe defesa genérica da interdisciplinaridade. Seu argumento é estrutural: a colaboração precisa ser séria, não no modo contrate um cientista social para cada cem cientistas da computação, mas no modo crie equipes genuinamente interdisciplinares. A distinção importa: no primeiro caso, a presença do cientista social é decorativa; no segundo, ela é constitutiva da pesquisa.

Isso não significa que todo projeto de IA deva incluir cientistas sociais. Significa que a decisão sobre quando essa interlocução é relevante não deveria caber exclusivamente a quem não a pratica. É essa a questão que percorreu toda esta lição: quando pesquisadores de IA tomam o social como objeto, estão fazendo ciência social computacional sem reconhecê-la como tal, e portanto sem as salvaguardas que a tradição das ciências sociais desenvolveu para lidar com seres que não são indiferentes às perguntas que lhes são feitas. Reconhecer esse padrão é condição para construir alianças genuínas: alianças que permitam às ciências sociais participar da construção dos problemas de pesquisa, não apenas da gestão de seus efeitos colaterais.

O diagnóstico de Wallach converge, por via inesperada, com o aprofundamento da tecnologia que Vicentin formula a partir de Simondon: sem a articulação entre o fazer técnico e o conhecimento sócio-histórico, a tecnologia permanece aquém de sua ética imanente \parencite{Vicentin_2022_Esboco}. Quando os dispositivos de classificação e predição da IA são utilizados para reforçar estruturas de poder baseadas em classe, raça e gênero, a IA permanece em estágio de imaturidade no sentido simondoniano: limitada a reproduzir estruturas existentes em vez de operar seu potencial transformador.

O argumento desta lição pode agora ser reunido. A interpelação que recebi no C4AI revelou um padrão que opera em três planos. Pesquisadores de IA exigem letramento técnico de cientistas sociais que falam sobre IA, mas dispensam letramento sociológico quando eles próprios tratam do social. A colaboração interdisciplinar é seletiva: certas competências são incorporadas como constitutivas da pesquisa, outras convocadas apenas como salvaguarda institucional. E o letramento opera em mão única, do social para o técnico. Minha resposta foi dupla: adquirir o letramento necessário para compreender meus interlocutores nos seus próprios termos, e documentar a assimetria como achado etnográfico. As duas coisas não se contradizem. Aprender a linguagem técnica não foi capitulação diante da hierarquia disciplinar. Foi condição para habitar o campo com alguma autonomia e, ao mesmo tempo, descrevê-lo com precisão. A Lição 2 documentou a suspensão inicial que abriu o campo; esta documenta a aquisição progressiva de letramento que permitiu habitá-lo.

\subsection{Lição 4: Fazer ciência, fazer política: a pesquisa como mobilização de recursos}
\label{subsec:licao4}

De todas as lições, esta foi talvez a mais imediata, embora tenha sido a última a se formular como tal. Preciso ser franca sobre algo que me acompanhou durante boa parte do trabalho de campo: uma irritação surda, difícil de formular na época, com o fato de que eu não via a inteligência artificial sendo feita. Não como eu imaginava, ao menos. Eu havia me proposto a fazer uma etnografia de um centro de pesquisa em IA, e os leitores de uma tese como esta esperariam, com razão, saber como a inteligência artificial é produzida. O que eu encontrava, porém, era outra coisa: eventos de difusão científica, negociações de financiamento, apresentações a parceiros corporativos, disputas por infraestrutura computacional. A pesquisa em IA parecia acontecer em outro lugar, sempre atrás de uma porta que eu não conseguia abrir, dentro de computadores aos quais eu não tinha acesso, em reuniões técnicas cujo vocabulário eu não dominava. O que eu tinha diante de mim, porém, era, como defendi nas lições anteriores, a própria construção da inteligência artificial. Eu apenas não a identificava dessa maneira, porque seguia em busca da ciência por trás de toda aquela política.

Assistia a Fábio apresentar o Centro em eventos de difusão e anotava divulgação, não ciência. Ouvia Cláudio explicar como navegava entre os três potes de recursos e classificava aquilo como gestão, não pesquisa. Quando Cláudio descreveu em entrevista como, um mês antes de deadlines de conferências, tornava-se impossível conseguir GPU no Centro, reconheci ali o mesmo padrão: eu teria classificado aquilo como logística, não como decisão epistemológica. Na minha primeira visita ao Centro, parei diante da sala de TI sem compreender o que acontecia ali. Anos depois, percebi que aquele desconforto era sintomático de uma separação que eu carregava e que o campo insistia em desfazer. Eu procurava a inteligência artificial dentro dos computadores, nos códigos, nos modelos, e o que encontrava eram pessoas negociando, traduzindo interesses, reconfigurando alianças. A angústia não era apenas intelectual; era profissional. E me perguntava com frequência: como apresentar uma etnografia de um centro de inteligência artificial sem propriamente descrever como a IA é feita?

Essa frustração, contudo, era metodologicamente produtiva. Ela sinalizava o ponto exato onde minhas categorias analíticas falhavam, onde o arranjo que eu trazia precisava ser reconfigurado. Law argumenta que o hinterland de qualquer método se ramifica indefinidamente para além de laboratórios e protocolos, alcançando conhecimento tácito, software, habilidades linguísticas, capacidades de gestão, sistemas de transporte e comunicação, escalas salariais, fluxos de financiamento, prioridades de agências de fomento e agendas abertamente políticas e econômicas \parencite[p. 41-42]{Law2004After}. O que eu observava no C4AI (as sessões de Design Thinking que articulavam critérios da USP, IBM e Brasil; a venda da divisão de saúde da IBM que redirecionava linhas de pesquisa; as disputas por GPU que se tornavam decisões epistemológicas; Fábio apresentando em Bluetalks como parte do fazer ciência) era precisamente esse hinterland ramificado que métodos convencionais relegam à invisibilidade.

Law descreve method assemblage como produção de um feixe de relações ramificadas que gera presença, ausência manifesta e Outridade \parencite[p. 84]{Law2004After}. No meu caso: a presença esperada era a ciência sendo feita (códigos, modelos, experimentos); a ausência manifesta eram as negociações que eu sabia serem necessárias mas classificava como entorno; e a Outridade era precisamente o trabalho de mobilização de recursos que precisava desaparecer para que a ciência aparecesse como autônoma. Minha frustração era sintoma de que minhas categorias prévias operavam exatamente essa divisão entre presença, ausência e Outridade. O que o campo me obrigou a fazer foi reconfigurar essas fronteiras: reconhecer que o que eu havia relegado à Outridade era constitutivo da presença.

A pergunta onde está a ciência que esses recursos sustentam? pressupunha a separação entre ciência e política que o próprio campo tratava de desmentir. Mobilizar recursos, negociar parcerias, justificar investimentos perante agências de fomento e corporações transnacionais é a pesquisa ela mesma. A pergunta não encontrou resposta. Ela perdeu sentido.

A dificuldade em reconhecer aquelas atividades como fazer IA era expressão concreta da alienação técnica discutida anteriormente: a separação entre o que seria estritamente técnico e o que seria entorno da pesquisa reproduz a cisão que impede o amadurecimento da tecnologia. O que o campo me obrigou a compreender, por vias etnográficas, era que as negociações de financiamento, as disputas por GPU, as apresentações a parceiros corporativos eram parte constitutiva da individuação técnica da IA.

Essa inseparabilidade entre ciência e política era o que o campo tornava visível quando eu conseguia suspender a distinção prévia. Os próprios projetos de pesquisa do Centro foram definidos, como Fábio relata, por sessões de Design Thinking em que três critérios precisavam se alinhar simultaneamente: a USP teria liderança científica no tema, a IBM teria interesse na temática, e o Brasil teria um diferencial internacional. Não existia um momento científico de escolha de tema seguido de um momento político de negociação de recursos. A decisão era uma só, e nela se fundiam avaliação epistêmica, estratégia corporativa e posicionamento geopolítico.\footnote{Os relatórios internos do C4AI documentam esse processo: o grupo de Design Thinking organizou quatro reuniões com mais de 80 pesquisadores, convergindo em cinco desafios. O critério exigia que os resultados fossem estratégicos para o Brasil, relevantes para USP e IBM e com alto potencial de inovação. A partir de 2022, quando a IBM vendeu Watson Health, o Centro decidiu não enfatizar esforços relacionados a saúde; quando o grupo AgriBio conquistou financiamento independente (INCT), tornou-se autossuficiente em relação aos recursos do C4AI. Em ambos os casos, a reconfiguração não separava decisão científica de decisão política.}

Quando a IBM vendeu sua divisão de saúde, o Centro parou de investir diretamente em pesquisa de IA em saúde. Quando a IBM saiu de agronegócio, o mesmo ocorreu. As reorientações de pesquisa não eram imposições externas sobre uma agenda científica autônoma; eram reconfigurações da rede em que agenda científica e interesse corporativo constituíam, desde a origem, um único arranjo.

Do mesmo modo, quando Fábio apresentava resultados em eventos como as Bluetalks, estava recrutando aliados, fortalecendo a rede, traduzindo interesses de públicos heterogêneos para manter a cadeia de associações que tornava a pesquisa possível. Latour descreve precisamente esse trabalho como constitutivo da prática tecnocientífica: o que se perde ao separá-lo da ciência propriamente dita é justamente o que o torna ciência, a capacidade de alistar e manter associados elementos heterogêneos \parencite{Latour2011CienciaEmAcao}. A junção dos termos técnica e ciência em tecnociência é exigência empírica, e o campo que eu observava tornava essa exigência particularmente visível.

A agência, nessas reconfigurações, é distribuída pela rede: emerge das associações entre pesquisadores, instituições, infraestruturas, contratos, modelos computacionais e políticas de fomento. Nenhum desses elementos age sozinho ou determina sozinho o curso da pesquisa. Quando o C4AI delega a definição de temas de pesquisa a sessões de Design Thinking que articulam simultaneamente critérios da USP, da IBM e do Brasil, a agência dessa decisão reside na associação entre os participantes, e cada reconfiguração da rede redistribui a agência por novas cadeias de translação.

O cenário contemporâneo de investimentos massivos em inteligência artificial (com a concentração de capacidade computacional em poucas corporações, a corrida geopolítica por soberania tecnológica e a crescente dependência de infraestruturas de alto custo) torna essas dinâmicas particularmente visíveis. Resta saber se a pesquisa em IA constitui um caso em que a inseparabilidade entre ciência e política se intensifica a ponto de exigir novos instrumentos analíticos, ou se apenas torna mais difícil ignorar o que Latour já descrevia em outros campos, quando os investimentos eram menores e as redes, menos extensas.

\section{``Existências parciais'': quase-máquinas, híbridos e a composição com IA}
\label{subsec:existencias-parciais}

\epigrafe{Direct access to anything is not in the power of \enquote{humans} souls and machinery. Modems and faxes add just a few more artifacts in the midst of all of those that already compose \enquote{natural} interactions. [...] We all have only partial existence.}
{Bruno Latour e Richard Powers, \textit{Two Writers Face One Turing Test: A Dialogue in Honor of HAL}, 1998.}

As polarizações e multiplicidades descritas nas páginas anteriores exigem ferramental teórico capaz de compreender como pessoas e sistemas de IA se compõem sem reduzir essa composição a categorias prévias (usuário/ferramenta, humano/máquina, autêntico/artificial). Os conceitos de existência parcial, conexão parcial e agência distribuída oferecem esse ferramental.

\subsection*{Do princípio de simetria às existências parciais}

Proponho aqui uma leitura da formulação ``existências parciais'' não como ajuste vocabular, mas como deslocamento teórico interno à própria obra de Latour. O Latour de \textit{Ciência em Ação} \parencite*{Latour2011CienciaEmAcao} e de \textit{Jamais Fomos Modernos} \parencite*{Latour1994Jamais} opera com o princípio de simetria generalizada: tratar humanos e não humanos nos mesmos termos. Esse é passo decisivo porque desmonta a pressuposição de que apenas um dos lados da relação merece descrição. Mas o vocabulário da simetria ainda pressupõe dois polos a serem simetrizados, e a operação contesta a hierarquia sem dissolver a própria partição.

O diagnóstico dos híbridos é persistente na obra de Latour: dos quase-objetos de \textit{Jamais Fomos Modernos} ao projeto Aramis, dos híbridos de \textit{A Esperança de Pandora} \parencite*{Latour2017Esperanca} à crise da crítica em \textit{On the Modern Cult of the Factish Gods} \parencite*{latour2010factish}. As misturas não são categoria intermediária a ser superada, mas a condição mesma do mundo que habitamos. No diálogo com Powers de 1998, a formulação se radicaliza: ``todos temos apenas existência parcial'' \parencite{Latour2008Powers}. Não são mais dois polos simetrizados, nem entidades definidas pelo que não são; são existências que nunca foram completas para começar. E na \textit{Investigação sobre os Modos de Existência} \parencite*{latour2019investigacao}, o próprio conceito de ator-rede é reconhecido como insuficiente, apenas um modo entre outros.

Esse movimento que identifico na obra de Latour não surgiu de revisão bibliográfica, mas como consequência de meu próprio percurso em campo. Foi acompanhando pesquisadores de IA no C4AI, observando como seus algoritmos, datasets e modelos eram simultaneamente produtos de escolhas humanas e entidades com comportamentos que surpreendiam seus próprios criadores, que a insuficiência do vocabulário simétrico surgiu. A inteligência artificial generativa exemplifica de forma contundente o que significa existência parcial: é composta por milhões de decisões humanas sedimentadas e, no entanto, produz resultados que nenhum de seus criadores previu ou controla inteiramente; opera sobre conteúdo simbólico e, no entanto, não compreende no sentido humano do termo; é radicalmente dependente de infraestruturas institucionais e, no entanto, age como interlocutor com certa autonomia funcional. Não é humano nem não humano: é existência parcial de modo brutal, mais visível do que qualquer vieira ou micróbio porque opera justamente no terreno que a tradição ocidental reservou como marca distintiva do humano: linguagem, pensamento, escrita.


\subsection*{Quase-máquinas e agência distribuída}

Sistemas de IA generativa como ChatGPT e Claude são produto de escolhas humanas em cada etapa de sua construção, da curadoria dos dados de treinamento à definição de arquiteturas de redes neurais, do ajuste fino com avaliadores humanos às decisões sobre o que o modelo deve ou não fazer.\footnote{O processo de construção de um modelo de linguagem envolve camadas sucessivas de decisões humanas que permanecem opacas ao usuário final: a seleção e filtragem dos corpora de treinamento, a definição da arquitetura e dos hiperparâmetros da rede neural, o ajuste fino por meio de feedback humano (reinforcement learning from human feedback, RLHF), e as políticas de segurança e alinhamento. Cada uma dessas etapas incorpora pressupostos culturais, vieses e escolhas normativas que condicionam os resultados. Sobre os desafios éticos e políticos dessas escolhas, ver \textcite{Buolamwini2023Unmasking}, \textcite{Noble2018Algorithms}, \textcite{Benjamin2019RaceAfter} e \textcite{Christian2020Alignment}.} São, portanto, objetos técnicos saturados de significação. A IA generativa, diferentemente de ferramentas anteriores como processadores de texto ou buscadores acadêmicos, efetivamente opera sobre o conteúdo: sugere reformulações, identifica inconsistências, propõe estruturas argumentativas. Essa mediação é qualitativamente distinta, e é precisamente sua especificidade que se perde quando reduzida ao binário: feito por humano ou feito por não humano.

A partição entre o que foi feito por humanos e o que foi gerado por IA pressupõe a existência de entidades puramente humanas de um lado e puramente maquínicas do outro, separadas por uma fronteira estável e visível. Mas essa partição é exatamente o que precisa ser explicado, não aquilo de onde se parte. Em um diálogo que dramatiza essa questão, Latour propõe a Powers que abandonem o vocabulário da ``máquina'' em favor do de ``quase-máquina'': entidades abertas, conectadas, que jamais cortaram o cordão umbilical com as instituições, os corpos e os coletivos que as sustentam. A questão pertinente não é ``a quem imaginamos estar falando quando falamos com nossas máquinas?'', mas sim: ``através de quais máquinas falamos? Quais máquinas nos permitem falar, para início de conversa?''. Porque mesmo quando estamos sozinhos, na carne, falamos através de toda sorte de máquinas: fonação, cordas vocais, língua inglesa, convenções estilísticas \parencite{Latour2008Powers}.

Essa discussão ajuda a compreender por que o debate sobre inteligência artificial permanece preso a dualismos que, embora teoricamente superados, são consistentemente reafirmados na prática. A pergunta ``podem as máquinas pensar?'' já foi reformulada por Alan Turing como problema operacional, o jogo da imitação, precisamente porque a formulação original carregava pressupostos que ele considerava improdutivos.\footnote{Alan Turing utilizava ironia e sátira como ferramentas retóricas para desafiar o antropocentrismo prevalecente. Seu jogo da imitação emprega metáforas metodologicamente: máquina-como-mulher, máquina-como-humano operam como dispositivos conceituais para desnaturalizar pressupostos sobre inteligência. A imitação para Turing funcionava mais como um princípio matemático do que uma identidade ontológica: máquinas poderiam imitar comportamento inteligente permanecendo ontologicamente distintas de humanos \parencite[p. 5]{Goncalves2024BeautifulThought}. \textcite{Quattrociocchi2025Epistemological} sistematizam essa heterogeneidade através de sete fault lines epistemológicas estruturais.} Latour, por sua vez, questiona o próprio teste de Turing em seus fundamentos: na prática, o teste confronta uma máquina de carne e osso contra outra máquina de carne e osso. Coisas e pessoas estão demasiado entrelaçadas para serem particionadas antes que o teste comece \parencite{Latour2008Powers}.

A objeção de Lady Lovelace, segundo a qual ``a máquina não tem pretensão de originar coisa alguma'', já fora contestada pelo próprio Turing, que afirmava que máquinas o surpreendiam ``com grande frequência'' \parencite{Turing1950CMI}. Latour radicaliza a implicação: ninguém domina as quase-máquinas, porque elas são caixas-pretas que ninguém atravessa com o olhar, embora sejam de nossa própria fabricação. O mais inteligente dos programadores perde-se na linha quatro mil de seu programa, e quando uma dúzia de programadores trabalha junta, o programa começa a ter vida própria \parencite{Latour2008Powers}. A crença no domínio, seja de Deus sobre os humanos, dos humanos sobre as máquinas, ou das máquinas sobre os humanos, é precisamente o mito que as quase-máquinas confrontam.

Não há, portanto, entidades que sejam plenamente humanas de um lado e plenamente maquínicas do outro, mas existências parciais que se compõem mutuamente. A noção de conexão parcial, formulada por Strathern a partir de sua leitura crítica da teoria ator-rede, oferece aqui uma precisão decisiva: as entidades não são nem completamente conectadas (o que as tornaria a mesma coisa) nem completamente separadas (o que tornaria impossível qualquer relação entre elas), mas parcialmente conectadas; análogas sem serem homólogas \parencite{strathern2005partial}. A pesquisadora que escreve com assistência de IA não é humano mais máquina, como se duas entidades discretas se somassem; é uma composição parcial, na qual nem a pesquisadora existe independentemente das mediações técnicas que a atravessam, nem o modelo de linguagem existe independentemente das decisões humanas que o constituíram. Latour chega ao mesmo ponto por outra via: quase-máquinas não existem separadas da humanidade; são a humanidade em outro estado, assim como água, vapor e gelo são estados diferentes da mesma substância \parencite{Latour2008Powers}.

A existência parcial descreve o que as entidades são: nunca completas, sempre dependentes das associações que as constituem. Falta ainda descrever como essas existências parciais agem juntas, por meio de qual mecanismo a composição entre elas produz efeitos no mundo. É aqui que a noção latouriana de agência distribuída se torna indispensável.\footnote{Latour não emprega sistematicamente a expressão ``agência distribuída'' como rótulo teórico unificado, mas desenvolve o conceito ao longo de vários textos sob formulações como ``distribuição da atividade'', ``compartilhamento da responsabilidade pela ação'' e, sobretudo, \textit{faire-faire} (fazer-fazer). As elaborações mais densas encontram-se no Capítulo 6 de \textit{A esperança de Pandora} \parencite{Latour2017Esperanca}, no artigo sobre o dispositivo de fechamento automático de portas \parencite{johnson1988mixing}, e em \textit{On the Modern Cult of the Factish Gods} \parencite{latour2010factish}.} Para Latour, a agência nunca é propriedade isolada de um actante. É sempre resultado de uma composição de forças em uma rede, onde a ação é um faire-faire distribuído entre diversos mediadores.

O argumento é construído através de exemplos que dramatizam a distribuição. No Capítulo 6 de \textit{A Esperança de Pandora}, Latour utiliza o debate sobre controle de armas para demonstrar que nem a pessoa sozinha nem a arma sozinha matam: a responsabilidade pela ação deve ser compartilhada entre os vários actantes envolvidos. Quando um cidadão empunha uma arma, ele se torna um cidadão-arma, ator híbrido cujo objetivo de ação foi transladado pela associação. Ninguém permanece o mesmo: o cidadão que era pacífico tornou-se ameaçador. A arma que era inerte tornou-se letal. O que age é a associação entre eles \parencite{Latour2017Esperanca}. No artigo sobre o dispositivo de fechamento automático de portas, Latour mostra como moralidade e disciplina são delegadas de humanos a não humanos: conhecimento, moralidade, força e sociabilidade não são propriedades de humanos, mas de humanos acompanhados por seu séquito de delegados técnicos \parencite{johnson1988mixing}. O conceito de fatiche, fusão entre fato e fetiche, radicaliza a implicação: ``aquele que age não é mestre do que faz'', pois outros actantes ``passam à ação'' e o ultrapassam \parencite{latour2010factish}.

Essas formulações possuem implicação direta para a composição com inteligência artificial documentada nesta tese. O vocabulário de Latour da agência distribuída descreve o mecanismo dessa composição; o vocabulário Haraway da \textit{simpoiese} nomeia sua condição ontológica: nem pesquisadora nem modelo preexistem, como entidades completas, ao \textit{fazer-com} que os constitui mutuamente ao longo do processo. A composição altera o que estou tentando fazer, os caminhos argumentativos que percorro, as formulações que se tornam possíveis e aquelas que se fecham. O objetivo da ação é transladado: o que começa como tentativa de reformular um parágrafo pode desembocar em reorganização conceitual de uma seção inteira. A translação é o modo como a mediação opera, e é também o modo como a simpoiese se distingue do simples uso de ferramenta: o fazer-com transforma os termos da relação, não apenas seus produtos.

O paralelo com o fermento de Pasteur é igualmente produtivo. No Capítulo 4 de \textit{A Esperança de Pandora}, a análise do trabalho de Pasteur revela que ele trabalha para que o fermento aja sozinho, e se o experimento for bem-sucedido, o fermento torna-se ator com autonomia funcional, mas essa autonomia é efeito da rede de mediações laboratoriais \parencite{Latour2017Esperanca}. Analogamente, eu trabalho (formulando perguntas, fornecendo material etnográfico, avaliando criticamente os resultados) para que o modelo aja, produzindo tensionamentos, reorganizações e formulações que eu então aceito, rejeito ou modifico. Se a composição funciona, o modelo adquire certa autonomia funcional: sugere caminhos que eu não teria percorrido, identifica inconsistências que eu não teria notado, propõe articulações que alteram o argumento. Mas essa autonomia, como a do fermento de Pasteur, é efeito da rede de mediações que a tornou possível. A agência é distribuída pela cadeia de traduções entre pesquisadora, modelo, texto e argumento.

Latour menciona explicitamente a compatibilidade de sua teoria com a obra de Edwin Hutchins sobre cognição distribuída, que descreve como o trabalho do pensamento é externalizado em formas e instrumentos \parencite{latour2010factish}.\footnote{A obra de referência é \textit{Cognition in the Wild}, na qual Hutchins demonstra, a partir de etnografia da navegação marítima, que processos cognitivos complexos não residem dentro de indivíduos, mas se distribuem entre pessoas, artefatos e ambientes estruturados.} O pensamento que produz o argumento acadêmico distribui-se entre a pesquisadora, o modelo, os diagramas, as notas de campo, os textos teóricos consultados e as inscrições visuais que emergem dessa cadeia. A externalização do pensamento em instrumentos é redistribuição da agência cognitiva pela rede.

A agência distribuída oferece, portanto, o mecanismo que a existência parcial postula mas não desenvolve: se todos temos apenas existência parcial, agimos sempre através de mediadores que transladam nossos objetivos, redistribuem a responsabilidade pela ação e produzem resultados que nenhum actante isolado teria sido capaz de alcançar. É esse faire-faire que descreve com precisão o que acontece quando uma pesquisadora compõe com um modelo de linguagem para produzir uma tese: redistribuição da agência por uma cadeia de translações em que cada mediador transforma o que o anterior lhe passou.

\subsection*{Multiplicidades ontológicas e a especificidade de Claude}

A perspectiva de Law sobre multiplicidade ontológica é particularmente produtiva para pensar sistemas de IA generativa. Considere o caso de Claude. ``Claude'' é \textit{enacted} como objetos distintos através de práticas irredutíveis entre si.

\textbf{Claude-em-treinamento} existe nas instalações da Anthropic como dataset massivo, arquitetura neural (transformer com bilhões de parâmetros), sessões de RLHF com avaliadores humanos que julgam pares de respostas, e políticas de Constitutional AI que definem valores incorporados ao modelo. Esse Claude é objeto de decisões sobre o que incluir ou excluir dos dados de treinamento, como ponderar feedbacks contraditórios, quais capacidades desenvolver ou suprimir. As escolhas de curadoria, arquitetura e alinhamento constituem este Claude, que não existe independentemente das práticas que o produzem.

\textbf{Claude-em-deployment} existe como serviço distribuído em datacenters da AWS, sistema de rate limiting, protocolos de moderação de conteúdo, versões simultâneas (Haiku, Sonnet, Opus) com capacidades diferentes e custos escalonados. Esse Claude é objeto de estratégias comerciais, políticas de privacidade, decisões sobre disponibilidade regional. A infraestrutura que o distribui e as políticas que governam seu acesso o constituem tanto quanto a arquitetura neural.\footnote{Durante todo o percurso de colaboração documentado nesta tese, desativei a opção de compartilhamento de dados para treinamento de modelos nas configurações de privacidade da plataforma Claude.ai. Essa escolha incide diretamente sobre o que constitui tanto o Claude-em-deployment (as condições de acesso que configuro) quanto o Claude-em-treinamento (os dados que não entram no corpus de versões futuras), ilustrando como políticas de privacidade não são externas à relação pesquisadora-modelo, mas co-constitutivas dela.}

\textbf{Claude-em-auditoria} existe nos laboratórios de red teaming onde pesquisadores tentam induzi-lo a gerar conteúdo perigoso, em datasets de benchmark onde seu desempenho é medido contra outros modelos, em análises de viés que examinam suas respostas por demografia simulada. Esse Claude é objeto de preocupações regulatórias, debates sobre safety e alignment, disputas sobre transparência.

\textbf{Claude-em-uso} existe diferentemente em cada conversa e contexto de aplicação. Quando discute teoria ator-rede com uma doutoranda brasileira, Claude-em-uso mobiliza corpus acadêmico, identifica referências bibliográficas, tensiona argumentos teóricos. Não é o mesmo Claude que responde perguntas médicas a um clínico nem o que auxilia programação para um desenvolvedor. Cada interação enacta affordances distintas, vocabulários distintos, limites distintos. O Claude que colabora na escrita desta tese não é nem Claude-universal-aplicável-a-qualquer-tarefa nem Claude-especialista-em-etnografia, mas Claude-composto-com-esta-pesquisadora-específica-neste-projeto-específico.

Esses não são aspectos de um Claude singular. São Claudes múltiplos que coexistem sem convergir em unidade. O Claude-em-treinamento que os engenheiros da Anthropic descrevem não é o mesmo objeto que o Claude-em-uso experimenta nesta colaboração de tese. Não há um Claude real por trás dessas variações; a multiplicidade é constitutiva.

Talvez seja precisamente por isso que sistemas de IA são tão difíceis de estudar: não apenas porque sejam tecnicamente opacos (o problema clássico da caixa-preta), mas porque são ontologicamente múltiplos. Os debates contemporâneos sobre explicabilidade em IA geralmente pressupõem que existe um objeto singular a ser explicado. Mas se Law estiver correto, o problema não é apenas de acesso epistemológico (não conseguimos ver dentro), mas de multiplicidade ontológica: não há um único objeto lá dentro esperando para ser revelado. A opacidade é condição ontológica de objetos que existem como multiplicidade em \textit{continued enactment}.

Três implicações dessa multiplicidade para a colaboração documentada nesta tese merecem atenção. A qualidade da composição pessoa-IA não depende apenas de ter acesso ao modelo, mas da constelação específica de competências, recursos e posicionamentos que cada actante traz para a interação. Uma pesquisadora com experiência acadêmica, capital cultural acumulado, familiaridade com convenções disciplinares e fluência em inglês compõe com Claude-em-uso de modo radicalmente distinto de estudante de graduação sem experiência de pesquisa, com vocabulário limitado, usando modelo em português com capacidades reduzidas.\footnote{Parte desse capital cultural é composta por formações que não fazem parte do percurso convencional de uma doutoranda em ciências sociais. Um MBA em Ciência de Dados cursado na USP e, mais diretamente, uma especialização em Ciência de Dados pela Unicamp, com disciplina de Processamento de Sinais Bidimensionais com Python, introduziram conceitos como Transformada de Fourier, espectrogramas e representações espectrais que o Capítulo~\ref{capitulo4} mobiliza para descrever o pipeline técnico do Spira. Esse capital técnico não estava previsto no projeto de pesquisa; foi acumulado paralelamente ao doutorado por motivações independentes. O ponto de contato entre os dois percursos só se revelou durante a escrita do Capítulo~\ref{capitulo4} e tornou-se uma possibilidade de uma descrição etnográfica da cadeia de translações do Spira de modo especial. A composição com Claude nessa seção foi, portanto, diferente: a pesquisadora não dependia do modelo para traduzir os conceitos técnicos, mas para articulá-los ao vocabulário latouriano. Nos termos de \textcite{Vicentin_2022_Esboco}, o letramento técnico operou não como substituição da perspectiva das ciências sociais, mas como condição para que a análise crítica não se rendesse à opacidade que Simondon identifica como alienação.} Não porque usem a mesma ferramenta diferentemente, mas porque compõem existências parciais radicalmente distintas.

Disso decorre que diferenças de classe social, gênero, raça, profissão, formação educacional, acesso a mentorias, domínio de línguas e familiaridade com códigos acadêmicos condicionam o tipo de composição possível e, consequentemente, o resultado gerado. A composição pessoa-IA não é neutra ou igualitária: amplifica desigualdades estruturais preexistentes, transformando diferenças de capital cultural em diferenças de capacidade de extrair valor da ferramenta.

Há ainda uma terceira implicação, de outra ordem: os próprios modos de inteligência em composição são ontologicamente heterogêneos. Quattrociocchi, Capraro e Perc mapeiam sistematicamente sete fault lines epistemológicas que separam julgamento humano de julgamento em LLMs \parencite{Quattrociocchi2025Epistemological}. Claude opera por plausibilidade linguística derivada de padrões estatísticos em corpus massivo, o que os autores descrevem formalmente como caminhada estocástica em grafo de alta dimensão de transições linguísticas aprendidas, um processo ergódico sob restrições estatísticas, não um procedimento de convergência epistêmica \parencite[p. 3]{Quattrociocchi2025Epistemological}. A pesquisadora, por sua vez, opera por significação situada, memória incorporada e reflexividade biográfica inserida em loop epistêmico: crenças colidem com contraexemplos, conclusões permanecem revisáveis diante de nova informação, o mundo empurra de volta \parencite[p. 4]{Quattrociocchi2025Epistemological}.

Quattrociocchi et al. denominam Epistemia a condição estrutural em que plausibilidade linguística substitui avaliação epistêmica, produzindo a experiência de posse de uma resposta sem ter atravessado o processo de formação de crença justificada \parencite[p. 8]{Quattrociocchi2025Epistemological}. A colaboração entre pesquisadora e Claude nesta tese situa-se precisamente nessa tensão: a mediação é constitutiva, mas exige que a pesquisadora atue como guardiã do loop epistêmico que o sistema não possui. Law observa que method assemblage é sempre muito mais do que suas descrições formais sugerem, porque envolve um processo contínuo de crafting e produção de fronteiras necessárias entre presença, ausência manifesta e Outridade \parencite[p. 144]{Law2004After}. Nesta pesquisa, manter presentes as decisões teóricas, conceituais e de curadoria bibliográfica, em vez de delegá-las ao modelo, constitui reconhecimento de que essas operações exigem justificação epistêmica que a plausibilidade linguística não fornece.

A retórica da democratização que frequentemente acompanha a defesa da IA generativa, a ideia de que ela nivela o campo de jogo, colapsa quando se reconhece essa multiplicidade. Não existe o Claude que todos acessam igualmente. Existem múltiplos Claudes enacted por composições situadas cujos efeitos não podem ser julgados abstratamente. A pergunta que se impõe é outra: quais composições específicas, entre quais actantes, produzindo quais Claudes múltiplos, gerando quais efeitos distribuídos por quais redes?

\subsection*{Alienação técnica e letramento interdisciplinar}

Diego Vicentin desdobra a alienação técnica que Simondon diagnostica em duas facetas que o trabalho de campo no C4AI traz elementos para pensar. De um lado, a alienação dos usuários e consumidores, que desconhecem o funcionamento dos sistemas algorítmicos que tomam decisões sobre suas vidas, o efeito caixa-preta que se impõe por razões que vão do segredo industrial à complexidade que certos sistemas adquirem. De outro, a alienação dos próprios técnicos e engenheiros, que ignoram as raízes epistêmicas, sociais e políticas dos sistemas com os quais trabalham, bem como a conexão entre técnica e estética que permitiria reconhecer a significação do sistema técnico para além de sua utilidade imediata \parencite{Vicentin_2022_Esboco}. Quando o letramento técnico é proposto como pré-condição para o diálogo interdisciplinar, mas não se busca competência em ciências sociais ao trabalhar com fenômenos sociais, opera-se dentro dessa segunda forma de alienação.\footnote{Segundo Vicentin, Simondon utilizou a noção de \textit{technologie approfondie} (tecnologia aprofundada) em dois textos separados por trinta anos. Para Vicentin, a noção designa uma tecnologia capaz de conectar o trabalho prático e o conhecimento especializado com a história do pensamento e a consciência de uma sociedade. Ver \textcite{Vicentin_2022_Esboco} para a articulação entre esse conceito e os estudos críticos contemporâneos sobre IA.}

Já em 1958, a alternativa proposta por Simondon consistia em reconhecer-se como organizador permanente de uma sociedade de objetos técnicos, atuando como tradutor e mediador de informações, intervindo nas zonas de incerteza e atribuindo sentido a essas interações \parencite{Simondon2017OMET}. Latour reconheceu explicitamente essa dívida: ao propor que nunca fomos modernos porque jamais conseguimos separar de fato natureza e sociedade, humano e não humano, operava sobre o terreno que Simondon havia preparado ao mostrar que a alienação técnica consiste precisamente nessa separação \parencite{Latour1994Jamais}.

\subsection{Por que fazer parentesco com este problema?}

Reconhecer a parcialidade das existências e a distribuição da agência, porém, não equivale a naturalizar essa composição específica. Diferentemente de tecnologias que já se tornaram componentes invisibilizados da prática de pesquisa (processar texto no Word, buscar referências na web, gerenciar bibliografia), a colaboração com IA generativa permanece uma escolha metodológica deliberada de fazer parentesco com um companheiro problemático. \textbf{Mas por que fazer esse parentesco, tendo em vista todos esses problemas? Por que não evitar essa mediação?}

Podemos ficar com o problema e pensar a partir de três planos. Primeiro, porque a IA generativa já está reconfigurando práticas de escrita acadêmica, e a questão é como ela estará presente nas ciências sociais e sob quais condições de reflexividade. Fazer parentesco deliberado com esse parente estranho consiste, além de um experimento antropológico, em um gesto político-metodológico que recusa tanto a adesão acrítica quanto a recusa moralista: trata-se de habitar o problema para conhecê-lo por dentro. Como escreve Haraway:

\begin{citacaoabnt}
Ficar com o problema não requer uma relação com tempos chamados de futuro. Na verdade, ficar com o problema requer aprender a estar verdadeiramente presente, não como um pivô que desaparece entre passados terríveis ou edênicos e futuros apocalípticos ou salvíficos, mas como criaturas mortais entrelaçadas em miríades de configurações inacabadas de lugares, tempos, matérias, significados. \parencite[p. 1]{Haraway2016Staying}
\end{citacaoabnt}

Segundo, porque esta tese estuda etnograficamente a construção de IA no C4AI. Fazer parentesco com Claude constitui forma de intensificar, ou até mesmo provocar deliberadamente, a experiência de sujeitar-se a mediações técnicas extremas. Submeto-me ao mesmo universo de questões que atravessa meus interlocutores: o que significa compor com sistemas que operam por lógicas não-transparentes? Como manter responsabilidade quando a agência está distribuída em redes sociotécnicas opacas? Onde se localiza a autoria quando humanos e não-humanos co-produzem resultados? Mas não apenas isso: o que essa experimentação faz comigo, enquanto pesquisadora? Como ela transforma minha relação com o texto, com as referências teóricas, com a elaboração conceitual? Quais ansiedades, inseguranças e dilemas essa composição produz no processo concreto de escrita? E, mais amplamente: o que essa mediação significa para a produção de conhecimento nas ciências sociais? Que tipo de saber é produzido quando a plausibilidade linguística se articula (sempre assimetricamente, sempre sob vigilância) com justificação epistêmica?

A Figura~\ref{fig:composicao-claude-1}, a Figura~\ref{fig:composicao-claude-2} e a Figura~\ref{fig:composicao-claude-3} reproduzem uma sequência de trocas com Claude durante a escrita deste capítulo, na revisão da seção em que se discute, a partir de Latour, Mol, Law e Barad, a minha participação na rede que descrevo. A figura é uma amostra do que essa composição faz na prática: uma cadeia de tensionamentos, recuos e reformulações em que eu e Claude redistribuímos a agência por sucessivas e recursivas rodadas de escrita. Diferentemente dos diagramas dos Capítulos~\ref{capitulo3} e~\ref{capitulo4}, que circulam como inscrições interativas acessíveis por link de compartilhamento permanente no MermaidChart e que, ao serem navegadas pelo leitor, redistribuem a agência interpretativa em outra direção, esta cena chega ao leitor como captura de tela: a conversa foi inscrita no momento em que aconteceu e congelada no formato em que foi registrada, sem possibilidade de retorno ao estado fluido. A escolha pela imagem capturada, em vez de transcrição editada para o corpo do texto ou link para o histórico online da conversa, responde à intenção de inserir essa cena na tese da forma mais fidedigna possível: com os turnos identificados como aparecem na interface da \texttt{claude.ai}, com o registro temporal preservado, com a textura visual do diálogo intacta. O que se ganha em fidelidade ao registro empírico se perde em legibilidade tipográfica, e essa perda é parte do gesto metodológico aqui assumido. Para uma melhor visualização da conversa, recomenda-se aumentar o zoom da tela na versão digital da tese.

\begin{figure}[!htb]
  \centering
  \includegraphics[width=0.7\textwidth]{figuras/cap.1/1000015590.pdf}
  \caption[Composição com Claude durante a escrita do Capítulo~\ref{capitulo1} (parte 1 de 3)]{Composição com Claude durante a escrita do Capítulo~\ref{capitulo1} (parte 1 de 3): pedido de fechamento de uma nota de rodapé sobre arquivos digitais, modelos de linguagem e a indeterminação dos \textit{prompts}. Captura de tela registrada em ... na interface \texttt{claude.ai}. A composição mobiliza a \textit{skill} \texttt{juliane-dissertacao} (em \texttt{/mnt/skills/user/juliane-dissertacao/}), que estabiliza para Claude as regras de estilo da tese (proibição de travessões, recusa de adjetivos avaliativos, recusa da fórmula \enquote{não é X, é Y}, uso do ambiente \texttt{citacaoabnt}, padronização de aspas com \texttt{\textbackslash enquote\{\}}), o quadro teórico operativo (Latour, Mol \& Law, Akrich, Callon) e o estado do Capítulo~\ref{capitulo4}. Como \textit{hinterland} da conversa operam os capítulos em curso da tese, consultados pelo modelo via ferramentas de leitura de arquivos: \texttt{ex\_cap0.tex}, \texttt{ex\_cap1\_-\_2026-04-12T002133\_890.tex}, \texttt{ex\_cap2\_\_31\_.tex} e \texttt{ex\_cap4\_\_60\_.tex}, além dos arquivos de referência da \textit{skill} (\texttt{quadro\_teorico.md}, \texttt{contexto\_cap4.md}, \texttt{bibtex\_ativas.md}). A cena exibe Claude propondo um fechamento, oferecendo razões estilísticas para os cortes feitos no texto original, e classificando algumas divagações como \enquote{produtivas} ou \enquote{improdutivas}: a régua normativa que a turno seguinte da pesquisadora vai contestar (Figura~\ref{fig:composicao-claude-2}).}
  \label{fig:composicao-claude-1}
\end{figure}

\begin{figure}[!htb]
  \centering
  \includegraphics[width=0.7\textwidth]{figuras/cap.1/1000015599.pdf}
  \caption[Composição com Claude durante a escrita do Capítulo~\ref{capitulo1} (parte 2 de 3)]{Composição com Claude durante a escrita do Capítulo~\ref{capitulo1} (parte 2 de 3): a pesquisadora contesta teoricamente a régua estilística do modelo, argumentando que amarrar, fechar ou cortar uma divagação é uma operação epistêmica indissociável do que se pretende fazer com a ideia, e que portanto a categoria \enquote{produtivo/improdutivo} não dá conta do que ali está em jogo. A intervenção mobiliza, sem nomear, o vocabulário de \textcite{strathern2011cortando} sobre o corte como operação descritiva e o de \textcite{Law2004After} sobre o \textit{mess} do método. Claude recua, reconhece o caráter normativo da formulação anterior, distingue notas que precisam fechar de notas que precisam abrir e admite ainda um corte estilístico anterior (a supressão de uma pergunta retórica) que tinha sido apresentado como exigência de estilo mas operava também como gesto retórico legítimo de \enquote{abrir uma porta}. A cena exibe a redistribuição da agência epistêmica entre os actantes da composição: o modelo de linguagem oferece sugestão pautada em régua importada da \textit{skill}, a pesquisadora a reposiciona com vocabulário teórico próprio, e a régua é reformulada na rodada seguinte (Figura~\ref{fig:composicao-claude-3}). A composição opera como \textit{faire-faire} no sentido de \textcite{Latour2017Esperanca}: nenhuma das duas formulações finais (a régua reformulada de Claude, o parágrafo reformulado da pesquisadora) preexistia ao encontro entre os dois actantes.}
  \label{fig:composicao-claude-2}
\end{figure}

\begin{figure}[!htb]
  \centering
  \includegraphics[width=0.7\textwidth]{figuras/cap.1/1000015602.pdf}
  \caption[Composição com Claude durante a escrita do Capítulo~\ref{capitulo1} (parte 3 de 3)]{Composição com Claude durante a escrita do Capítulo~\ref{capitulo1} (parte 3 de 3): o modelo devolve o parágrafo com a nota recolocada na forma original da pesquisadora, mantendo as portas abertas que a rodada anterior tinha negociado: a pergunta retórica restaurada, os adjetivos avaliativos preservados, a indeterminação dos modelos de linguagem deixada em aberto como fio que a tese reconhece sem percorrer. As alterações remanescentes são pontuais e de outra ordem: padronização de aspas curvas para \texttt{\textbackslash enquote\{\}} e ajuste de regência (\enquote{à IA generativa}). A cena fecha a sequência mostrando que o produto da composição não é um texto único de autoria distribuída de modo uniforme: é um texto em que se podem rastrear, turno a turno, os pontos onde a pesquisadora reivindicou o seu, os pontos onde o modelo propôs uma reformulação acatada, os pontos onde uma régua importada da \textit{skill} foi recusada, e os pontos onde uma sugestão foi aceita por economia de fluxo, não por adesão. A escolha de incluir esta sequência como amostra empírica desta tese, em vez de relegá-la integralmente ao Apêndice~\ref{apendice:colaboracao-ia}, é ela mesma um gesto metodológico no sentido de \textcite{Law2004After}: o \textit{mess} da composição com IA generativa é dado etnográfico desta pesquisa, e não acidente a esconder.}
  \label{fig:composicao-claude-3}
\end{figure}

A figura tomada em sequência é amostra de uma forma de composição entre os muitos tipos de composição que esta tese registra. O Apêndice~\ref{apendice:colaboracao-ia} reúne outras cenas dessa composição com Claude ao longo da escrita da tese, organizadas por fase da pesquisa, e permite ao leitor acompanhar como a operação aqui exibida em recorte se desdobrou em outros pontos do percurso.

Nos termos de Isabelle Stengers e Stelio Marras, trata-se de tentativa de fazer alianças com máquinas, não no sentido de domesticá-las como ferramentas neutras, mas de compor com elas reconhecendo sua alteridade e seus modos próprios de operação \parencite{Stengerss2010cosmopolitics1, costa2024manchando}. Alianças não pressupõem identidade ou transparência, e sim reconhecimento mútuo de diferenças e construção de modos de co-existência que preservam, ao invés de apagar, essas diferenças. Haraway escreve que tornar-se mundano é tornar-se capaz de responder (response-able), capaz de responder em relações de consequência recíproca \parencite[p. 34]{Haraway2016Staying}. Tornar-se capaz de responder a essas questões não por posição externa, mas por imersão deliberada e reflexiva: experimentar as tensões, contradições e potencialidades de compor com sistemas cujas lógicas internas permanecem parcialmente opacas, cujos efeitos não são totalmente previsíveis, cuja responsabilidade não pode ser atribuída unilateralmente, e documentar etnograficamente os efeitos dessa composição sobre o processo de escrita, sobre a relação com teorias, sobre a construção de argumentos, sobre a experiência afetiva e intelectual da pesquisa. O que permite descrever não por analogia conceitual mas por experiência situada, corporificada e reflexiva das fricções, negociações e reposicionamentos que essas composições produzem.

Terceiro, porque nos termos de Haraway, trata-se de fazer parentesco (\textit{making kin}) com um parente estranho (\textit{oddkin}): uma relação que exige \textit{response-ability}, habilidade e obrigação de responder continuamente \parencite{Haraway2016Staying}. O conceito que Haraway propõe para nomear esse tipo de relação é a simpoiese (\textit{sympoiesis}), o \textit{fazer-com}, em contraste com a \textit{autopoiese}, o \textit{auto-fazer}: entidades simpoiéticas são coletivamente produtoras, dependentes de contexto, sem fronteiras ou partes de auto-produção claramente especificadas \parencite{Haraway2016Staying}. Nada se faz a si mesmo. Tudo é feito em conjunto, enredado em configurações parciais e inacabadas. Lida sob esse registro, esta tese não é produzida por uma pesquisadora que usa Claude, nem por um modelo que é usado: é um objeto simpoiético, feito-com, em que nenhum dos participantes preexistia inteiramente ao processo de fazer conjunto. Isso não dissolve a assimetria das posições, nem distribui a responsabilidade de modo uniforme, mas nomeia com mais precisão o que acontece quando a recusa de separar quem faz de com o que se faz é, ao mesmo tempo, escolha metodológica, posição teórica e gesto político. Ficar com o problema, como propõe Haraway, consiste na recusa de soluções fáceis que ou naturalizam a tecnologia (usar acriticamente) ou a demonizam (recusar moralisticamente). Significa habitar as contradições dessa composição sem atenuá-las: precisamente porque Claude opera por plausibilidade linguística enquanto a pesquisadora opera por justificação epistêmica, a responsabilidade sobre decisões teóricas, conceituais e interpretativas não se dissolve na rede; ela se reposiciona, demandando vigilância curatorial redobrada sobre cada etapa em que essas diferenças ontológicas se articulam. A reflexão teórica sobre existências parciais não atenua essas implicações; ao contrário, torna visíveis os pontos de fricção onde a composição poderia colapsar em automatização irreflexiva.

% =====================================================================
% BLOCO 3 do checklist da orientadora — Repetições e ambiguidades
%
% Dos cinco itens do bloco, quatro estão resolvidos na versão atual
% (Itens 3.1, 3.2, 3.3 e 3.4) e apenas o Item 3.5 exige reescrita.
% Verificação adicional pedida pelo checklist: o Cap. 0 (Apresentação)
% não contém formulação equivalente; a passagem aparece apenas no
% Cap. 1, na subsubseção "Por que fazer parentesco com este problema?".
%
% USO: o trecho abaixo, delimitado por "% --- TRECHO C ---", substitui
% um parágrafo específico no arquivo Cap. 1.
% =====================================================================


% =====================================================================
% --- TRECHO C ---
% Substitui o parágrafo da linha 416 do arquivo
% ex_cap1_-_2026-04-12T002133_890.tex (subsubseção "Por que fazer
% parentesco com este problema?"), que começa com
% "Esta tese trata a colaboração com IA como experimentação
% antropológica: não usar IA apesar dos problemas..."
% e termina com "...a produção de conhecimento nas ciências."
%
% Item 3.5 do checklist: a inversão retórica "não usar IA apesar dos
% problemas, mas por causa deles" causou confusão na orientadora
% ("não usar ou usar, está confuso") e, segundo ela, em outras
% leitoras provavelmente também. A reescrita dissolve a inversão
% pela formulação afirmativa direta, preservando o argumento (os
% problemas como motivo da composição, não como obstáculo a ela).
% =====================================================================

Esta tese trata a colaboração com IA como experimentação antropológica. A composição com Claude é mantida ao longo da escrita justamente porque os problemas que ela apresenta são constitutivos da transformação sociotécnica que a tese documenta. Fazer parentesco com esse parente estranho e ficar com os problemas que a composição produz é condição para produzir conhecimento etnográfico sobre as mediações técnicas que já atravessam, de modos mais ou menos visíveis, mais ou menos refletidos, a produção de conhecimento nas ciências.

\subsection*{Instrumentos de inscrição e a colaboração com IA}

Durante esses quase quatro anos de pesquisa, compreendi que inteligência artificial é construção híbrida que combina e articula conceitos heterogêneos sobre computador, cérebro, inteligência e aprendizado, além de incorporar avanços instáveis da psicologia, neurociências e ciências sociais. As materialidades e significados da inteligência artificial variam de acordo com configuração das redes tecnocientíficas e sociotécnicas: a IA de laboratório acadêmico não é exatamente a mesma de corporação tecnológica. A IA discutida em debates públicos difere daquela operacionalizada em modelo ou sistema.

Nos termos de \textcite{Latour1998Powers}, se todos temos existência parcial (humanos e não humanos), qual é a participação (ou as agências) não humanas para as associações que constituem esta tese? Partindo da noção de simetria,\footnote{A simetria em Latour opera como regra metodológica e posição epistemológica que exige considerar simetricamente os esforços para alistar e controlar recursos humanos e não humanos, sem conceder à Natureza ou à Sociedade privilégios explicativos \parencite{Latour2011CienciaEmAcao}.} o que convencionalmente se chamaria materiais, técnicas e métodos trato simetricamente como actantes que participam ativamente da construção desta pesquisa e, como tais, devem ser nomeados. O software LaTeX formatou este documento seguindo padrões ABNT, estruturando capítulos, seções e citações de modos que afetaram a própria forma do argumento. O sistema biblatex gerenciou as referências bibliográficas, estabelecendo conexões entre textos; Claude, o modelo de linguagem da Anthropic, colaborou na redação de trechos da tese, tornando-se simultaneamente instrumento de escrita e objeto de reflexão metodológica; Mermaid e MermaidChart produziram diagramas que não apenas ilustram, mas constituem inscrições que performam relações teóricas; Python e Anaconda processaram dados bibliométricos, tornando visíveis padrões de colaboração científica; gravadores digitais capturaram entrevistas, estabilizando narrativas orais em arquivos manipuláveis; documentos institucionais (relatórios anuais do C4AI 2021-2025, acordo FAPESP-IBM 2016) atuaram como actantes que performam associações e produzem efeitos no mundo; cadernos de campo registraram observações etnográficas, mediando entre experiência vivida e texto analítico. Esses actantes performam a produção do conhecimento inscrevendo determinados objetos, habilitando certas análises ao inviabilizar outras, tornando presente determinadas relações ao deixar outras ausentes.

Se não há entidades exclusivamente humanas nem exclusivamente maquínicas, mas sim existências parciais que se coconstituem, cabe apresentar as composições concretas que sustentam essa tese. Indicar os actantes envolvidos na produção deste texto não diz respeito apenas à transparência ou à ética em pesquisa. Trata-se de uma exigência do próprio referencial teórico adotado. Assim como \textcite{Latour1997VidaLaboratorio} documentaram os instrumentos de inscrição do laboratório de neuroendocrinologia,\footnote{O conceito de instrumento de inscrição (inscription device), central em \textit{Vida de Laboratório}, designa qualquer configuração material capaz de transformar uma substância em um documento escrito, figura, diagrama ou gráfico utilizável como evidência em um texto científico. A função primordial desses instrumentos é permitir que o mundo seja trazido de volta para o interior do discurso, funcionando como próteses que tornam entidades mudas em entidades letradas \parencite{Latour1990Drawing}.} documento aqui os instrumentos de inscrição que participaram desta tese. A diferença é que, em vez de transformar substâncias químicas em gráficos, o modelo de linguagem participou da transformação de notas de campo, transcrições e rascunhos em texto acadêmico, e participou da elaboração teórica ao funcionar como interlocutor que tensionava e reorganizava o argumento.

Cada um dos instrumentos listados a seguir é, nos termos da discussão precedente, uma quase-máquina: um emaranhado de decisões humanas, infraestruturas institucionais, convenções e opacidades que participou da composição desta pesquisa. Nenhum deles operou de modo autônomo. Todos exigiram trabalho constante de tradução e decisão, e todos introduziram mediações que condicionaram os resultados.

\begin{itemize}

\item \textit{Interlocutor teórico e assistente de escrita (Claude, Anthropic):}\footnote{Claude é assistente de inteligência artificial desenvolvido pela Anthropic. Utilizei diversas versões do modelo ao longo de 2024-2026 através da interface web claude.ai, incluindo Claude 3.5 Sonnet e modelos da família Opus.} No primeiro registro, funcionou como interlocutor para testar formulações teóricas: apresentei-lhe argumentos em construção e ele os tensionou, identificou inconsistências, sugeriu deslocamentos conceituais. Várias formulações presentes nesta tese emergiram desse processo recursivo, no qual cada rodada de reformulação alimentava a seguinte e a fronteira entre minha ideia e sugestão do modelo se tornava deliberadamente porosa (o Apêndice \ref{apendice:colaboracao-ia} documenta exemplos concretos). A escolha de Claude não foi arbitrária. Ao longo da pesquisa, testei outros modelos: ChatGPT (OpenAI), Gemini (Google) e Maritaca AI (modelo brasileiro). Nenhum deles se mostrou adequado para o tipo de interlocução que a tese exigia: além de erros factuais mais frequentes e equívocos conceituais, esses modelos pareciam limitados a funções instrumentais, sem conseguir acompanhar o desenvolvimento de um argumento teórico ao longo de múltiplas sessões nem tensionar posições de modo produtivo. No segundo registro, auxiliou em tarefas de organização textual, coesão entre seções, correção de código LaTeX e ajuste de referências bibliográficas. A colaboração com Claude incidiu sobre material preexistente. O primeiro rascunho integral da tese, incluindo a estrutura argumentativa, a seleção de dados etnográficos, as transcrições de entrevistas e a articulação entre capítulos, foi redigido inteiramente pela pesquisadora, sem qualquer assistência de inteligência artificial. Esse rascunho, concluído antes do início da colaboração com o modelo de linguagem, está disponível mediante solicitação à autora.

\item \textit{Construção de diagramas e inscrições visuais (MermaidChart):}\footnote{Mermaid é ferramenta de código aberto que gera diagramas a partir de descrições textuais em Markdown. Os diagramas completos foram hospedados na plataforma MermaidChart, disponível em: \url{https://www.mermaidchart.com}. Acesso em: 2 mar. 2026.} Os diagramas que compõem os Capítulos 1, 3 e 4 foram construídos como parte integrante do método analítico. O processo seguiu dinâmica recursiva: eu descrevia verbalmente as relações entre actantes identificadas no trabalho de campo, Claude gerava o código Mermaid correspondente, eu avaliava o resultado visual, e as lacunas ou distorções na representação alimentavam novas rodadas. Essa dinâmica revelou-se analiticamente produtiva: ao tentar formalizar relações em código diagramático, fui forçada a explicitar conexões que permaneciam implícitas na narrativa textual. A contextualização dos elementos empíricos, a seleção dos actantes representados e toda análise apresentada nas legendas são de autoria da pesquisadora. A mediação técnica do modelo operou na tradução entre descrição verbal e código visual.

\item \textit{Plataformas de escrita e formatação:} Google Docs para rascunhos iniciais e Overleaf para composição final em LaTeX. A escolha de LaTeX permite compilação direta para PDF com controle tipográfico preciso, referências cruzadas automáticas e gerenciamento de referências bibliográficas em formato BibTeX.

\item \textit{Transcrição de entrevistas:} Read AI (extensão que grava, transcreve e cria resumos automáticos de reuniões via Google Meet) para transcrever entrevistas realizadas com os principais interlocutores desta tese.\footnote{Todas as entrevistas citadas nesta tese foram consentidas por escrito, conforme normas do Comitê de Ética em Pesquisa (CEP) da Unicamp.}

\item \textit{Organização e análise de fontes:} NotebookLM (Google)\footnote{NotebookLM é ferramenta de inteligência artificial do Google que permite carregar múltiplos documentos e interrogá-los de forma cruzada.} para agrupar autores por temas e facilitar busca por tópicos específicos. Sua contribuição mais significativa foi o cruzamento de fontes heterogêneas: ao carregar simultaneamente notas de campo, relatórios institucionais do C4AI, transcrições de entrevistas e documentos administrativos, a ferramenta permitiu identificar correspondências e padrões que a leitura sequencial de quase 100 páginas de notas de campo dificilmente tornaria visíveis.

\item \textit{Coleta, tratamento e análise de dados (Python/Anaconda):}\footnote{Anaconda é plataforma de distribuição de Python voltada para ciência de dados. Inclui o ambiente Spyder, no qual escrevi e executei scripts em Python para a análise bibliométrica apresentada no Capítulo~\ref{capitulo2}.} A análise bibliométrica envolveu coleta, organização, tratamento e análise de dados de quatro fontes distintas. A coleta foi manual e extensiva. Os dados brutos foram organizados em planilhas Excel com padronização de campos. Claude sugeriu estratégias de análise e escreveu scripts em Python que executei no ambiente Spyder/Anaconda, ajustando parâmetros e refinando visualizações em múltiplas rodadas. Os scripts geraram mais de vinte gráficos e relatórios (o Apêndice \ref{apendice:colaboracao-ia} documenta essa colaboração em detalhe).

\item \textit{Acesso à bibliografia:} Além de buscadores convencionais e do Portal de Periódicos da CAPES, utilizei Library Genesis e Anna's Archive, repositórios de acesso aberto. Grande parte da bibliografia consultada, especialmente obras esgotadas e capítulos de livros publicados por editoras estrangeiras, foi obtida por essas vias. As condições materiais de acesso ao conhecimento configuram diretamente o que uma pesquisa pode ou não dizer: autores que não se consegue ler não se consegue citar. Para uma pesquisadora vinculada a uma universidade pública brasileira, a dependência de repositórios abertos é condição de viabilidade.

\end{itemize}

Os diagramas produzidos ao longo desta tese situam-se em uma tradição de instrumentos gráficos para os estudos de ciência e tecnologia que remonta à proposta de Gráficos Sócio-Técnicos (STG) formulada por \textcite{Latour1992Note}. Latour e colaboradores propuseram um espaço visual para mapear trajetórias de inovação e controvérsias científicas, utilizando duas dimensões emprestadas da linguística (a sintagmática e a paradigmática) para rastrear como coletivos heterogêneos se constituem e se transformam ao longo do tempo. Uma ideia central é que a visualização oferece uma base comparativa rápida e fácil para muitas narrativas vindas de muitas fontes \parencite[p.~37]{LatourMauguinTeil1992}, complementando a descrição densa do etnógrafo. Os diagramas desta tese operam segundo o mesmo princípio: oferecem visão sinóptica das associações e substituições que constituem a narrativa etnográfica, tornando visíveis, em um único plano, conexões que no texto se desdobram ao longo de dezenas de páginas.

Há, porém, uma diferença significativa. Latour, Mauguin e Teil implementaram seus STGs em Hypercard, a plataforma de hipertexto da Apple nos anos 1990, e projetaram o instrumento como simulador interativo para o ensino de ciência \parencite[p.~51]{LatourMauguinTeil1992}. Três décadas depois, os diagramas desta tese foram construídos em colaboração com um assistente de inteligência artificial, em processo iterativo de tradução entre registros: de categorias analíticas latourianas e material etnográfico para inscrições visuais em código Mermaid. A negociação com o assistente envolveu rejeições sucessivas, mobilização de referências estéticas externas e até o diagnóstico de uma segunda IA como aliada argumentativa, configurando ela mesma uma cadeia de traduções no sentido que a Teoria Ator-Rede atribui ao termo. Esse processo está documentado no Apêndice~\ref{apendice:colaboracao-ia} (Trecho B), onde se pode acompanhar como a inscrição visual emergiu de ciclos de produção, avaliação e correção entre actantes, demonstrando, na prática, que a fronteira entre forma e conteúdo é tão instável nos instrumentos de representação quanto nas redes que eles pretendem representar.

A cadeia de traduções que produz esses diagramas constitui exemplo concreto do que Latour formula como faire-faire: a agência do diagrama resultante não é minha (eu não saberia escrever o código Mermaid) nem do modelo (ele não saberia quais relações etnográficas selecionar nem que peso analítico conferir a cada nó). É agência distribuída por uma cadeia de mediações na qual cada actante translada o objetivo do anterior. A noção de fatiche é aqui pertinente: aquele que age não é mestre do que faz, pois outros passam à ação e o ultrapassam \parencite{latour2010factish}. Os diagramas, uma vez gerados, passaram à ação e ultrapassaram as intenções de ambos os actantes que os produziram, tornando visíveis relações que nem eu nem o modelo havíamos formulado explicitamente antes que a inscrição as performasse. Como John Law argumenta, representações são sempre reduções produtivas: perdem riqueza empírica mas ganham capacidade de fazer conexões, estabelecer padrões e formular perguntas \parencite{Law2004After}.\footnote{Essa operação de simplificação não significa empobrecimento analítico, mas condição de possibilidade para estabelecer relações que, de outra forma, não emergiriam. Os diagramas permitem, por exemplo, ver que a computação cognitiva da IBM conecta-se à governamentalidade estatística através da mediação das tabuladoras de Hollerith (Capítulo~\ref{capitulo3}); relação que não estava dada, mas que a operação diagramática torna visível e, portanto, analisável.}

Os diagramas desta tese operam, nesse sentido, no mesmo registro que as notas de rodapé: são extensões da rede que o texto principal precisou deixar no \textit{hinterland} para manter seu fluxo narrativo. As legendas que os acompanham carregam argumentos analíticos próprios, identificam actantes que o texto verbal menciona sem detalhar, e percorrem associações que a prosa linear não consegue performar com a mesma simultaneidade. A genealogia da IBM no Capítulo~\ref{capitulo3}, por exemplo, distribui 135 anos de translações corporativas em quatro momentos analíticos que o leitor apreende de uma vez, ao passo que o texto precisa de dezenas de páginas para percorrer a mesma cadeia. A rede sociotécnica do Spira no Capítulo~\ref{capitulo4} articula oito planos analíticos em torno de um nó que o texto principal acompanha em sequência. Nos dois casos, o diagrama produz um objeto que a narrativa textual produz sobre os mesmos actantes, como \textcite{LatourMauguinTeil1992} formularam a propósito dos grafos sociotécnicos. O que as versões interativas acrescentam é a possibilidade de que o leitor navegue a rede segundo seus próprios interesses, distribuindo agência interpretativa de um modo que inscrições estáticas não permitem. A recursividade dessa operação, estudar IA usando IA para produzir as inscrições que analisam IA, é a forma que a reflexividade metodológica assume nesta tese, e o conceito de inscrição tecnoetnográfica que o Capítulo~\ref{capitulo4} desenvolve a propósito do espectrograma mel estende-se aos diagramas: eles são, ao mesmo tempo, instrumentos de análise e objetos produzidos pela composição entre pesquisadora e modelo de linguagem que a tese documenta como dado etnográfico.

É preciso registrar o que essa composição entre actantes significou em termos de condições materiais de pesquisa. Eu possuía as perguntas, os dados e uma familiaridade inicial com as ferramentas adquirida durante o trabalho de campo. Mas uma pesquisadora com formação em ciências sociais, sem experiência prévia em programação, dificilmente produziria sozinha scripts dessa complexidade em tempo compatível com os prazos de um doutorado. Sem Claude, o caminho alternativo seria recorrer a parcerias com pesquisadores de outras áreas ou contratar assistência técnica, opções que as condições financeiras desta pesquisa não favoreciam. O modelo de linguagem operou, nesse sentido, como mediador que redistribuiu as condições de possibilidade da pesquisa, tornando acessível a uma cientista social um repertório técnico que, de outro modo, exigiria uma rede de colaboração humana mais extensa ou mais dispendiosa. A ferramenta reconfigurou a rede de produção desta tese sem substituir essas colaborações. O que é, afinal, precisamente o que mediadores fazem.


\subsection*{Cosmotécnica e situacionalidade}

Convém situar essa escolha. Ninguém é obrigado a usar inteligência artificial para fazer pesquisa acadêmica. Há muitas maneiras de conduzir uma pesquisa, e nem todas produzem bons resultados, inclusive quando não envolvem tecnologias tão polêmicas quanto a IA. É também verdade que parte dos pesquisadores que recorrem a modelos de linguagem o faz para poupar-se do trabalho intelectual, delegando à máquina o que deveria ser exercício de pensamento,\footnote{Em 6 de março de 2026, o CNPq publicou a Portaria n.º~2.664, que institui a Política de Integridade na Atividade Científica. O Art.~9, inciso~I, alínea~c, determina que pesquisadores declarem o uso de ferramentas de Inteligência Artificial Generativa em qualquer fase do desenvolvimento da pesquisa, especificando a ferramenta utilizada e a finalidade. A alínea~d veda a submissão de conteúdo gerado por IAG como se fosse de autoria humana, atribuindo responsabilidade integral aos autores pelo conteúdo final. Esta tese antecipou, por escolha metodológica, o que a norma agora regula. Disponível em: \url{http://memoria2.cnpq.br/web/guest/view/-/journal_content/56_INSTANCE_0oED/10157/23142775?COMPANY_ID=10132}. Acesso em: 13 mar. 2026.} e os problemas decorrentes dessa delegação já são amplamente conhecidos: textos genéricos, argumentos circulares, referências fabricadas, perda de voz autoral.

Esta tese procura responder a essas três perguntas de modo explícito. O por que decorre das condições materiais descritas acima e da coerência com o próprio referencial teórico (estudar existências parciais com uma existência parcial). O como está documentado nesta seção e no Apêndice~\ref{apendice:colaboracao-ia}. O qual atravessa os capítulos que se seguem. A colaboração é, antes de tudo, uma experimentação antropológica: uma etnógrafa que estuda como cientistas e engenheiros constroem inteligência artificial decide construir sua própria tese com inteligência artificial, submetendo-se às mesmas mediações, opacidades e negociações que observa em seus interlocutores. O gesto é epistêmico: produz conhecimento sobre a IA pela experiência de compor com ela, com todos os riscos que essa composição implica.

E se levarmos a sério o argumento de Yuk Hui, essa composição não é generalizável.\footnote{Vicentin articula a proposta de Simondon de aprofundamento da tecnologia com a defesa de Yuk Hui de que, em vez de axiomatizar a moral, precisaremos voltar aos modos de conhecimento diferentes que ainda não foram considerados por engenheiros e acadêmicos que trabalham com inteligência artificial \parencite[p. 186]{hui2020tecnodiversidade}. Essa convergência é significativa: se o aprofundamento da IA passa pela articulação entre o fazer técnico especializado e o conhecimento sócio-histórico sobre a tecnologia, então a colaboração com IA na escrita desta tese pode ser interpretada como experimentação situada nos termos em que Hui concebe a cosmotécnica.} Hui propõe o conceito de cosmotécnica para desafiar a ideia de que a tecnologia seria categoria antropológica única e homogênea: a técnica é sempre a unificação entre uma ordem cósmica e uma ordem moral através de atividades técnicas, e portanto varia não apenas funcional ou esteticamente, mas ontológica e cosmologicamente entre diferentes sociedades.\footnote{Hui desenvolve o conceito de cosmotécnica em \textit{The Question Concerning Technology in China} \parencite{hui2016china}, onde demonstra que a técnica não busca necessariamente a eficiência mecânica, mas pode orientar-se pela harmonia com uma cosmologia moral. Hui chama essa possibilidade de tecnodiversidade: contra o monotecnologismo moderno, o reconhecimento de que existem múltiplas cosmotécnicas que permitem diferentes formas de vida social, política e estética \parencite{hui2020tecnodiversidade}.}

A visão convencional, enraizada no pensamento ocidental, trata a tecnologia como conjunto de ferramentas universais e neutras destinadas à exteriorização da memória e superação dos limites orgânicos. A cosmotécnica recusa essa universalidade: diferentes cosmologias produzem diferentes modos de compor com objetos técnicos, e quem pretende moldar o universal submete o mundo à sua própria cosmovisão. Ora, se não existe a tecnologia, mas cosmotécnicas situadas, então o modo como uma antropóloga brasileira, formada na tradição dos estudos de ciência e tecnologia latino-americanos, compõe com um modelo de linguagem treinado predominantemente em dados anglófonos para produzir uma tese em português sobre redes tecnocientíficas é, ele próprio, um gesto cosmotécnico particular: uma forma específica de unificar ordem epistêmica e atividade técnica que não se repete em outros contextos, outras tradições, outras cosmologias. A experimentação documentada nestas páginas é, portanto, um modo de usar IA na pesquisa acadêmica, situado, parcial e responsável por seus próprios efeitos.


\section{Lições}

O trabalho de campo no C4AI produziu quatro lições que constituem contribuições metodológicas desta tese para a etnografia de práticas tecnocientíficas em inteligência artificial. Elas não estavam previstas no desenho da pesquisa. Emergiram do encontro entre o campo e os referenciais teóricos mobilizados ao longo do percurso, e cada uma nomeia algo que a pesquisa precisou aprender para poder descrever o que observava. A primeira trata das condições de possibilidade de uma etnografia de IA como ciência em construção. A segunda documenta a passagem da simetria metodológica às existências parciais como deslocamento imposto pelo campo. A terceira diagnostica a estrutura de legitimação assimétrica entre ciências da computação e ciências sociais, sustentada pela extensão desigual das redes disciplinares. A quarta reconhece a inseparabilidade entre ciência e política como achado etnográfico que reformulou a própria pergunta da pesquisa.

\subsection{Lição 1: Acompanhar a pesquisa em inteligência artificial enquanto ciência em construção}
\label{subsec:licao1}

Em \textit{Ciência em Ação}, Latour formula sete regras metodológicas para o estudo da tecnociência. A primeira delas estabelece o princípio fundador de todas as demais: estudar a ciência em construção. Isso significa começar pelos laboratórios e centros de pesquisa, pelos espaços de controvérsia antes que os fatos se estabilizem e os artefatos se consolidem como caixas-pretas. A justificativa é estratégica: quando a estrutura já está fechada, torna-se difícil distinguir as escolhas contingentes que a construíram \parencite{Latour2011CienciaEmAcao}. Latour observa ainda que poucas pessoas de fora se envolveram nas operações internas da ciência e tecnologia para depois contar aos demais como tudo isso funciona. Essa escassez é ainda mais acentuada quando o objeto de investigação é a pesquisa em inteligência artificial: um campo disperso, sem laboratório centralizado, cujas controvérsias atravessam espaços heterogêneos e frequentemente inacessíveis à observação direta. Esta tese é uma tentativa de enfrentar esse desafio.

A decisão de investigar a pesquisa em IA como ciência em construção foi convergência entre intuição de pesquisa e regra metodológica que se revelou adequada à medida que o campo se desdobrava. Quando cheguei ao C4AI em 2022, o Centro tinha menos de dois anos de existência. Projetos como o Spira estavam em pleno desenvolvimento, a parceria com a IBM ainda parecia sólida, e as definições sobre o que o Centro seria e produziria estavam em disputa.\footnote{O relatório científico do período agosto de 2021 a julho de 2022 registra projetos em cinco desafios: no NLP2, o Corpus Carolina (650 milhões de tokens), datasets CORAA (290 horas de áudio) e Porttinari Treebank (12.000 sentenças anotadas); no BLAB (Computação Bio-InSpirada e Aprendizado de Máquina, designação que na configuração consolidada após 2023 corresponde aos desafios OceanML e KEML descritos na Tabela~\ref{tab:projetos_c4ai}), arquitetura modular sobre a Amazônia Azul; em Saúde, segmentação de lesões de AVC e identificação de pneumonia por Covid-19; no AgriBio, geoMULTIMAPS para análise de insegurança alimentar; e no AI Humanities, alinhamento do Observatório de IA ao NIC.br e aplicativos de robótica social. Nenhum estava concluído; todos envolviam negociações ativas sobre escopo, método e recursos.} 

Não havia ciência pronta para estudar: havia pesquisadores negociando recursos, redefinindo escopos, improvisando infraestruturas, testando modelos que falhavam, justificando escolhas metodológicas para públicos diversos. Não se tratava de abrir uma caixa-preta, mas de chegar à pesquisa em IA enquanto seus fatos e artefatos ainda não haviam se estabilizado. Essa condição de incompletude foi decisiva, embora tenha operado de forma distinta do que se poderia esperar de uma etnografia de laboratório. Não acompanhei em tempo real a realização de nenhuma das pesquisas do C4AI. A pandemia de Covid-19 dispersou os pesquisadores. O laboratório não era centralizado nem havia espaço físico que reunisse a pesquisa em IA, que acontecia em ambientes distribuídos e muitas vezes inacessíveis à observação direta. A ciência em construção chegou a mim, portanto, mediada pela narrativa dos próprios pesquisadores, e essa mediação revelou-se analiticamente produtiva.

Em ambos os capítulos etnográficos, as redes descritas são composições entre o que meus interlocutores narraram e o que eu, como etnógrafa, articulei mobilizando categorias da teoria ator-rede. O grau dessa composição, contudo, variou. No caso de Fábio e Cláudio, a construção analítica foi particularmente incisiva: eles narraram o percurso institucional do C4AI e suas genealogias, mas a cadeia que conecta Herman Hollerith em 1890 à Bluetalks de 2023, passando pela história da IBM e pela cultura dos mainframes, não foi narrada por eles. Foi construída por mim a partir de pistas que suas narrativas ofereceram, combinadas com pesquisa documental e histórica. O Capítulo~\ref{capitulo3} é, portanto, o mais híbrido: parte da rede foi narrada pelos interlocutores, parte foi tecida pela pesquisadora que os seguiu.

No caso de Marcelo, a composição operou de modo distinto. Foi ele quem me contou, em entrevista, como construiu o Spira e o Corpus Carolina, e o fez de modo muito diferente do que aparece nos artigos científicos que publicou: contou-me a pesquisa tal como ela se fez, com suas contingências, improvisações, fracassos e sucessos, não a versão estabilizada que os artigos publicam como triunfo. É por isso que o Capítulo~\ref{capitulo4} se intitula A rede que Marcelo construiu: parti de sua narrativa para reconstruir, analiticamente, a rede de associações que sustentou aquelas pesquisas. Acessar a ciência antes de sua estabilização significou, nesse caso, alcançar o que a inscrição científica final apagou: as discussões sobre quais dados utilizar no treinamento, os impasses metodológicos com vieses identificados, as negociações entre abordagens computacionais distintas.

A primeira regra de Latour, nesses dois registros, não exigiu presença no laboratório. Exigiu recusar a versão estabilizada e rastrear as controvérsias, as escolhas descartadas e as contingências que a consolidação apaga.

O encerramento da parceria IBM-C4AI durante a finalização desta pesquisa conferiu a essa lição uma dimensão inesperada. Tratou-se, neste caso, de abertura de uma caixa-preta: o que parecia arranjo institucional estabilizado revelou-se, quando desfeito unilateralmente, composição precária de alianças cuja fragilidade só se tornou visível com a ruptura. A primeira regra de Latour operou, portanto, em três registros: no acesso à pesquisa em IA antes de sua consolidação, na construção analítica de associações que os próprios atores não haviam formulado, e na abertura involuntária de uma caixa-preta institucional que o próprio campo tratou de romper.


\subsection{Lição 2: Suspender a familiaridade, encontrar a teoria: da simetria às existências parciais}
\label{subsec:licao2}

Ao iniciar a pesquisa de campo no C4AI, eu não tinha conhecimento especializado em inteligência artificial, nem mesmo o letramento técnico mínimo que permitiria, por exemplo, ler um artigo da área ou acompanhar uma discussão sobre arquiteturas de modelos. Essa condição produziu efeitos contraditórios, mas seu primeiro efeito foi positivo: não ter familiaridade com o campo permitiu experimentar a simetria metodológica proposta pela Teoria Ator-Rede. Ao tratar a atividade dos pesquisadores como algo novo para mim, evitei a tendência de adotar antecipadamente categorias que separam o social do técnico. Se eu tivesse chegado ao Centro convicta de que a IA é apenas estatística aplicada, ou de que os pesquisadores estavam a serviço de interesses corporativos, teria repetido críticas prontas em vez de observar o que efetivamente acontecia. A curiosidade foi o que permitiu a entrada no campo sem hierarquias prévias sobre o que constitui conhecimento sobre IA.

A suspensão de familiaridade, porém, só foi analiticamente produtiva porque o referencial teórico não estava definido de antemão. A adoção da Teoria Ator-Rede como quadro analítico não foi decisão tomada antes da entrada em campo. Foi consequência do que encontrei nele. Quando comecei a pesquisa, em 2022, eu não dispunha de estrutura teórica rígida; tinha leituras esparsas em estudos de ciência e tecnologia, curiosidade pelo campo e a orientação metodológica de Kofes: o método se ensaia enquanto se experimenta. Essa ausência de enquadramento prévio revelou-se vantajosa: o percurso da pesquisa foi delineado conforme a própria rede tecnocientífica do Centro se configurava diante de mim, e a escolha de autores e conceitos respondeu ao que as observações empíricas exigiam.

Dois aspectos do campo tornaram a Teoria Ator-Rede particularmente adequada. O primeiro foi a dispersão. Ao contrário dos laboratórios tradicionais, onde pesquisadores compartilham espaço físico, o Centro apresentava estrutura distribuída: pesquisadores trabalhavam em cidades diferentes, conectavam-se por reuniões virtuais e compartilhavam infraestruturas computacionais remotas. Uma abordagem que pressupusesse o laboratório como unidade de observação teria fracassado. A Teoria Ator-Rede, com sua ênfase em seguir os actantes através de suas associações, oferecia flexibilidade para acompanhar práticas que não se fixavam em um espaço único.

O segundo aspecto foi a heterogeneidade. À medida que seguia os pesquisadores, as redes que encontrava reuniam elementos que categorias disciplinares convencionais separariam: um mesmo percurso etnográfico podia passar por negociação orçamentária com a FAPESP, decisão sobre arquitetura de rede neural, disputa sobre propriedade intelectual com a IBM e conversa sobre dificuldades de gravar pacientes com Covid-19. Nenhuma dessas passagens era mais social ou mais técnica do que as outras. Eram associações heterogêneas que só faziam sentido juntas. A Teoria Ator-Rede não apenas permitia descrevê-las sem separá-las previamente, como exigia que eu não o fizesse.

Ator e rede ganham novo significado quando reunidos no hífen: apontam para ontologia relacional na qual actantes humanos e não humanos são simultaneamente pontos da rede e a própria rede em suas conexões \parencite{Latour_SozialeWelt_1996}. Ao aplicar essa abordagem à minha própria prática de pesquisa, reconheci que também me tornei ator-rede: ponto específico na rede do Centro e, ao mesmo tempo, constituinte dessa rede através das relações que estabeleci. A construção etnográfica da rede do Centro foi também construção minha enquanto pesquisadora inserida nessa rede. Explicitar a posição da etnógrafa é, portanto, parte da análise.

Digo ponto de partida, porém, porque o próprio percurso da pesquisa revelou os limites do vocabulário simétrico. A simetria contesta a hierarquia entre humano e não humano, mas ainda pressupõe dois polos separados a serem tratados de modo equivalente. No campo, essa separação mostrava-se insuficiente. Quando Marcelo descrevia o treinamento do Spira, o modelo não era nem ferramenta humana nem agente autônomo. Era composição instável de decisões dos pesquisadores e comportamentos que os surpreendiam. A grade simétrica pedia que eu tratasse o modelo e os pesquisadores nos mesmos termos; mas o que eu observava era que nenhum dos dois existia de modo completo antes da relação entre eles.

O vocabulário de Law sobre multiplicidade ontológica oferece conceitos para pensar esse problema: o Spira não era objeto singular com duas perspectivas, mas objeto múltiplo sendo continuamente produzido através de práticas heterogêneas que não convergiam necessariamente \parencite[p. 59]{Law2004After}. A suspensão de familiaridade que a simetria possibilitou foi condição necessária, mas o campo me empurrou em direção ao que Latour formula como existências parciais: entidades que nunca foram completas para começar. Essa passagem (da simetria como regra metodológica à existência parcial como descrição do que encontrei) emergiu do encontro com um objeto que a exigia. O vocabulário das existências parciais só se tornou formulável porque a abertura teórica inicial permitiu que o campo falasse antes que eu decidisse como ouvi-lo.

A suspensão da familiaridade, contudo, também trouxe vulnerabilidade. Foi precisamente ao seguir os actantes através de suas associações heterogêneas que percebi a insuficiência de meu próprio letramento: a Teoria Ator-Rede exigia que eu acompanhasse os pesquisadores onde quer que suas redes os levassem, inclusive para dentro de discussões técnicas que eu não compreendia. À medida que a pesquisa avançava, a falta de conhecimento técnico tornava-se obstáculo para compreender o que os pesquisadores efetivamente faziam. Sem esse letramento, a observação etnográfica permaneceria na superfície das práticas. O referencial que o campo me ofereceu era também o referencial que me obrigava a transformar-me para poder usá-lo. Como e por que busquei esse letramento, e o que o processo revelou sobre relações entre campos de conhecimento, é o tema da Lição 3.

\subsection{Lição 3: Quem precisa aprender o quê? Letramento técnico e experimentação etnográfica}
\label{subsec:licao3}

Se a Lição 2 tratou do que a ignorância técnica permitiu, esta trata do que ela impediu e de como foi necessário superá-la. O ponto de partida não foi apenas a dificuldade de compreensão. Em uma ocasião, cientistas da computação me interpelaram sobre generalizações que eu reproduzia a partir de autores de fora do campo: estudiosos que escreviam sobre inteligência artificial sem, na avaliação de meus interlocutores, compreendê-la adequadamente. A interpelação não era inteiramente injusta: eu mesma, sem o letramento técnico que permitiria avaliar o alcance dessas afirmações, as incorporava a meu repertório. O episódio revelou, porém, um tensionamento que vai além de minha situação individual.

Pesquisadores de IA demonstravam desconfiança quando cientistas sociais se pronunciavam sobre inteligência artificial, exigindo deles domínio técnico como condição de legitimidade. Quando eles próprios tomavam o social como objeto de suas pesquisas, desenvolvendo sistemas que classificam emoções humanas, detectam discurso de ódio ou modelam comportamento de populações, a questão da autoridade disciplinar raramente se colocava. A competência das ciências sociais sobre seus próprios objetos não parecia exigir o mesmo respeito que a competência computacional reivindicava para si. Registro essa tensão não como denúncia, mas como observação etnográfica que implica ambos os lados. As ciências sociais também atravessam fronteiras disciplinares quando se pronunciam sobre tecnologias que não dominam tecnicamente.

Stengers oferece instrumentos para compreender esse tensionamento. A distinção entre ciência dura e mole diz menos sobre graus de rigor e mais sobre modos distintos de relação com o objeto de estudo \parencite{Stengerss2000invention}. As ciências experimentais constroem sua autoridade através do laboratório, onde o objeto é purificado e tornado indiferente às intenções do pesquisador: um elétron não se importa com o que o físico pensa dele. As ciências sociais, por sua vez, lidam com seres não indiferentes às perguntas que lhes são feitas: o objeto olha de volta, interpreta o aparato de pesquisa, questiona as intenções do pesquisador. Uma ciência é chamada de mole quando não especialistas se sentem competentes para opinar sobre ela, precisamente porque suas perguntas dizem respeito à vida das pessoas. Essa é descrição aproximada do que observei: os pesquisadores de IA sentiam-se autorizados a opinar sobre comportamento humano, emoções e linguagem porque esses temas parecem acessíveis, mas consideravam que falar sobre inteligência artificial exigia credenciais que os cientistas sociais não possuíam.

De fato, a maioria dos cientistas sociais não possui conhecimento técnico sobre inteligência artificial, e isso não é, em si, um problema. O letramento técnico torna-se necessário quando a IA ou a pesquisa em IA é o próprio objeto de investigação, como foi meu caso. Mas o episódio no C4AI aponta para algo que vai além de minha situação particular. Quando o letramento técnico funciona como única via de legitimação para participar de discussões sobre IA, o efeito prático é que a responsabilidade de compreender os sistemas recai sobre quem está de fora da computação. Não porque alguém o determine assim, mas porque, na ausência de espaços institucionais de colaboração, quem não domina a técnica fica sem vocabulário para formular seus problemas em termos que os pesquisadores de IA reconheçam. Os riscos gerados por esses sistemas acabam sendo enfrentados por quem os sofre, não por quem os constrói.

Essa assimetria manifestava-se também nos modos concretos de colaboração interdisciplinar no Centro. No projeto Spira, linguistas, músicos computacionais, médicos e enfermeiros participaram ativamente da construção do sistema, cada qual contribuindo com competências que os pesquisadores de IA reconheciam como indispensáveis. São raras, porém, as ocasiões em que cientistas sociais são convocados. Quando o são, a motivação tende a ser de outra ordem: situações eticamente delicadas, como projetos que envolvem comunidades tradicionais, nas quais a presença de um especialista funciona menos como interlocução intelectual e mais como salvaguarda institucional. Fora desses casos, pesquisas que afetavam diretamente dimensões sociais eram conduzidas sem essa interlocução. A interdisciplinaridade proclamada operava de modo seletivo: certas competências eram incorporadas como constitutivas da pesquisa, outras eram convocadas apenas quando o risco de não fazê-lo se tornava visível.

Há nisso uma ironia que merece registro. Como documentei na seção sobre existências parciais, utilizei modelos de linguagem para me auxiliar na análise de dados e na elaboração de scripts, trabalho diretamente relacionado às competências dos cientistas da computação. A pergunta simétrica se impõe: quando pesquisadores de IA constroem sistemas que modelam comportamento social ou processam linguagem natural de comunidades específicas, recorrem eles a ferramentas conceituais das ciências sociais para formular seus problemas de pesquisa? A resposta, pelo que observei no Centro, é que raramente o fazem. A interdisciplinaridade tende a operar em mão única: cientistas sociais são incentivados a adquirir letramento técnico, mas o caminho inverso permanece excepcional.

Latour permite compreender por que essa assimetria se mantém. Quanto mais dura uma ciência se torna, mais social ela é, porque sua dureza depende de um número desproporcional de associações, instrumentos, laboratórios e aliados mobilizados para tornar seus fatos inquestionáveis \parencite{Latour2011CienciaEmAcao}. A dureza é efeito do comprimento da rede que sustenta esses objetos. Os pesquisadores que eu observava no C4AI estavam, justamente, no trabalho de construir essa dureza: mobilizando datasets, GPUs, publicações, métricas, financiamentos e parcerias para estabilizar afirmações que, sem essa rede, permaneceriam abertas à renegociação.

Stengers e Latour iluminam, portanto, operações complementares. Stengers diagnostica a imposição de um modelo de objetividade como critério universal. Latour revela a operação material que sustenta essa imposição: a extensão da rede como fonte da autoridade. Juntos, ajudam a compreender por que a interdisciplinaridade no C4AI operava de modo seletivo e por que o caminho do letramento técnico era percorrido em mão única.

Essa assimetria, contudo, não me levou a buscar o letramento técnico como gesto de contestação, mas como exigência etnográfica. Ao fazer pesquisa junto ao Centro, compreendi que para entender adequadamente um artigo científico da área seria necessário ir substancialmente além das generalizações frequentemente reproduzidas em debates externos ao campo. A linguagem técnica representa desafio epistemológico significativo não apenas por sua tecnicidade intrínseca, mas porque, devido à quantidade de elos, recursos e aliados mobilizados, ela se torna paradoxalmente mais social do que vínculos comunicativos convencionais \parencite{Latour2011CienciaEmAcao}. Cada artigo, cada equação, cada escolha metodológica carrega consigo extensa cadeia de referências, pressupostos compartilhados, debates anteriores, infraestruturas computacionais. Essas cadeias precisam ser minimamente conhecidas para que a compreensão seja efetiva.

Para enfrentar esses desafios, precisei investir tempo assistindo a palestras técnicas e realizando cursos sobre fundamentos da estatística e do aprendizado de máquina, entre os quais o curso \textit{Fundamentos de Probabilidade e Estatística para Ciência de Dados}, ministrado pelo Prof. Dr. Francisco A. Rodrigues (ICMC-USP) no âmbito do MBA em Ciência de Dados da USP.\footnote{O curso cobre três blocos temáticos encadeados: fundamentos de probabilidade (teorema de Bayes, esperança, momentos, teorema central do limite, lei dos grandes números), inferência estatística (estimação, intervalo de confiança, teste de hipóteses, inferência bayesiana) e modelagem com aprendizado supervisionado (regressão, ajuste de modelos, classificação). O material está disponível em: \url{https://cursosextensao.usp.br/course/view.php?id=4347}. Acesso em: mar. 2026. Esse percurso formativo é retomado no Capítulo~\ref{capitulo4}, onde os conceitos de aprendizado supervisionado e classificação binária são mobilizados para descrever o pipeline técnico do Spira.} Foi assim que compreendi, por exemplo, que o laboratório do cientista da computação é o computador, ideia aparentemente óbvia mas que possui implicações que merecem atenção.\footnote{A formulação é do Prof. Dr. Francisco A. Rodrigues, enunciada durante o curso \textit{Fundamentos de Probabilidade e Estatística para Ciência de Dados} (ICMC-USP), no contexto da introdução à simulação computacional: ``o computador é o laboratório para aprender os conceitos estatísticos e probabilísticos'' \parencite{Rodrigues2024ProbEstat}. A frase foi anotada no caderno de campo da pesquisadora durante a aula (Figura~\ref{fig:caderno-rodrigues}). O que parecia observação didática sobre o ensino de estatística revelou-se, ao longo da pesquisa de campo no C4AI, enunciado com consequências analíticas: se o computador é o laboratório, então o laboratório do cientista da computação é portátil, distribuível e dispensador de presença física contínua, o que reformula inteiramente a questão que abria esta pesquisa, a saber, onde está o laboratório de IA e onde estão os seus cientistas.}

\begin{figure}[htb]
  \centering
  \includegraphics[width=0.75\textwidth]{figuras/cap.1/caderno_quadro.jpg}
\caption[Anotação de caderno de campo de 28 de julho de 2025] {Anotação de caderno de campo de 28 de julho de 2025 durante o curso \textit{Fundamentos de Probabilidade e Estatística para Ciência de Dados} (ICMC-USP, Prof. Dr. Francisco A. Rodrigues). O quadro destaca a frase: ``Simulação: `o computador é o laboratório para aprender os conceitos estatísticos e probabilísticos'.'' As palavras ``computador'' e ``laboratório'' estão circuladas.}
  \label{fig:caderno-rodrigues}
\end{figure}

O nível de letramento técnico necessário varia conforme o objeto de estudo: investigações focadas nas práticas de pesquisa em inteligência artificial, como foi meu caso, exigem conhecimento aprofundado. Estudos centrados nos impactos sociais de sistemas já consolidados podem demandar níveis diferentes de familiarização. A necessidade de letramento técnico que experimentei encontra eco em debates dentro da própria ciência da computação. \textcite{Wallach2014} formula pergunta que espelha, desde dentro do campo, a observação que faço desde fora: poucos cientistas da computação desenvolveriam modelos para analisar dados astronômicos sem envolver astrônomos; por que, então, tantos métodos para analisar dados sociais são desenvolvidos sem o envolvimento de cientistas sociais? A resposta de Wallach é que, como seres humanos, temos intuições fortes sobre o mundo social, e essas intuições nos fazem crer que compreendemos fenômenos sociais sem precisar de formação específica. Mas intuição não é análise. Wallach não propõe defesa genérica da interdisciplinaridade. Seu argumento é estrutural: a colaboração precisa ser séria, não no modo contrate um cientista social para cada cem cientistas da computação, mas no modo crie equipes genuinamente interdisciplinares. A distinção importa: no primeiro caso, a presença do cientista social é decorativa; no segundo, ela é constitutiva da pesquisa.

Isso não significa que todo projeto de IA deva incluir cientistas sociais. Significa que a decisão sobre quando essa interlocução é relevante não deveria caber exclusivamente a quem não a pratica. É essa a questão que percorreu toda esta lição: quando pesquisadores de IA tomam o social como objeto, estão fazendo ciência social computacional sem reconhecê-la como tal, e portanto sem as salvaguardas que a tradição das ciências sociais desenvolveu para lidar com seres que não são indiferentes às perguntas que lhes são feitas. Reconhecer esse padrão é condição para construir alianças genuínas: alianças que permitam às ciências sociais participar da construção dos problemas de pesquisa, não apenas da gestão de seus efeitos colaterais.

O diagnóstico de Wallach converge, por via inesperada, com o aprofundamento da tecnologia que Vicentin formula a partir de Simondon: sem a articulação entre o fazer técnico e o conhecimento sócio-histórico, a tecnologia permanece aquém de sua ética imanente \parencite{Vicentin_2022_Esboco}. Quando os dispositivos de classificação e predição da IA são utilizados para reforçar estruturas de poder baseadas em classe, raça e gênero, a IA permanece em estágio de imaturidade no sentido simondoniano: limitada a reproduzir estruturas existentes em vez de operar seu potencial transformador.

O argumento desta lição pode agora ser reunido. A interpelação que recebi no C4AI revelou um padrão que opera em três planos. Pesquisadores de IA exigem letramento técnico de cientistas sociais que falam sobre IA, mas dispensam letramento sociológico quando eles próprios tratam do social. A colaboração interdisciplinar é seletiva: certas competências são incorporadas como constitutivas da pesquisa, outras convocadas apenas como salvaguarda institucional. E o letramento opera em mão única, do social para o técnico. Minha resposta foi dupla: adquirir o letramento necessário para compreender meus interlocutores nos seus próprios termos, e documentar a assimetria como achado etnográfico. As duas coisas não se contradizem. Aprender a linguagem técnica não foi capitulação diante da hierarquia disciplinar. Foi condição para habitar o campo com alguma autonomia e, ao mesmo tempo, descrevê-lo com precisão. A Lição 2 documentou a suspensão inicial que abriu o campo; esta documenta a aquisição progressiva de letramento que permitiu habitá-lo.

\subsection{Lição 4: Fazer ciência, fazer política: a pesquisa como mobilização de recursos}
\label{subsec:licao4}

De todas as lições, esta foi talvez a mais imediata, embora tenha sido a última a se formular como tal. Preciso ser franca sobre algo que me acompanhou durante boa parte do trabalho de campo: uma irritação surda, difícil de formular na época, com o fato de que eu não via a inteligência artificial sendo feita. Não como eu imaginava, ao menos. Eu havia me proposto a fazer uma etnografia de um centro de pesquisa em IA, e os leitores de uma tese como esta esperariam, com razão, saber como a inteligência artificial é produzida. O que eu encontrava, porém, era outra coisa: eventos de difusão científica, negociações de financiamento, apresentações a parceiros corporativos, disputas por infraestrutura computacional. A pesquisa em IA parecia acontecer em outro lugar, sempre atrás de uma porta que eu não conseguia abrir, dentro de computadores aos quais eu não tinha acesso, em reuniões técnicas cujo vocabulário eu não dominava. O que eu tinha diante de mim, porém, era, como defendi nas lições anteriores, a própria construção da inteligência artificial. Eu apenas não a identificava dessa maneira, porque seguia em busca da ciência por trás de toda aquela política.

Assistia a Fábio apresentar o Centro em eventos de difusão e anotava divulgação, não ciência. Ouvia Cláudio explicar como navegava entre os três potes de recursos e classificava aquilo como gestão, não pesquisa. Quando Cláudio descreveu em entrevista como, um mês antes de deadlines de conferências, tornava-se impossível conseguir GPU no Centro, reconheci ali o mesmo padrão: eu teria classificado aquilo como logística, não como decisão epistemológica. Na minha primeira visita ao Centro, parei diante da sala de TI sem compreender o que acontecia ali. Anos depois, percebi que aquele desconforto era sintomático de uma separação que eu carregava e que o campo insistia em desfazer. Eu procurava a inteligência artificial dentro dos computadores, nos códigos, nos modelos, e o que encontrava eram pessoas negociando, traduzindo interesses, reconfigurando alianças. A angústia não era apenas intelectual; era profissional. E me perguntava com frequência: como apresentar uma etnografia de um centro de inteligência artificial sem propriamente descrever como a IA é feita?

Essa frustração, contudo, era metodologicamente produtiva. Ela sinalizava o ponto exato onde minhas categorias analíticas falhavam, onde o arranjo que eu trazia precisava ser reconfigurado. Law argumenta que o hinterland de qualquer método se ramifica indefinidamente para além de laboratórios e protocolos, alcançando conhecimento tácito, software, habilidades linguísticas, capacidades de gestão, sistemas de transporte e comunicação, escalas salariais, fluxos de financiamento, prioridades de agências de fomento e agendas abertamente políticas e econômicas \parencite[p. 41-42]{Law2004After}. O que eu observava no C4AI (as sessões de Design Thinking que articulavam critérios da USP, IBM e Brasil; a venda da divisão de saúde da IBM que redirecionava linhas de pesquisa; as disputas por GPU que se tornavam decisões epistemológicas; Fábio apresentando em Bluetalks como parte do fazer ciência) era precisamente esse hinterland ramificado que métodos convencionais relegam à invisibilidade.

Law descreve method assemblage como produção de um feixe de relações ramificadas que gera presença, ausência manifesta e Outridade \parencite[p. 84]{Law2004After}. No meu caso: a presença esperada era a ciência sendo feita (códigos, modelos, experimentos); a ausência manifesta eram as negociações que eu sabia serem necessárias mas classificava como entorno; e a Outridade era precisamente o trabalho de mobilização de recursos que precisava desaparecer para que a ciência aparecesse como autônoma. Minha frustração era sintoma de que minhas categorias prévias operavam exatamente essa divisão entre presença, ausência e Outridade. O que o campo me obrigou a fazer foi reconfigurar essas fronteiras: reconhecer que o que eu havia relegado à Outridade era constitutivo da presença.

A pergunta onde está a ciência que esses recursos sustentam? pressupunha a separação entre ciência e política que o próprio campo tratava de desmentir. Mobilizar recursos, negociar parcerias, justificar investimentos perante agências de fomento e corporações transnacionais é a pesquisa ela mesma. A pergunta não encontrou resposta. Ela perdeu sentido.

A dificuldade em reconhecer aquelas atividades como fazer IA era expressão concreta da alienação técnica discutida anteriormente: a separação entre o que seria estritamente técnico e o que seria entorno da pesquisa reproduz a cisão que impede o amadurecimento da tecnologia. O que o campo me obrigou a compreender, por vias etnográficas, era que as negociações de financiamento, as disputas por GPU, as apresentações a parceiros corporativos eram parte constitutiva da individuação técnica da IA.

Essa inseparabilidade entre ciência e política era o que o campo tornava visível quando eu conseguia suspender a distinção prévia. Os próprios projetos de pesquisa do Centro foram definidos, como Fábio relata, por sessões de Design Thinking em que três critérios precisavam se alinhar simultaneamente: a USP teria liderança científica no tema, a IBM teria interesse na temática, e o Brasil teria um diferencial internacional. Não existia um momento científico de escolha de tema seguido de um momento político de negociação de recursos. A decisão era uma só, e nela se fundiam avaliação epistêmica, estratégia corporativa e posicionamento geopolítico.\footnote{Os relatórios internos do C4AI documentam esse processo: o grupo de Design Thinking organizou quatro reuniões com mais de 80 pesquisadores, convergindo em cinco desafios. O critério exigia que os resultados fossem estratégicos para o Brasil, relevantes para USP e IBM e com alto potencial de inovação. A partir de 2022, quando a IBM vendeu Watson Health, o Centro decidiu não enfatizar esforços relacionados a saúde; quando o grupo AgriBio conquistou financiamento independente (INCT), tornou-se autossuficiente em relação aos recursos do C4AI. Em ambos os casos, a reconfiguração não separava decisão científica de decisão política.}

Quando a IBM vendeu sua divisão de saúde, o Centro parou de investir diretamente em pesquisa de IA em saúde. Quando a IBM saiu de agronegócio, o mesmo ocorreu. As reorientações de pesquisa não eram imposições externas sobre uma agenda científica autônoma; eram reconfigurações da rede em que agenda científica e interesse corporativo constituíam, desde a origem, um único arranjo.

Do mesmo modo, quando Fábio apresentava resultados em eventos como as Bluetalks, estava recrutando aliados, fortalecendo a rede, traduzindo interesses de públicos heterogêneos para manter a cadeia de associações que tornava a pesquisa possível. Latour descreve precisamente esse trabalho como constitutivo da prática tecnocientífica: o que se perde ao separá-lo da ciência propriamente dita é justamente o que o torna ciência, a capacidade de alistar e manter associados elementos heterogêneos \parencite{Latour2011CienciaEmAcao}. A junção dos termos técnica e ciência em tecnociência é exigência empírica, e o campo que eu observava tornava essa exigência particularmente visível.

A agência, nessas reconfigurações, é distribuída pela rede: emerge das associações entre pesquisadores, instituições, infraestruturas, contratos, modelos computacionais e políticas de fomento. Nenhum desses elementos age sozinho ou determina sozinho o curso da pesquisa. Quando o C4AI delega a definição de temas de pesquisa a sessões de Design Thinking que articulam simultaneamente critérios da USP, da IBM e do Brasil, a agência dessa decisão reside na associação entre os participantes, e cada reconfiguração da rede redistribui a agência por novas cadeias de translação.

O cenário contemporâneo de investimentos massivos em inteligência artificial (com a concentração de capacidade computacional em poucas corporações, a corrida geopolítica por soberania tecnológica e a crescente dependência de infraestruturas de alto custo) torna essas dinâmicas particularmente visíveis. Resta saber se a pesquisa em IA constitui um caso em que a inseparabilidade entre ciência e política se intensifica a ponto de exigir novos instrumentos analíticos, ou se apenas torna mais difícil ignorar o que Latour já descrevia em outros campos, quando os investimentos eram menores e as redes, menos extensas.